% =====================================================================
%  Horizontal Product Branch (HPB) -- a practitioner's methodology paper
%  Author: Achraf Bouchtarouif El Khiat
%  Engine: pdfLaTeX (Overleaf / arXiv, TeX Live 2023 or later).
%  No shell-escape, no external images, no fontspec: figures are TikZ.
% ---------------------------------------------------------------------
%  NOTES FOR SECTION WRITERS AND THE INTEGRATOR
%  * Theorem-like environments defined below:
%      \begin{definition}[Name] ... \end{definition}   -> "Definition 1"
%      \begin{hpbrule}[Name]    ... \end{hpbrule}      -> "Rule R1"
%      \begin{property}[Name]   ... \end{property}     -> "Property P1"
%      \begin{remark}           ... \end{remark}       -> "Remark 1"
%  * IMPORTANT: there is NO environment called "rule". The name clashes
%    with the LaTeX primitive \rule, so the rules R1--R7 must be written
%    with \begin{hpbrule} ... \end{hpbrule}. Rules are numbered R1, R2,
%    ... and properties P1, P2, ... automatically; do not type the
%    "R"/"P" prefix by hand inside the optional argument. The rules are
%    defined ONCE, formally, in appendix-a-formal.tex (labels rule:origin,
%    rule:closure, rule:time, rule:gate, rule:lineage, rule:typed,
%    rule:discussion); the main text refers to them with \cref.
%  * Cross references: always \cref{...} / \Cref{...} (cleveref knows
%    Definition, Rule, Property, Remark, Algorithm, Figure, Table,
%    Section, Appendix). Citations: \citep{key} / \citet{key}, keys
%    from references.bib only.
%  * Label prefixes: fig:, tab:, sec:, app:, def:, rule:, prop:, alg:,
%    eq:, lst:.
%  * Code listings: the default style is "hpblatex" (LaTeX code for
%    hpb.sty). HPB Markup is typeset with
%      \begin{lstlisting}[style=hpb] ... \end{lstlisting}
%    (or \lstinline[style=hpb]|...| inline).
% =====================================================================
\documentclass[11pt,a4paper]{article}

% ---------------------------------------------------------------- layout
\usepackage[a4paper,margin=2.5cm]{geometry}

% ---------------------------------------------------------- fonts & text
\usepackage[utf8]{inputenc}
\usepackage[T1]{fontenc}
\usepackage{lmodern}
\usepackage{microtype}

% ---------------------------------------------------------------- maths
\usepackage{amsmath}
\usepackage{amssymb}
\usepackage{amsthm}

% ------------------------------------------------- figures, tables, lists
\usepackage{graphicx}
\usepackage{booktabs}
\usepackage{tabularx}
\usepackage{enumitem}
\usepackage{xcolor}
\usepackage{caption}
\usepackage{subcaption}
\captionsetup{font=small,labelfont=bf}
% float placement: allow tall figures/tables at the top of text pages and
% avoid float-only pages that are mostly empty
\renewcommand{\topfraction}{0.9}
\renewcommand{\bottomfraction}{0.5}
\renewcommand{\textfraction}{0.1}
\renewcommand{\floatpagefraction}{0.75}
\setcounter{topnumber}{3}
\setcounter{totalnumber}{4}

% ----------------------------------------------------------- algorithms
\usepackage{algorithm}
\usepackage{algpseudocode}

% ------------------------------------------------------------- listings
\usepackage{listings}
% (1) LaTeX code that uses hpb.sty (default style, Appendix D)
\lstdefinestyle{hpblatex}{
  language=[LaTeX]TeX,
  basicstyle=\ttfamily\footnotesize,
  keywordstyle=\color{blue!55!black},
  commentstyle=\color{black!55}\itshape,
  columns=fullflexible,
  keepspaces=true,
  breaklines=true,
  frame=single,
  rulecolor=\color{black!25},
  framesep=4pt,
  xleftmargin=6pt,
  xrightmargin=6pt,
  showstringspaces=false,
  upquote=true,
  moretexcs={hpbgraph,hpblane,hpbmainline,hpborigin,hpbrelease,hpbplanned,
    hpbhardware,hpbbranch,hpbstart,hpbstep,hpblabel,hpbnow,hpbcollapsed,
    hpbeos,hpblinefork,hpbcallout,hpbgate,hpbtarget,hpbstream,hpbprocess,
    hpbmilestone,hpbdepends,hpbpropagates,hpbpoststep,hpbbundle,hpbfit,
    hpbseparator,hpbset}
}
% (2) HPB Markup, the text language of Section 9 and Appendix C
\lstdefinelanguage{hpb}{
  sensitive=false,
  alsoletter={-.},
  morekeywords=[1]{hpb,config,line,use,from,on,now,eos,collapse,note},
  morekeywords=[2]{release,plan,merge,merges,into,discard,target,targets,
    force,step},
  morekeywords=[3]{improvement,imp,feature,issue,bug,defect,task,lab,
    evolution},
  morekeywords=[4]{firmware,software,hardware,other,done,todo,pass,fail},
  morecomment=[l]{\#},
  morecomment=[l]{\%\%},
  morestring=[b]",
}
\lstdefinestyle{hpb}{
  language=hpb,
  basicstyle=\ttfamily\footnotesize,
  keywordstyle=[1]\color{blue!55!black}\bfseries,
  keywordstyle=[2]\color{teal!60!black}\bfseries,
  keywordstyle=[3]\color{violet!70!black}\bfseries,
  keywordstyle=[4]\color{orange!65!black},
  stringstyle=\color{black!70},
  commentstyle=\color{black!50}\itshape,
  columns=fullflexible,
  keepspaces=true,
  breaklines=true,
  frame=single,
  rulecolor=\color{black!25},
  framesep=4pt,
  xleftmargin=6pt,
  xrightmargin=6pt,
  showstringspaces=false,
  upquote=true,
  numbers=none
}
\lstset{style=hpblatex}

% ------------------------------------------------------- plots & HPB glyphs
\usepackage{pgfplots}
\pgfplotsset{compat=1.18}
\usepackage{hpb}
% one canvas for every HPB diagram (layout rule L1): column, lane, key width
\hpbset{col=0.75cm, lane=1cm, legend=1.7cm}

% ----------------------------------------------------------- bibliography
\usepackage[numbers,sort&compress]{natbib}

% ---------------------------------------------------------- hyperlinks
\definecolor{hpbLinkBlue}{HTML}{1E3A8A}
\usepackage[colorlinks=true,linkcolor=hpbLinkBlue,citecolor=hpbLinkBlue,
  urlcolor=hpbLinkBlue,filecolor=hpbLinkBlue,
  pdfauthor={Achraf Bouchtarouif El Khiat},
  pdftitle={Horizontal Product Branch (HPB): A Practitioner's Methodology
    for Planning and Tracing Product Evolution across Hardware, Firmware
    and Software},
  pdfkeywords={product management, product evolution, hardware, firmware,
    software, release planning, traceability, branch-and-merge,
    configuration management, visual notation, diagrams as code}]{hyperref}
\usepackage[capitalise,noabbrev]{cleveref}

% ----------------------------------------------- theorem-like environments
\theoremstyle{definition}
\newtheorem{definition}{Definition}
\newtheorem{hpbrule}{Rule}
\renewcommand{\thehpbrule}{R\arabic{hpbrule}}
\newtheorem{property}{Property}
\renewcommand{\theproperty}{P\arabic{property}}
\theoremstyle{remark}
\newtheorem{remark}{Remark}

\crefname{definition}{Definition}{Definitions}
\Crefname{definition}{Definition}{Definitions}
\crefname{hpbrule}{Rule}{Rules}
\Crefname{hpbrule}{Rule}{Rules}
\crefname{property}{Property}{Properties}
\Crefname{property}{Property}{Properties}
\crefname{remark}{Remark}{Remarks}
\Crefname{remark}{Remark}{Remarks}
\crefname{algorithm}{Algorithm}{Algorithms}
\Crefname{algorithm}{Algorithm}{Algorithms}
\crefname{appendix}{Appendix}{Appendices}
\Crefname{appendix}{Appendix}{Appendices}
\crefname{lstlisting}{Listing}{Listings}
\Crefname{lstlisting}{Listing}{Listings}

% ------------------------------------------------------------- shorthands
\newcommand{\HPB}{HPB}
% Public repository of HPB Markup and the open-source renderer.
% Same URL as "repository" in every package.json of the toolkit. If it is
% emptied, the paper prints a neutral wording instead of a link.
\def\hpbrepourl{https://github.com/AchrafBr1/horizontal-product-branch}
\newcommand{\hpbrepo}{\ifx\hpbrepourl\empty
  the HPB Markup repository (its address is given with the first public
  release)\else\expandafter\url\expandafter{\hpbrepourl}\fi}
\providecommand{\keywords}[1]{%
  \par\medskip\noindent{\small\textbf{Keywords:} #1}\par}

% ================================================================ front
\title{Horizontal Product Branch (HPB): A Practitioner's Methodology for
  Planning and Tracing Product Evolution across Hardware, Firmware and
  Software}
\author{Achraf Bouchtarouif El Khiat\\\texttt{achrafbouchtarouif@gmail.com}}
\date{November 2026\\[4pt]\small Preprint, version 1.0\quad
  DOI: \href{https://doi.org/10.5281/zenodo.24355435}{10.5281/zenodo.24355435}}

\begin{document}

\maketitle

\begin{abstract}
% ---------------------------------------------------------------------
%  00-abstract.tex -- abstract text only (wrapped by main.tex)
% ---------------------------------------------------------------------
For more than seven years I have managed complete products as a senior
product manager: hardware, firmware and software, across a large portfolio
with constant rotation of references, improvement requests and competitive
pressure. Backlogs, Kanban, Scrum, roadmaps, Stage-Gate and PLM change
orders each solve part of that job, but none answers the questions I am
asked every week: from which version did this work start, what else is
running on the same product, and which release actually closed it.
Software teams answer exactly these questions for code with
branch-and-merge. The \emph{Horizontal Product Branch} (HPB) methodology
applies branch-and-merge to the product, not to the code: every piece of
work is a typed branch that forks from a released version and ends in a
named release or in a discard with a written reason. This paper defines
HPB through seven rules, a standard notation, an operating guide and
metrics that fall out of the graph. It also presents HPB Markup, an
open-source text language that draws HPB diagrams from a few lines of text
in HTML, React, Next.js, Vite or the command line, and a prototype built
into a product management portal.

\end{abstract}

\keywords{product management; product evolution; hardware, firmware and
  software; release planning; traceability; branch-and-merge;
  configuration management; visual notation; diagrams as code}

% ================================================================= body
% ---------------------------------------------------------------------
%  01-introduction.tex
% ---------------------------------------------------------------------
\section{Introduction}
\label{sec:introduction}

I have spent more than seven years as a senior product manager responsible
for complete products: the hardware, the firmware that runs on it and the
software around it. The portfolio is large and it moves all the time. New
references come in, others are replaced or reach end of life, customers
and sales ask for improvements, competitors force changes, and the product
is expected to get better continuously, not once a year. That kind of job
is less about any single project and more about keeping many product
states, and the work around them, under control at the same time.

Over those years I used every method that product organizations rely on
today: backlogs, Kanban boards, Scrum, roadmaps and Gantt charts,
Stage-Gate reviews, engineering change orders in PLM, and, on the software
side, git and git-flow. Each of them is good at something. None of them
answered the questions I actually had to answer, every week, about my
products.

\subsection{The questions I could not answer}
\label{sec:intro-questions}

Take an illustrative product, the DMS-4000, a fictitious smart door
controller with its own hardware, firmware and companion software. On an
ordinary day its product manager is asked things like these:
\begin{itemize}[leftmargin=1.5em,itemsep=2pt,topsep=3pt]
  \item \emph{A customer reports a defect.} Was it seen on firmware v1.3 or
        on v1.4? Is it already fixed in a later release, or is someone
        still working on it?
  \item \emph{The next release is being planned.} Two improvements and
        three defects are open. Which of them are committed to v1.5? Can
        v1.5 ship if one of them is not finished?
  \item \emph{A component shortage forces a hardware revision.} Which open
        work moves to the new reference? Which stays with the units already
        installed? How long do we keep maintaining the old one?
  \item \emph{An auditor, a key customer or the support team asks} what
        went into v1.4, and why a reported problem was closed without a
        fix.
\end{itemize}

None of these questions is unusual. Yet with the usual tools the answer
means opening a tracker, a spreadsheet, a PLM change order and somebody's
memory, and stitching them together by hand. What is missing is always the
same thing: the link between a piece of work and the \emph{product states}
it belongs to. I kept running into three blind spots:
\begin{enumerate}[label=(\roman*),leftmargin=2.2em,itemsep=1pt,topsep=3pt]
  \item \textbf{Origin.} From which version or hardware reference did this
        work start?
  \item \textbf{Parallelism.} What else is running on the same product, and
        on which line, at the same time?
  \item \textbf{Closure.} Which concrete version incorporated the result,
        or, if none did, why was it abandoned?
\end{enumerate}
Version fields, release milestones and change-order references capture
fragments of this. They are attributes of individual records, not a
structure you can look at, check and plan with. The result is familiar to
anyone in the role: release notes written by hand, defects that cannot be
traced to the version in which they escaped, versions shipped while
committed work was quietly still pending, and hardware revisions whose
relation to the previous reference lives only in people's heads.

\paragraph{When the gap became obvious.}
% AUTHOR: the vignette below is a composite written from the situations
% described above, with no company, product or number. Replace it with a
% real, anonymised story from your own practice if you have one you can
% tell (4-6 sentences, first person, no confidential details).
The gap did not show up in a single project; it showed up as a pattern.
A composite of it, with every detail removed, goes like this. A unit comes
back from the field with a fault, on a hardware reference that had been
revised a few months earlier. Support asks a simple question: is the
firmware fix for this fault already in the build that ships with the new
revision? The fix exists, the ticket is closed and the revision is
released, yet nobody in the room can say whether the one went into the
other without opening three tools and asking two people. Every time I saw
that scene repeat, the missing piece was the same: a record of which
product state each piece of work started from and ended in.

\subsection{What HPB is}
\label{sec:intro-what}

Software engineering solved a structurally identical problem long ago, for
code. In branch-and-merge version control~\citep{chacon2014progit} every
line of work forks from a known commit, runs in parallel with others and
is merged back when it is done; branching models such as
git-flow~\citep{driessen2010gitflow} turned that graph into a picture every
developer recognizes. My proposal is simple: \emph{apply the same model to
the product, not to the code}.

The \emph{Horizontal Product Branch} (HPB) methodology draws each product
as a horizontal main line of baselines (its released and planned
versions). Every improvement, defect, task or laboratory evaluation is a
typed branch that forks from a released baseline and ends in exactly one
of two ways: merged into a named version, or discarded with a written
reason. A hardware revision that creates a new reference forks the line
itself, so the new reference keeps its ancestry. A planned version is a
gate: it does not ship while committed work is still unverified. The name
comes from the drawing: time runs left to right, and the product is the
horizontal line from which all work branches off and to which it returns.

An HPB branch is not a code branch. It is a unit of \emph{change intent}
on the product, which may be implemented by many commits in several
repositories, by a laboratory campaign or by a hardware redesign. HPB
records where that intent started, what happened to it and where it ended.
Throughout the paper, a \emph{product line} means one product reference
followed over time, not a family of variants in the sense of software
product line engineering.

HPB did not come from a single insight. It is the result of accumulating
those blind spots over years of running a real portfolio, trying to patch
them with each tool in turn, and noticing that the patches always pointed
to the same missing structure. I have since built a working prototype of
it inside the product management portal I use, and the improvement points
it captures can be applied well beyond my own products: to any
organization that ships identifiable versions of a product over time,
and especially to products where hardware, firmware and software evolve
together and several references coexist in the field.

\subsection{What this paper gives you}
\label{sec:intro-gives}

This is a methodology paper. It describes a method and everything you need
to apply it consistently; it does not report a controlled study. It
contains:
\begin{enumerate}[label=\textbf{\arabic*.},leftmargin=2em,itemsep=2pt,topsep=3pt]
  \item \textbf{The method.} A small model (product lines, baselines,
        typed branches, merges, discards, line forks and release gates)
        and seven rules, R1--R7, that keep the graph honest
        (\cref{sec:onepage,sec:model}). The full formal model, with
        properties and proof sketches, is in \cref{app:formal}.
  \item \textbf{A standard notation.} Lanes, glyphs and colors with
        redundant encoding, plus layout rules, so that an HPB diagram looks
        the same whoever draws it (\cref{sec:notation}), shown on the
        situations a product manager meets in practice
        (\cref{sec:practice}).
  \item \textbf{An operating guide and metrics.} Practices, cadence, roles,
        coexistence with Scrum, Kanban, Stage-Gate and PLM, an adoption
        path you can start on Monday (\cref{sec:operating}), and flow and
        quality metrics computed directly from the graph
        (\cref{sec:metrics}).
  \item \textbf{Tooling.} HPB Markup, a text language in the spirit of
        Mermaid that turns a few lines of text into an HPB diagram, with an
        open-source renderer for plain HTML, React, Next.js, Vite and the
        command line, and a prototype built into a product management
        portal (\cref{sec:tooling}).
\end{enumerate}
All examples, including the DMS-4000 and ``Product A'', are fictitious and
only illustrate the method. I make no claims about measured effectiveness;
\cref{sec:limitations} is explicit about what has not been shown and how
an organization can evaluate HPB for itself.

\subsection{How to read this paper}
\label{sec:intro-read}

If you have ten minutes, read \cref{sec:gaps} for the problem and
\cref{sec:onepage} for the whole method on one page. If you want to apply
it, add \cref{sec:notation,sec:practice,sec:operating}. If you want to
build tools for it, go to \cref{sec:tooling} and
\cref{app:formal,app:markup,app:package}. \Cref{app:card} is a
quick-reference card of the notation you can print and pin next to your
desk.

% ---------------------------------------------------------------------
%  02-gaps.tex -- where today's methods fall short
% ---------------------------------------------------------------------
\section{Where Today's Methods Fall Short}
\label{sec:gaps}

HPB does not replace the methods below; I still use most of them. It sits
on top of them and fills a specific gap that each one leaves open. This
section goes through them one family at a time, the way they show up in a
product manager's week, and says plainly what each does well and where it
stops. A comparison table and a three-point gap statement close the
section.

\subsection{Backlogs and Kanban boards}
\label{sec:gaps-boards}

A backlog is an ordered list of what should be done next. A Kanban
board~\citep{anderson2010kanban} shows where each item is in the workflow,
limits work in progress (WIP) and makes flow visible. Both are excellent at
answering ``what should the team do now?'' and ``where is it stuck?''.
Reinertsen's work on product development flow~\citep{reinertsen2009flow}
explains why that matters: queues and WIP drive lead time far more than
most people expect, and Little's law~\citep{little1961} ties them
together.

What a board does not record is the \emph{product state}. A card in
``Done'' tells me the work is finished, not which version contains it. A
card in ``In progress'' does not say whether the defect was seen on v1.3 or
v1.4. Two cards for the same product sit next to two cards for another
product, in the same column, with nothing that shows they belong to
different lines. And once the card is archived, the link between the work
and the version it went into is gone. In practice the board answers
\emph{where is the work}, never \emph{where is the product}.

\subsection{Scrum}
\label{sec:gaps-scrum}

Scrum~\citep{schwaber2020scrum} gives a team a rhythm: an ordered product
backlog, time-boxed sprints and a usable increment at the end of each one.
For a software team shipping continuously, the increment and the release
can be almost the same thing.

For a product that combines hardware and firmware they are not. A sprint
increment is not a firmware release that goes through certification, nor a
hardware revision that goes through a pilot lot. Several sprints, sometimes
of several teams, feed one product version; some work lives in the lab and
never enters a sprint at all. Scrum organizes the team's time well, but it
does not say which product baseline an item started from or ended in, and
it was never meant to.

\subsection{Roadmaps and Gantt charts}
\label{sec:gaps-roadmaps}

Roadmaps and Gantt charts answer ``when''. They are the tool management
asks for, and they are good at showing dates, dependencies between
activities and the big picture of a portfolio over quarters.

Their weakness is that they plan activities, not product states. A bar
that says ``low-power mode, March to May'' does not say which version will
carry it, and when the date slips nobody can see which release is now
exposed. They are also forward-looking: once the period is over, a roadmap
is rarely updated to show what actually happened, so it cannot be used for
audit. In my experience the roadmap and the reality drift apart quickly,
and the reconciliation is done in a meeting, not in the artifact.

\subsection{Stage-Gate}
\label{sec:gaps-stagegate}

Stage-Gate~\citep{cooper1990stagegate} brought something very valuable
into product development: explicit decision points. At each gate a
cross-functional review decides, against known criteria, whether to
continue. The phases-and-decisions view of product
development~\citep{ulrich2016pdd} is still how most hardware programs are
run.

The gates, however, apply to \emph{projects}. Once the product is on the
market and evolving through firmware releases, minor hardware revisions
and a steady stream of improvements and fixes, there is no project left to
gate, only a sequence of versions. What I needed was the same discipline
applied to each release: this version does not ship until the work
committed to it is done and verified, or explicitly deferred.

\subsection{PLM and engineering change orders}
\label{sec:gaps-plm}

On the hardware side, PLM systems~\citep{stark2015plm} and engineering
change management~\citep{jarratt2011ec} give strong traceability: change
requests and orders, item revisions, effectivity, approvals. They build on
the configuration management discipline of
baselines, change control and status
accounting~\citep{ieee828,iso10007}, which is very close in spirit to HPB.

The gap is scope and form. Change orders are document-centric and live in
a different system from the tracker where firmware and software defects and
improvements are managed. The PM ends up with hardware lineage in one tool
and the rest of the product's evolution in another, and nothing that shows
both on the same timeline. Nor does a change order give a quick picture of
everything open on a product at once.

\subsection{Git and git-flow}
\label{sec:gaps-git}

Version control is where the right model already exists. In
Git~\citep{chacon2014progit} every line of work has an explicit origin,
parallel lines are first-class, and every merge is a recorded event, so the
content of each release can be derived from the graph. Git-flow~\citep{driessen2010gitflow}
added roles (feature, release, hotfix) and a horizontal diagram that
developers read at a glance; trunk-based development and continuous
delivery~\citep{humble2010cd} pushed integration further. Versioning
conventions such as Semantic Versioning~\citep{semver} even encode the
nature of a change in the version name.

But it all stops at the repository. A git graph does not contain the
customer's defect report, the lab evaluation of a new sensor, the
hardware revision forced by a shortage, or the decision to drop an
improvement because the bill of materials got too expensive. And a product
with firmware, an app and a cloud service has several repositories, none of
which is ``the product''. The model is right; it is applied at the wrong
level.

\subsection{Issue trackers with version fields}
\label{sec:gaps-trackers}

A word on the tool most teams use every day. Issue trackers usually offer
\emph{affects version} and \emph{fix version} fields, and release pages
that list unresolved items. Used with discipline, they record origin and
closure for each item. They still give no structural view across items: no
picture of what is running in parallel on a product, no lineage between
hardware references, and nothing that stops a version from being marked
released while committed items are open. HPB can be built directly on top
of such a tracker (\cref{sec:implementation}).

\subsection{Comparison and the gap}
\label{sec:gaps-comparison}

\Cref{tab:comparison} rates these approaches against eight capabilities
that follow from the questions in \cref{sec:introduction}, listed in its
caption: from recording where work starts and which version closes it, to
hardware lineage, release gates, maintained versions and configurations.
Each approach is rated as it is commonly practiced, not as it could be
extended; most tools can be configured to cover more.

\begin{table}[t]
  \centering
  \small
  \caption{How common approaches cover the eight capabilities a product
    manager needs, as they are usually practiced. (a)~fork baseline
    recorded; (b)~closing version recorded; (c)~parallel work per product
    visible; (d)~hardware lineage preserved; (e)~release gate on committed
    work; (f)~one visual artifact for planning and audit; (g)~several
    maintained versions in parallel; (h)~configurations and variants.}
  \label{tab:comparison}
  \setlength{\tabcolsep}{3pt}
  \begin{tabularx}{\linewidth}{@{}>{\raggedright\arraybackslash}X cccccccc@{}}
    \toprule
    Approach & (a) & (b) & (c) & (d) & (e) & (f) & (g) & (h) \\
    \midrule
    Backlog or list (e.g., Scrum product backlog~\citep{schwaber2020scrum})
      & -- & -- & -- & -- & -- & -- & -- & -- \\
    Kanban board~\citep{anderson2010kanban}
      & -- & -- & $\sim$ & -- & -- & $\sim$ & -- & -- \\
    Roadmap or Gantt chart
      & -- & $\sim$ & $\sim$ & $\sim$ & $\sim$ & $\sim$ & $\sim$ & -- \\
    Tracker with affects/fix version fields
      & \checkmark & \checkmark & $\sim$ & -- & $\sim$ & -- & $\sim$ & $\sim$ \\
    ALM or requirements traceability tool
      & $\sim$ & \checkmark & $\sim$ & -- & $\sim$ & -- & $\sim$ & $\sim$ \\
    Stage-Gate process~\citep{cooper1990stagegate}
      & -- & -- & -- & -- & $\sim$ & -- & -- & -- \\
    PLM engineering change management~\citep{stark2015plm,jarratt2011ec}
      & \checkmark & \checkmark & $\sim$ & \checkmark & \checkmark & -- & $\sim$ & \checkmark \\
    Git with a branching model, for code~\citep{driessen2010gitflow}
      & \checkmark & \checkmark & \checkmark & -- & $\sim$ & $\sim$ & \checkmark & -- \\
    \midrule
    \textbf{HPB (by design)}$^{\dagger}$
      & \checkmark & \checkmark & \checkmark & \checkmark & \checkmark
      & \checkmark & \checkmark & $\sim$ \\
    \bottomrule
  \end{tabularx}

  \smallskip
  \begin{minipage}{\linewidth}
    \footnotesize\emph{Note.} \checkmark: supported as commonly practiced;
    $\sim$: partial, or only through configuration or convention; --: not
    supported. Ratings are my qualitative assessment of typical use, not
    the result of a systematic tool review. Roadmaps get $\sim$ for (d)
    because they often show a successor reference as a new bar, without a
    recorded link. Stage-Gate gets $\sim$ for (e) because its gates apply
    to projects, not to each release of a product in the market. PLM change
    orders gate the release of affected items and record effectivity, but
    they rarely cover firmware and software work and give no overview,
    hence -- for (f). Git models get $\sim$ for (f) because their graph
    covers code only, and $\sim$ for (e) because protected branches and
    checks gate merges, not product releases. $^{\dagger}$HPB is rated by
    construction, not by evidence from use (\cref{sec:limitations});
    configurations are an optional extension and variants are drawn as
    sibling lines, hence $\sim$ for (h).
  \end{minipage}
\end{table}

Read across, the table says two things. First, almost every capability
already exists somewhere: Git has origin, parallelism and closure for
code; trackers record versions item by item; PLM has lineage, gates and
effectivity for hardware; Stage-Gate has decision points; Kanban makes flow
visible. Second, they are almost never available in one place, at the
level of the \emph{product}, across hardware, firmware and software, which
is exactly the level at which a product manager plans releases and is held
accountable for them. HPB does not claim to invent these capabilities. Its
contribution is to put them together in one typed graph, with checkable
rules, a standard picture and an operating routine.

\paragraph{The gap, in three points.}
\begin{itemize}[leftmargin=1.5em,itemsep=2pt,topsep=3pt]
  \item \textbf{No product-level graph.} Origin, parallelism and closure
        are recorded for code, not for the product; work on the product is
        tracked as items, not as branches of a line of versions.
  \item \textbf{Hardware and software live apart.} Hardware lineage sits in
        PLM, firmware and software work in trackers and repositories, and
        nothing shows both on the same timeline.
  \item \textbf{Releases are not gates.} Commitments to a version are
        implicit, so a version can ship while committed work is still open,
        and nobody can later prove what was in it and why.
\end{itemize}

Research on software evolution, traceability and release planning covers
parts of this; its oldest lesson, that a system in use must keep changing
or become less useful~\citep{lehman1980}, is the daily reality of a
product manager. HPB records the link between work and product states
\emph{when each decision is made}, for the whole product, hardware
included.

% ---------------------------------------------------------------------
%  03-hpb-in-one-page.tex -- the whole method on one page
% ---------------------------------------------------------------------
\section{HPB in One Page}
\label{sec:onepage}

\paragraph{The idea.}
Draw the product, not the work. Each product reference is a thick
horizontal line, the \emph{main line}, running left to right in time. Its
versions sit on it as \emph{baselines}: released ones on the left of the
NOW marker, planned ones on the right. Everything you do to the product is
a \emph{branch}: it leaves the main line at the released baseline it
started from and comes back to the version that closes it, or ends in a
cross with a written reason. Improvements and lab work run above the line,
defects and tasks below. When a hardware change creates a new reference,
the line itself forks and the new reference starts its own main line,
linked to its parent. A planned version is a \emph{gate}: it is released
only when the work committed to it has passed verification.

That is the whole method. \Cref{fig:anatomy} shows every element on a
single product line, ``Product A''. Read it from left to right: v1.0 is
released; defect X-31 forks from it and is fixed in v1.0.1; improvement
I-12 forks from v1.0.1 and is merged into v1.1; task T-7 is discarded
because it was superseded; lab evaluation L-3 is still open, with one test
passed and one failed; and defect X-40 is committed to the planned v1.2,
whose gate is still closed.

% F1 -- Anatomy of an HPB diagram: the glyphs of a single-line diagram.
% Illustrative data. Requires hpb.sty v1.1. Canvas: document defaults.
% Callout leaders never cross a label or a branch other than their target.
\begin{figure}[t]
\centering
\hpbfit{%
\begin{hpbgraph}
  % ---- lanes: proactive work above, reactive work below the main line
  \hpblane{2}{lab}{Lab (L)}{0}{19.6}
  \hpblane{1}{imp}{Improvement (I)}{0}{19.6}
  \hpblane{-1}{iss}{Issue (X)}{0}{19.6}
  \hpblane{-2}{task}{Task (T)}{0}{19.6}
  % ---- main line: history (solid) up to NOW, plan (dashed) afterwards
  \hpbmainline{0}{16}
  \hpbmainline*{16}{19.6}
  \hpborigin{0}{Product A}
  \hpbrelease{1}{v1.0}   \hpbcollapsed{1}{12}
  \hpbrelease{4}{v1.0.1}
  \hpbrelease{11}{v1.1}
  \hpbplanned{18}{v1.2}{Mar 2027}
  % ---- branches (upward forks from v1.0.1 clear its wide tag)
  \hpbbranch[label={X-31 Crash}]{iss}{1}{-1}{2}{3}{merge}{4}
  \hpbbranch[forkoff=0.5, label={T-7 Migrate logs}]{task}{4}{-2}{5}{14}{discard}{}
  \hpbbranch[forkoff=0.95, mergeoff=0.6, label={I-12 Faster pairing}]{imp}{4}{1}{6}{10}{merge}{11}
  \hpbbranch[forkoff=0.7, label={L-3 New sensor}]{lab}{4}{2}{7}{16}{open}{}
  \hpbbranch[label={X-40 Battery drain}]{iss}{11}{-1}{12}{16}{target}{18}
  % ---- steps on branches
  \hpbstep{lab}{8}{2}{pass}
  \hpbstep{task}{9}{-2}{done}
  \hpbstep{lab}{13}{2}{fail}
  \hpbstep{iss}{15}{-1}{todo}
  % ---- rationale of the discard, gate of the planned baseline, NOW
  \hpblabel[below center]{14}{-2}{\itshape rationale: superseded}
  \hpbgate[1.6]{18}{1}{-1}{gate v1.2}
  \hpbnow{16}{2}{-2}
  % ---- callouts, upper row (targets reachable without crossing a path)
  \hpbcallout{-1.15,3.35}{lane legend}{-1.15,2.3}
  \hpbcallout{1.25,3.35}{main line\\(the product)}{0.45,0.05}
  \hpbcallout{3.6,3.35}{version tag\\(released)}{1.15,0.95}
  \hpbcallout{6.2,3.35}{fork from\\baseline}{4.7,1.3}
  \hpbcallout{11.0,3.35}{lab test steps\\pass / fail}{13,2.15}
  \hpbcallout{14.3,3.35}{open branch\\(head at NOW)}{16,2.2}
  \hpbcallout{18.0,3.35}{planned baseline\\and target date}{18,1.0}
  % ---- callouts, lower rows
  \hpbcallout{-1.15,-3.25}{lane guide}{-0.05,-2}
  \hpbcallout{2.6,-3.25}{start node\\with type letter}{2,-1.2}
  \hpbcallout{7.6,-3.25}{sub-task step\\(done)}{9,-2.15}
  \hpbcallout{11.4,-3.25}{discard with\\rationale}{13.1,-2.45}
  \hpbcallout{16.3,-3.25}{pending step\\(todo)}{15,-1.15}
  \hpbcallout{0.4,-4.15}{collapsed history:\\12 merged branches}{1,-0.6}
  \hpbcallout{4.6,-4.15}{merge into\\baseline}{3.5,-0.6}
  \hpbcallout{18.45,-4.15}{targeted branch\\(dashed to gate)}{17.5,-0.6}
\end{hpbgraph}}
\caption{Anatomy of an HPB diagram, on one illustrative product line. The
thick line is the product and its versions; improvements and lab work run
above it, issues and tasks below. Every branch starts at the version it
forked from and ends in a merge (check mark), a discard (gray cross, with its
reason) or an open head on NOW, continued as a dashed line into the gate of
its planned target. Glyphs for several lines or streams are shown in
\cref{fig:hardware-fork,fig:backport,fig:shared} and in \cref{app:card}.}
\label{fig:anatomy}
\end{figure}


\paragraph{Four kinds of branch.}
Every branch has exactly one type, shown by its lane, its color and the
letter in its start node: \textbf{I}~improvement (a change that makes the
product better for users), \textbf{L}~laboratory evolution (evaluation of
a new component, design or technology), \textbf{X}~issue (a defect) and
\textbf{T}~task (engineering or operational work that keeps the product
correct and operable). The split is deliberate: I and L move the product
forward and sit above the line; X and T keep it healthy and sit below.

\paragraph{The seven rules.}
HPB is defined by seven rules. They are short because each one closes a
specific hole I kept falling into. Their formal statements are in
\cref{app:formal}; here is what they say and why they exist.
\begin{description}[leftmargin=1.2em,style=nextline,itemsep=3pt,topsep=4pt]
  \item[\Cref{rule:origin}: single origin.]
    Every branch forks from exactly one released baseline, fixed when the
    work is accepted. For a defect, that is the version where it was seen.
    \emph{Why:} without it, ``relative to which version?'' has no answer
    and you cannot tell in which version a defect escaped.
  \item[\Cref{rule:closure}: explicit closure.]
    Every closed branch ends in exactly one way: merged into one baseline,
    or discarded with a written reason. Nothing is closed silently.
    \emph{Why:} ``closed'' with no version and no reason is the most common
    way traceability is lost. Six months later nobody can tell whether it
    was fixed, dropped or forgotten.
  \item[\Cref{rule:time}: time moves forward.]
    A branch merges only into a baseline that comes after its origin, on
    the same line or a descendant of it. \emph{Why:} a version cannot
    contain work that was requested after it was frozen, so change logs
    contain only what really went in.
  \item[\Cref{rule:gate}: gate integrity.]
    A planned version is released only when every branch committed to it
    has passed verification against a candidate of that version, and
    everything it depends on is released too. Work that is not ready is
    re-targeted or discarded, explicitly. \emph{Why:} it stops the most
    expensive habit in product work, shipping a version while committed
    work is quietly still pending.
  \item[\Cref{rule:lineage}: lineage and life cycle.]
    A hardware change that is not interchangeable in form, fit or function
    creates a new line with a permanent link to its parent; references are
    never edited in place, and a retired reference does not accumulate new
    work. \emph{Why:} units already in the field keep their history, and
    you always know what to maintain and for how long.
  \item[\Cref{rule:typed}: typed work.]
    Every branch has one type; sub-tasks and test runs are steps of a
    branch, not branches of their own. \emph{Why:} otherwise the same piece
    of work is counted twice, and the metrics lie.
  \item[\Cref{rule:discussion}: traceable discussion.]
    Every comment, decision or piece of evidence is attached to one
    element of the graph (a line, a baseline or a branch), and the record
    is append-only. \emph{Why:} decisions made in chat threads and meetings
    cannot be audited; attached to the branch they concern, they are found
    exactly where the diagram is read.
\end{description}

\paragraph{HPB 1.0: what is normative.}
This section, together with the quick-reference card in \cref{app:card},
is the normative definition of HPB~1.0: the idea, the four branch types,
the seven rules and the notation. Everything else in this paper explains,
illustrates or supports it, and the formal model of \cref{app:formal}
states the same rules precisely. The method will change only through
versioned releases: HPB~1.x for additions that keep every existing diagram
valid, a new major version for anything that would not. The text of the
method and its notation are offered under the Creative Commons Attribution
4.0 license (CC~BY~4.0): copy, translate and adapt them, with credit. The
\LaTeX{} package \texttt{hpb.sty} and the HPB Markup toolkit are released
under the MIT license.

\paragraph{What you get.}
With these rules in place, the questions from \cref{sec:introduction} are
answered by looking at the picture: where a defect was seen (its fork),
whether it is fixed and in which version (its terminal), what else is open
on the product (the branches crossing NOW), what a planned version is
waiting for (the branches entering its gate), and what moved to a new
hardware reference (the line fork). The rest of the paper makes each of
these elements precise (\cref{sec:model}), fixes how they are drawn
(\cref{sec:notation}), and shows them in the situations a product manager
meets every week (\cref{sec:practice}).

% ---------------------------------------------------------------------
%  04-model.tex -- the HPB model in plain language
%  The precise definitions, rules R1--R7 and properties live in
%  appendix-a-formal.tex (labels def:*, rule:*, prop:*).
% ---------------------------------------------------------------------
\section{The HPB Model}
\label{sec:model}

HPB has very few concepts, on purpose. Every concept had to earn its place
by answering a question that I was asked in real product reviews and could
not answer from the tools I had: \emph{what did this version contain,
where was this defect seen, what is still open, and where did this
reference come from?} This section explains the concepts in plain
language, with the symbols introduced only where later sections use them.
The precise definitions, the rules R1--R7 in their exact form and the
properties that follow from them are collected in \cref{app:formal}; a
tool builder should read that appendix, a product manager can stop here.

\subsection{Product lines and baselines}
\label{sec:model-lines}

\paragraph{Product line.}
A \emph{product line} $P$ is one product reference followed over time: a
device model with its manufacturer reference, a firmware, an application
or a service (\cref{def:line}). This is narrower than the ``product line''
of software product line engineering, which means a family of variants; in
HPB a line is \emph{one} reference and its history. The portfolio is the
set of all lines, and lines are related by descent: when a hardware
revision creates a new reference, the new line records the old one as its
parent. The result is a family tree of references, the \emph{line
forest}.

\paragraph{Baseline.}
A \emph{baseline} $v$ is a named, identifiable state of a line:
\texttt{v2.3.1}, \texttt{Rev~B}, \texttt{PC-7} (\cref{def:baseline}). I
use the word in its configuration-management sense, a reference
configuration against which change is
defined~\citep{ieee828,iso10007}. Every baseline has
\begin{itemize}[leftmargin=1.5em,itemsep=1pt,topsep=2pt]
  \item a \emph{kind}: firmware, software, hardware, or \emph{other}, the
        last one used for \emph{process baselines}, i.e., a change to a
        manufacturing process, test procedure or supplier that changes the
        product as built without a new design revision;
  \item a \emph{status}: \emph{planned} (a target with a date) or
        \emph{released} (a fact with a date);
  \item a name that is unique within its line. Any versioning convention
        works, for example Semantic Versioning~\citep{semver}.
\end{itemize}
A status only moves from planned to released, and a released baseline is
immutable. That one constraint is what makes everything else in HPB
trustworthy: if a released version could still absorb work, no change log
and no defect count would ever be final.

\paragraph{What ``released'' means.}
Every time-based metric depends on it, so HPB fixes a default per kind.
Firmware and software are released when the build is published to
customers; hardware when its engineering change order is approved
\emph{and} its production effectivity is recorded, with the change-order
identifier; a process change when it becomes effective in production. You
may define it differently, but write your definition down.

\paragraph{New revision or new reference?}
This question comes back with every hardware revision, and HPB settles it
with the test that engineering change management already
uses~\citep{jarratt2011ec,stark2015plm}: \emph{form, fit and function}. A
revised unit that is not interchangeable with the previous one in form,
fit or function gets a new reference and starts a new line (a \emph{line
fork}, below). An interchangeable change, such as a component substitution
that keeps the reference, is simply a hardware baseline on the existing
line.

\paragraph{Release streams.}
Products are rarely maintained on a single sequence of versions. Each line
has a \emph{main stream} and, when needed, \emph{maintenance streams} that
leave a released baseline of another stream, for example a long-term
support branch of a firmware (\cref{def:mainline}). The ordered baselines
of the main stream form the \emph{main line} $M(P)$, the thick horizontal
stroke of every HPB diagram. Release dates order baselines only
\emph{within} a stream: a patch \texttt{v1.4.1} on a stream forked at
\texttt{v1.4} follows \texttt{v1.4}, even if it ships after \texttt{v2.0}
on the main stream, and it does not contain anything that only went into
\texttt{v2.0}. A stream never changes the product reference; that is what
line forks are for (\cref{fig:backport}).

\paragraph{Variants and configurations.}
Regional or certification variants, OEM builds and members of a product
family are \emph{sibling} lines: distinct references, each with its own
history. Work that applies to all of them belongs on the line of the shared
element, typically a firmware or software line, instead of being
duplicated. Each firmware or software baseline records the hardware
references it is released to run on (its \emph{applies-to} set), and a
\emph{configuration} is a combination of baselines that is actually
possible in the field, such as (\texttt{Rev~A}, firmware \texttt{v2.0},
app \texttt{5.1}) (\cref{def:config}). The applies-to relation is the
compatibility matrix of the portfolio.

For a single device with its own firmware, a \emph{combined} line with
both hardware revisions and firmware releases is simplest, and most
examples here use it. Switch to separate lines as soon as one firmware
build runs on several hardware references (\cref{fig:shared}).

\subsection{Branches}
\label{sec:model-branches}

A \emph{branch} $b$ is one unit of change intent on a product: an
improvement to make, a defect to fix, an engineering task to do, a lab
evolution to qualify (\cref{def:branch}). It is \emph{not} a code branch.
A single HPB branch may be implemented by many commits in several
repositories, by a laboratory campaign or by a hardware redesign. HPB
records where the intent started, where it ended, and what happened in
between, and nothing more.

\paragraph{Four types.}
Every branch has exactly one type $\tau(b)$, as in \cref{sec:onepage}:
$\mathsf{I}$ improvement (from customers, sales, competitor analysis or the
roadmap), $\mathsf{L}$ laboratory evolution (evaluation, qualification or
prototype work), $\mathsf{X}$ issue (a defect observed in a released
baseline) and $\mathsf{T}$ task (work that keeps the product operable, such
as a certification renewal or a last-time buy). In a portfolio with
constant rotation, the balance between making the product better
($\mathsf{I}$, $\mathsf{L}$) and keeping it correct ($\mathsf{X}$,
$\mathsf{T}$) is one of the first things management asks about, so the
notation puts the two groups on opposite sides of the main line.
Sub-tasks are not branches: they are \emph{steps} of the branch they
belong to, so that each piece of intent is counted once (\cref{rule:typed}).

\paragraph{The fork baseline.}
Every branch starts from exactly one released baseline, its \emph{fork
baseline} $f(b)$, fixed when the request is accepted at triage
(\cref{rule:origin}). For an issue, it is the version in which the defect
was \emph{observed}, as reported. For the other types, it is the version
the change is specified against, by default the latest release of the main
stream. When a defect was reported against a combination of baselines,
the full observed configuration is kept and one of its baselines is the
primary one, on whose line the branch is drawn.

The fork baseline records what was known at intake and is never
rewritten by what is learned later: other affected versions go to an
\emph{also affects} list, and the version that \emph{introduced} the
defect is recorded separately. I insist on this because the first thing
lost when tickets are edited over months is the one fact quality
reporting needs: where the problem was actually seen. An inferred fork
baseline, for example from old tickets loaded into HPB, is flagged as
such.

\paragraph{Steps.}
What happens inside a branch is a short, time-ordered list of
\emph{steps}: done, to do, test passed, test failed, each optionally with
evidence such as a report, a measurement or a build identifier. Steps are
where laboratory results and verifications live, and they are what the
release gate checks.

\subsection{How a branch lives and ends}
\label{sec:model-closing}

A branch is in one of four states (\cref{def:state},
\cref{fig:lifecycle}):
\begin{itemize}[leftmargin=1.5em,itemsep=1pt,topsep=2pt]
  \item \textsf{open}: accepted, not yet committed to a version;
  \item \textsf{targeted}: committed to a planned baseline, its
        \emph{target} $g(b)$. The target can be changed (re-target) or
        withdrawn (un-target), and the change remains visible;
  \item \textsf{merged}: its result is part of a released baseline,
        $m(b)$;
  \item \textsf{discarded}: closed without a baseline, with a written
        rationale.
\end{itemize}
Merged and discarded are final. Reopening a closed branch exists, but only
as an exceptional, audited operation.

% F7 -- Life cycle of a branch: the state machine of Definition "Branch state".
% Requires hpb.sty v1.0 (palette and glyph styles). Plain TikZ, units in cm.
\begin{figure}[t]
\centering
\resizebox{\linewidth}{!}{%
\begin{tikzpicture}[
    x=1cm, y=1cm,
    font=\sffamily\fontsize{7}{8.4}\selectfont,
    lc/state/.style   = {rectangle, rounded corners=5pt, draw=hpbInk,
                         line width=1pt, minimum width=3.3cm,
                         minimum height=1.45cm, inner sep=0pt},
    lc/final/.style   = {double=white, double distance=1.3pt, line width=0.8pt},
    lc/name/.style    = {font=\sffamily\bfseries\fontsize{8}{9.6}\selectfont,
                         text=hpbInk, anchor=west, inner sep=0pt},
    lc/inv/.style     = {font=\sffamily\fontsize{6.5}{7.8}\selectfont,
                         text=hpbMuted, anchor=west, inner sep=0pt},
    lc/arrow/.style   = {draw=hpbInk, line width=1pt, rounded corners=6pt,
                         -{Stealth[length=5pt,width=4.5pt]}},
    lc/audit/.style   = {lc/arrow, dash pattern=on 4pt off 2.5pt},
    lc/tr/.style      = {font=\sffamily\fontsize{7}{8.4}\selectfont,
                         text=hpbInk, fill=white, inner sep=1.5pt},
    lc/trb/.style     = {lc/tr, font=\sffamily\bfseries\fontsize{7}{8.4}\selectfont},
  ]
  % ---- the open set W(t): states in which a branch counts as work in progress
  % (dotted frame: deliberately not the style of an HPB gate band)
  \draw[draw=hpbDiscInk, line width=0.8pt, rounded corners=4pt,
        dash pattern=on 0.1pt off 2pt, line cap=round]
    (1.45,-1.02) rectangle (10.15,1.02);
  \node[lc/inv, anchor=south west, font=\sffamily\itshape\fontsize{6.5}{7.8}\selectfont]
    at (1.5,1.06) {open set $W(t)$};

  % ---- states: neutral fills (not the baseline color of the notation);
  %      the glyph in each box is the HPB terminal that represents the state
  \node[lc/state, fill=white]               (open) at (3.35,0)  {};
  \node[lc/state, fill=white]               (targ) at (8.25,0)  {};
  \node[lc/state, lc/final, fill=white]     (merg) at (13.65,0) {};
  \node[lc/state, lc/final, fill=white]     (disc) at (8.25,-3) {};

  % glyphs (shown for an improvement)
  \node[hpb/head ring, draw=hpbImp] at (2.2,0) {};
  \node[hpb/event, fill=white]      at (2.2,0) {};
  \draw[hpb/path, hpb/path target] (7.3,0) -- (7.83,0);
  \node[hpb/head ring, draw=hpbImp] at (7.1,0) {};
  \node[hpb/event, fill=white]      at (7.1,0) {};
  \node[hpb/event, fill=hpbImp]     at (12.5,0) {};
  \draw[draw=hpbInk, line width=1.3pt, line cap=round, line join=round]
    (12.5,0) ++(-2.6pt,0.1pt) -- ++(1.8pt,-2pt) -- ++(3.5pt,4.2pt);
  \node[hpb/event, fill=hpbDisc, draw=hpbMuted] at (7.1,-3) {};
  \draw[draw=hpbInk, line width=1.3pt, line cap=round]
    (7.1,-3) ++(-2.2pt,-2.2pt) -- ++(4.4pt,4.4pt) ++(-4.4pt,0) -- ++(4.4pt,-4.4pt);

  % names and invariants
  \node[lc/name] at (3.0,0.22)   {Open};
  \node[lc/inv]  at (3.0,-0.2)   {$g(b)$ unset};
  \node[lc/name] at (7.95,0.22)  {Targeted};
  \node[lc/inv]  at (7.95,-0.2)  {$g(b)$ planned;};
  \node[lc/inv]  at (7.95,-0.47) {verified at $t_r(b)$};
  \node[lc/name] at (13.4,0.22)  {Merged};
  \node[lc/inv]  at (13.4,-0.2)  {$m(b)=g(b)$};
  \node[lc/name] at (7.95,0.22-3) {Discarded};
  \node[lc/inv]  at (7.95,-0.2-3) {rationale $\rho(b)$};

  % ---- creation: a branch is opened by forking from a released baseline,
  %      drawn as the released-baseline glyph of the notation
  \node[hpb/baseline] (init) at (-0.35,0) {};
  \draw[lc/arrow] (init.east) -- (open.west)
    node[lc/trb, midway, above=2pt] {fork}
    node[lc/inv, midway, below=3pt, anchor=north] {from $f(b)$};

  % ---- ordinary transitions
  \draw[lc/arrow] ([yshift=3.5mm]open.east) -- ([yshift=3.5mm]targ.west)
    node[lc/trb, midway, above=1pt] {target};
  \draw[lc/arrow] ([yshift=-3.5mm]targ.west) -- ([yshift=-3.5mm]open.east)
    node[lc/trb, midway, below=1pt] {un-target};
  \draw[lc/arrow] (7.75,0.725) -- (7.75,1.3) -- (8.75,1.3) -- (8.75,0.725);
  \node[lc/trb, anchor=south] at (8.25,1.32) {re-target};
  \draw[lc/arrow] (targ.east) -- (merg.west)
    node[lc/trb, midway, above=2pt] {merge}
    node[lc/inv, midway, below=3pt, anchor=north, align=center]
      {on publication\\of $g(b)$, if verified};
  \draw[lc/arrow] ([xshift=6mm]open.south) |- (disc.west)
    node[lc/trb, pos=0.75, above=1pt] {discard};
  \draw[lc/arrow] (targ.south) -- (disc.north)
    node[lc/trb, midway, right=2pt] {discard};

  % ---- exceptional, audited transitions (dashed)
  \draw[lc/audit] (merg.north) -- (13.65,1.95) -- (3.95,1.95) -- (3.95,0.725);
  \node[lc/trb, anchor=south] at (8.8,1.97) {reopen\quad
    \textmd{\textcolor{hpbMuted}{(exceptional, audited)}}};
  \draw[lc/audit] (disc.south) -- (8.25,-4.15) -- (2.75,-4.15) -- (2.75,-0.725);
  \node[lc/trb, anchor=north] at (5.5,-4.17) {reopen\quad
    \textmd{\textcolor{hpbMuted}{(exceptional, audited)}}};
\end{tikzpicture}}
\caption{Life cycle of a branch (\cref{def:state}). A branch is
\textsf{open} until it is committed to a planned version (\textsf{targeted});
it merges only from \textsf{targeted}, when that version is published and
after it has been verified; it can be discarded from either state, always
with a reason. \textsf{Merged} and \textsf{discarded} are final; dashed
transitions are exceptional and audited. Each box shows the glyph of the
state; the dotted frame marks the states counted as work in progress,
$W(t)$.}
\label{fig:lifecycle}
\end{figure}


\paragraph{Merge.}
A branch is merged only from \textsf{targeted}, and only when its target
baseline is published (\cref{def:merge}). The relation ``branch $b$ is
closed by baseline $v$'' is what produces the change log of $v$. This
looks strict, and it is meant to be: every merge fulfills a commitment that
was written down beforehand, and released versions are never altered by
ordinary work. For an urgent fix, targeting, verification and publication
can follow each other within hours.

A targeted branch is \emph{verified} when it has a passed test step whose
evidence identifies a candidate of its target: a release-candidate build,
a prototype or a pilot lot. Verification is checked at the gate.

\paragraph{Discard.}
A discard closes a branch without a baseline and always carries a
rationale from one of two classes. \emph{Abandoned}: not viable,
duplicate, cannot reproduce, superseded, out of scope, will not fix.
\emph{Resolved without a product baseline}: already resolved in an
existing version, contained without design change (for example by
reworking stock), resolved by the supplier. Keeping the two classes apart
matters for quality reporting: a problem solved outside the product is not
the same as a problem nobody solved. A defect whose solution changes
\emph{how} the product is built rather than its design is not discarded;
it is merged into a process baseline.

\subsection{Line forks, gates and the life of a reference}
\label{sec:model-forks}

\paragraph{Line fork.}
A \emph{line fork} is a hardware baseline that creates a new reference
$P'$ (\cref{def:linefork}). The hardware baseline is the first baseline of
the new line, and the new line remembers its parent and the parent
baseline it was derived from. The branches that produce the new reference,
typically lab evolutions and improvements, fork from the parent line and
merge into that hardware baseline; later work forks from the new line.
Two consequences are worth stating because they are exactly the questions
that come up in practice. First, later versions of the parent do
\emph{not} contain the work that created the new reference. Second, a fix
made on the parent after the fork is \emph{not} in the child unless it is
carried over explicitly (\cref{fig:hardware-fork}).

\paragraph{Release gate.}
Every planned baseline is a \emph{gate} (\cref{def:gate}). Its
\emph{targeted set} $\mathcal{B}(g)$ is the set of branches currently
committed to it. Publishing the baseline is a single operation that
releases it and merges every branch of its targeted set. Before that,
every member must be verified, and anything it depends on must already be
released or be released in the same operation; whatever is not ready is
re-targeted or discarded first (\cref{rule:gate}). This is the
Stage-Gate idea~\citep{cooper1990stagegate} applied to versions instead of
projects, and it removes the most common source of misleading release
notes: work that was ``in'' the version on the plan but quietly did not
make it (\cref{fig:release-gate}).

\paragraph{Dependencies, propagations and launch bundles.}
Three optional relations cover the cases where versions and branches of
different lines are tied together (\cref{def:relations}). A
\emph{dependency} says that a branch can only be delivered with or after
another branch or baseline, for example a firmware improvement that needs
hardware \texttt{Rev~B}, or work blocked by a supplier sample. A
\emph{propagation} says that a branch carries the change of another branch
to a different stream or line: a back-port or forward-port. A change that
must reach several versions is therefore one branch per target, linked by
propagations, never one branch with several endings (\cref{rule:closure});
that is how a missing back-port becomes visible. A \emph{launch bundle}
is a set of planned baselines of different lines that are published
together or not at all, for example hardware \texttt{Rev~B} with firmware
\texttt{v2.0} and app \texttt{5.0}.

\paragraph{Life cycle of a reference.}
In a portfolio with high rotation, references are retired as often as
they are created, and the rules for what may still change on an old
reference are usually tribal knowledge. HPB makes them explicit. Every
line has up to four dated milestones: end of sale, end of production,
start of security-only maintenance and end of service
(\cref{def:lifecycle}). They divide its life into phases, and
\cref{tab:phases} in \cref{app:formal} says which branches and baselines
each phase still accepts (\cref{rule:lineage}). Supplier notices are a frequent trigger in
the later phases: a product change notice or an end-of-life notice for a
component usually opens a lab evolution (qualify an alternative) or a task
(last-time buy), targeted at a hardware baseline that is either
interchangeable or, if form, fit or function change, a line fork.


\paragraph{Deployment overlay.}
The questions a product manager is asked most often are about customers:
\emph{who is exposed to this defect, who has not upgraded, which promise
has slipped?} An optional \emph{deployment overlay} (\cref{def:deployment})
records the rollout status and installed base of each released baseline,
the configuration each customer runs, and the commitments made on each
branch. ``Customers exposed to X-118'' is then simply the set of customers
whose configuration contains a version in which X-118 occurs.

\subsection{What you get for free}
\label{sec:model-free}

The point of the rules is that the documents a product manager usually
compiles by hand become \emph{queries on the graph}. Four of them are
stated precisely, with proof sketches, in \cref{app:formal}:
\begin{itemize}[leftmargin=1.5em,itemsep=1pt,topsep=2pt]
  \item \textbf{Change log.} The content of a version is the set of
        branches it closed; grouped by type, it is the release note
        (\cref{prop:changelog}).
  \item \textbf{Escaped defects.} Every issue belongs to exactly one
        version: the one in which it was observed (\cref{prop:escaped}).
  \item \textbf{Work in progress.} At any moment, the open and targeted
        branches are a well-defined set, split by type, to which Little's
        law~\citep{little1961} applies (\cref{prop:wip}).
  \item \textbf{Lineage.} Every reference has a single ancestry and every
        version a single chain of predecessors, so ``is this fix in that
        version?'' always has a yes-or-no answer (\cref{prop:lineage}).
\end{itemize}
The last point has a practical corollary that I consider one of the most
useful things HPB shows. A fix merged on a maintenance stream, or on a
parent line after a line fork, is not contained in the main stream or the
child line. The diagram makes the question ``must this fix be carried
over?'' visible, and the rules force the answer to be recorded: either as
a propagated branch or as an explicit decision (\cref{fig:backport}).

% ---------------------------------------------------------------------
%  05-notation.tex -- the HPB notation standard (condensed)
%  Glyphs in tab:notation are drawn with hpb.sty v1.0 (same code as the
%  figures). The anatomy figure (fig:anatomy) is input in section 03.
% ---------------------------------------------------------------------
\section{The HPB Notation}
\label{sec:notation}

A methodology spreads when people can draw it on a whiteboard and others
read the drawing the same way. Git-flow~\citep{driessen2010gitflow} is a
good example: one diagram, copied everywhere, and every engineer knows how
to read it. HPB needs the same thing, so the notation is fixed: the canvas,
where each element goes, the glyphs and a short list of layout rules. All
diagrams in this paper are drawn with the \LaTeX{} package
\texttt{hpb.sty} (\cref{app:package}), and the open-source renderer of
\cref{sec:tooling} copies the same glyphs, so a diagram looks the same in
a paper, a wiki or a product portal. A quick-reference card is in
\cref{app:card}.

Four goals shaped the notation:
\begin{enumerate}[label=\textbf{G\arabic*},leftmargin=*,itemsep=1pt,topsep=3pt]
  \item \textbf{One concept, one glyph.} Every element a reader must
        recognize has exactly one glyph, and no glyph is decorative.
  \item \textbf{Position carries meaning.} \emph{When} and \emph{what kind
        of work} are read from horizontal and vertical position, before
        color or shape.
  \item \textbf{Never color alone.} Every category is encoded at least
        twice, so diagrams survive grayscale printing and color-vision
        deficiency.
  \item \textbf{Familiar to engineers.} A horizontal main line, version tags
        above the nodes, branches that leave and rejoin it: the vocabulary
        of the branch diagrams R\&D teams already read.
\end{enumerate}

\subsection{Canvas, lanes and reading order}
\label{sec:notation-canvas}

Time runs left to right. The vertical axis carries no quantity; it is
divided into horizontal \emph{lanes} of equal height. Lane~$0$ is the
\emph{main line}: a thick (3\,pt) dark stroke that \emph{is} the product.
Released baselines sit on it as nodes, in order; planned baselines sit to
the right of the present moment on a dashed continuation of the same
stroke. Each lane has a faint dotted guide across the canvas and a legend
box at its left end that names the type in words; the legend boxes and the
product box of the line align in one column, so the left margin reads as
the key.

\paragraph{Lane placement.}
Branches run on the lane of their type, by a fixed \emph{lane map}:
\begin{equation}
  \lambda(\mathsf{I}) = +1,\qquad
  \lambda(\mathsf{L}) = +2,\qquad
  \lambda(\mathsf{X}) = -1,\qquad
  \lambda(\mathsf{T}) = -2 .
  \label{eq:lanemap}
\end{equation}
Work that makes the product better (improvements, lab evolutions) goes
\emph{above} the main line; work that keeps it correct (issues, tasks)
goes \emph{below}. On each side, the type that most directly changes what
customers receive sits nearest the line. When several branches of a type
are open at the same time, extra lanes of that type are stacked outward
next to the existing ones, and the order $\mathsf{I}$ before $\mathsf{L}$
above and $\mathsf{X}$ before $\mathsf{T}$ below is kept
(\cref{fig:scale}). When a diagram shows several lines or streams, each
gets its own \emph{lane band}, separated by a thin rule, and above or below
is read relative to that band's main line
(\cref{fig:hardware-fork,fig:backport,fig:shared}).

\paragraph{Reading order.}
Every HPB diagram reads the same way: \emph{above the line, how the product
gets better; below it, what is wrong and how it is fixed; left of NOW,
history, which does not change; right of NOW, the plan, which is expected
to change.} The NOW marker is a dashed vertical line in its own color,
labeled \textsf{NOW}, across all lanes. Left of it all strokes are solid,
right of it all strokes are dashed: the drawing shows directly that
released versions are immutable. \Cref{fig:anatomy} annotates every glyph
of a single-line diagram.

\subsection{Ordinal time}
\label{sec:notation-time}

By default HPB uses an \emph{ordinal} time axis: every event (a release, a
branch start, a step, a terminal) gets its own column of equal width, in
time order, followed by the NOW column and then the planned baselines in
order of target date (\cref{def:ordinal}). Calendar months appear only as
tick labels on an axis below the diagram.

I chose ordinal time because real product histories are lumpy. A single
week can bring a hotfix, a release and three new issues, and then nothing
happens for a month. On a calendar scale the busy week becomes an
unreadable knot and the quiet month wastes half the page. With one column
per event, dense periods stay readable, glyphs never collide, and labels
can be placed by a fixed rule. The \emph{linear} alternative, where
horizontal position is proportional to calendar time as in a Gantt chart,
is still allowed when the distance to a date is itself the message, as in
a roadmap; it must then be stated in the caption.

\begin{remark}
  \label{rem:ordinal}
  In an ordinal diagram, horizontal length is \emph{not} duration. A long
  branch is one that was open while many other things happened, not
  necessarily one that was open for long. Durations such as lead time are
  computed from the model (\cref{sec:metrics}) and, when needed, printed
  as labels; they are never read off the drawing.
\end{remark}

\subsection{Glyphs}
\label{sec:notation-glyphs}

\Cref{tab:notation} lists the core glyphs and the model element each one
stands for; they are drawn with the same macros as the figures, so the
table is an exact sample of the notation. A few rules complete it.

\emph{Baselines.} When the target date of a planned baseline has passed, an
exclamation badge appears and its tag and date turn red. A hardware
baseline is a diamond with an indigo tag; with the caption \emph{new ref.}\
it is the first node of a new line, reached from the parent line by a
vertical lineage connector. A process baseline is a square. A maintenance
stream leaves its baseline with a rounded corner and continues in its own
band. Life-cycle milestones are thin labeled ticks, and a bar labeled
\textsf{EoS} ends a retired line. The dot at the left end of a main line
marks the creation of the reference; it is not a baseline.

\emph{Branches and terminals.} A branch is an orthogonal path with rounded
corners: a short run from the fork baseline, a drop (or rise) to the lane
of its type, and a run along the lane. Its label, above the start node,
gives the tracker reference and a short title. A discarded branch turns
gray, never rejoins the line and carries its rationale. A re-targeted
branch keeps its old target as a faded dashed continuation, so the change
of plan stays visible (\cref{fig:discard}).

\emph{Markers.} A step after closure, such as an effectiveness check on a
corrective action, sits on a dotted continuation (\cref{fig:rma}). A dark
badge under a baseline counts merged branches whose drawing was collapsed.
A dash-dotted frame around planned baselines of several lines is a launch
bundle. Thin straight connectors are the only lines that may cross lanes
diagonally: dotted with an open arrowhead for a dependency, dashed with a
filled arrowhead for a propagation.

% ---------------------------------------------------------------------
% Glyph table.  Glyph column: small hpbgraph pictures (hpb.sty v1.0).
% ---------------------------------------------------------------------
% \HPBglyph[<hpbgraph keys>]{<hpb macros>}: one glyph sample with 3pt of
% vertical padding (glyph sizes are normative; only spacing is reduced).
\providecommand{\HPBglyph}[2][]{%
  \begin{hpbgraph}[col=0.55cm,lane=0.9cm,legend=1.3cm,#1]#2%
    \path ([yshift=3pt]current bounding box.north)
          ([yshift=-3pt]current bounding box.south);%
  \end{hpbgraph}}

\begin{table}[p]
  \centering
  \small
  \caption{The HPB notation: glyph, element and meaning, drawn with
    \texttt{hpb.sty} at normative size. Color is never the only carrier of
    meaning (\cref{sec:notation-a11y}). The full card, with the macro that
    draws each glyph, is in \cref{app:card}.}
  \label{tab:notation}
  \renewcommand{\arraystretch}{1.12}
  \renewcommand{\tabularxcolumn}[1]{m{#1}}
  \begin{tabularx}{\linewidth}{@{}>{\centering\arraybackslash}m{3.55cm}
      >{\raggedright\arraybackslash}m{2.75cm}
      >{\raggedright\arraybackslash}X@{}}
    \toprule
    Glyph & Element & Meaning \\
    \midrule
    % main line
    \HPBglyph{\hpbmainline{0}{2.2}\hpbmainline*{2.2}{4.2}
      \draw[hpb/now] (2.2,0.35) -- (2.2,-0.35);}
    & Main line
    & The product line and its versions. Solid left of NOW (history),
      dashed right of NOW (plan). \\
    % origin
    \HPBglyph{\hpbmainline{0}{2}\hpborigin{0}{Product A}}
    & Line origin
    & Creation of the reference (not a version), with the product
      reference in a box aligned with the lane legends. \\
    % released baseline
    \HPBglyph{\hpbmainline{-1.2}{1.2}\hpbrelease{0}{v1.1}}
    & Released baseline
    & A released version; the tag holds its name. Immutable. \\
    % planned baseline
    \HPBglyph{\hpbmainline*{-1.8}{1.2}\hpbplanned{-0.7}{v1.2}{Mar 2027}}
    & Planned baseline
    & A planned version with its target date, and a gate. Badge ``!'',
      red tag and date: target date passed. \\
    % hardware baseline
    \HPBglyph{\hpbmainline{-1.8}{1.2}\hpbhardware{-0.7}{Rev B}}
    & Hardware baseline
    & A hardware revision; with the caption \emph{new ref.}, the first
      version of a new line (line fork). \\
    % gate band
    \HPBglyph[lane=0.75cm]{\hpbmainline*{1.3}{4.6}
      \node[hpb/planned] at (4,0) {};
      \hpbgate[0.8]{4}{0}{-1}{}
      \hpbterminal{target}{iss}{1.5}{-1}
      \hpbtarget{1.5}{-1}{4}}
    & Gate band
    & The targeted set $\mathcal{B}(g)$: every dashed continuation entering
      the band must be verified before the version ships
      (\cref{rule:gate}). \\
    % start nodes
    \HPBglyph[col=0.75cm]{\hpbstart{imp}{0}{0}{I}\hpbstart{lab}{1}{0}{L}
      \hpbstart{iss}{2}{0}{X}\hpbstart{task}{3}{0}{T}}
    & Start node
    & Opening of a branch; fill color and letter give its type. \\
    % branch path from a baseline
    \HPBglyph[lane=0.75cm]{\hpbmainline{-0.6}{3.8}
      \node[hpb/baseline] at (0,0) {};
      \hpbbranch{iss}{0}{-1}{1.3}{3.8}{none}{}}
    & Branch (fork)
    & The path leaves its fork baseline $f(b)$ and runs on the lane of its
      type. \\
    % merge
    \HPBglyph[lane=0.75cm]{\hpbmainline{-0.6}{4.6}
      \node[hpb/baseline] at (0,0) {};
      \node[hpb/baseline] at (4,0) {};
      \hpbbranch{imp}{0}{1}{1.2}{2.7}{merge}{4}}
    & Merge terminal
    & Merged: the path rejoins the version $m(b)$ that closed it. \\
    % discard
    \HPBglyph[lane=0.75cm]{\hpbmainline{-0.6}{3.8}
      \node[hpb/baseline] at (0,0) {};
      \hpbbranch{task}{0}{-1}{1.3}{3.3}{discard}{}}
    & Discard terminal
    & Discarded: closed without a version, with a mandatory rationale;
      gray path, never rejoins. \\
    % open head
    \HPBglyph[lane=0.75cm]{\hpbmainline{-0.6}{3.8}
      \node[hpb/baseline] at (0,0) {};
      \hpbbranch{lab}{0}{1}{1.3}{3.3}{open}{}
      \draw[hpb/now] (3.3,1.6) -- (3.3,0.4);}
    & Open branch
    & Open: the head lies on the NOW marker. \\
    % targeted
    \HPBglyph[lane=0.75cm]{\hpbmainline{-0.6}{2.7}\hpbmainline*{2.7}{5}
      \node[hpb/baseline] at (0,0) {};
      \node[hpb/planned] at (4.4,0) {};
      \hpbbranch{iss}{0}{-1}{1}{2.7}{target}{4.4}}
    & Targeted branch
    & Committed: dashed continuation from the head to the planned target
      $g(b)$. \\
    % steps
    \HPBglyph[col=0.6cm,lane=0.7cm]{\draw[hpb/path] (-0.4,0) -- (4.4,0);
      \hpbstep{imp}{0}{0}{done}\hpbstep{imp}{1.33}{0}{todo}
      \hpbstep{lab}{2.67}{0}{pass}\hpbstep{lab}{4}{0}{fail}}
    & Steps
    & Done, pending, test or verification passed, test failed. \\
    % NOW
    \HPBglyph[lane=0.7cm]{\hpbmainline{-1.6}{0}\hpbmainline*{0}{1.6}\hpbnow{0}{0}{0.2}}
    & NOW marker
    & The present moment: history to its left, plan to its right. \\
    \bottomrule
  \end{tabularx}
\end{table}

\subsection{Never color alone}
\label{sec:notation-a11y}

Color is lost in grayscale prints and photocopies, and many readers have
a color-vision deficiency. Color universal design therefore recommends
never using color as the only carrier of
information~\citep{okabe2008cud}. HPB encodes every category at least
twice:
\begin{itemize}[leftmargin=1.5em,itemsep=1pt,topsep=3pt]
  \item \emph{branch type}: lane position, letter in the start node, legend
        box in words, and fill color;
  \item \emph{branch state}: terminal shape (check, cross, hollow head),
        line style (solid, dashed, gray) and position relative to NOW;
  \item \emph{baseline status}: filled or white node, solid or dashed
        outline, presence of a target date, position relative to NOW;
  \item \emph{baseline kind}: shape (circle, diamond, square), tag color,
        and the caption \emph{new ref.}\ for a line fork;
  \item \emph{overdue}: the ``!'' badge in addition to the red tag;
  \item \emph{test outcome}: check or cross in addition to green or red.
\end{itemize}
The palette uses light fills with one dark ink color for every outline,
path and text. All fills keep a contrast ratio of at least $4.4{:}1$
against the ink, so letters, checks and crosses stay readable in
grayscale. Not every pair of fills is distinguishable without color
vision: task green and lab teal are close. That is acceptable only because
the two types are never adjacent (one is below the main line, the other
above it) and their letters and legend boxes differ. In the spirit of
Tufte~\citep{tufte2001}, there are no gridlines other than the faint lane
guides, no separate legend panel, and the thickest stroke is reserved for
the product itself.

\subsection{Layout rules}
\label{sec:notation-layout}

Glyphs alone do not make diagrams comparable; their arrangement must be
fixed too. A diagram \emph{conforms} to the HPB notation when it uses the
glyphs of \cref{tab:notation} and follows these rules. Every figure in
this paper conforms, \texttt{hpb.sty} implements the rules by construction
or by documented convention, and the renderer of \cref{sec:tooling}
applies them automatically, with the one gap noted there (the lane
preference of \ref{lay:onelane}).

\begin{enumerate}[label=\textbf{L\arabic*},ref=L\arabic*,leftmargin=*,
    itemsep=2pt,topsep=3pt]
  \item \label{lay:ordinal}\textbf{Ordinal columns.} Unless the caption
        says linear time: one event per column, equal column width, and one
        canvas (column width, lane height, key width) for all diagrams of a
        document. One exception is allowed: the merge terminals of the
        branches that a single release closes may share one column, just
        before that release, as long as they sit on different lanes. Calendar periods appear only as ticks below the diagram,
        and the axis may be left out when dates are not the message.
  \item \label{lay:lanes}\textbf{Fixed lanes.} Branches run on the lanes
        of \cref{eq:lanemap}; extra lanes of a type stack outward next to
        that type. Every lane has a guide and a legend box, aligned in the
        key column with the origin boxes of all lines.
  \item \label{lay:onelane}\textbf{One branch per lane at a time.} A lane
        is reused only after the previous branch on it has ended. Branches
        are placed in order of their fork. A branch already closed by NOW
        takes the innermost free lane of its type; a branch still open or
        targeted at NOW takes the outermost one.
  \item \label{lay:fork}\textbf{Clean verticals.} A branch leaves its fork
        baseline half a column to the right of the node and enters its
        merge baseline half a column to the left, so verticals never cross
        a tag. When several branches share a baseline and a side, the
        farther lane uses the smaller offset, so corners nest and paths
        never cross.
  \item \label{lay:tags}\textbf{One tag per baseline.} Directly above the
        node, with an arrow to it. Tags never overlap each other, a path or
        a label: make the column at least as wide as the widest tag, or
        insert a spacer column between consecutive baselines. In linear
        time, a colliding tag is raised one tag height.
  \item \label{lay:labels}\textbf{Labels.} A branch label (reference and
        short title) is left-aligned at the start node, above the lane, on
        a white halo, and ends before the next vertical and before NOW. The
        only text beside a terminal is the discard rationale. Nothing is
        written on a main line, and callouts never cover a path or a label
        other than the one they point to.
  \item \label{lay:terminal}\textbf{Exactly one terminal.} Every branch
        ends in exactly one merge, discard or open head, mirroring
        \cref{rule:closure}.
  \item \label{lay:now}\textbf{History is not redrawn.} Left of NOW,
        strokes are solid and drawn as recorded; right of NOW, only planned
        versions, targeted continuations and pending steps appear, all
        dashed or hollow. An audited reopen is a new event in its own
        column, never an edit of the past.
  \item \label{lay:linefork}\textbf{Several lines and streams.} Each child
        line, maintenance stream or line of a configuration gets its own
        lane band, separated by a thin rule, with its legend boxes in the
        key column. A child line starts at its hardware diamond, reached by
        the lineage connector; a stream leaves its fork baseline. A retired
        parent line ends with an \textsf{EoS} bar.
  \item \label{lay:fonts}\textbf{Legibility.} Scale diagrams uniformly,
        never in one direction only. After scaling, no text is smaller than
        6\,pt and glyphs keep their relative sizes.
\end{enumerate}

Rules \ref{lay:ordinal}--\ref{lay:onelane} also make the layout
\emph{computable}: columns come from sorting the events, lanes from
packing the branch intervals of each type. That is the basis of the
layout algorithm in \cref{app:layout}, and the reason a text description
of a product history is enough to produce a conforming diagram.

% ---------------------------------------------------------------------
%  06-practice.tex -- HPB in practice: canonical situations
%  Every condition figure is \input exactly once, next to its narrative.
%  Each situation follows the same pattern:
%    situation -> how HPB draws it -> decision it supports.
% ---------------------------------------------------------------------
\section{HPB in Practice: Canonical Situations}
\label{sec:practice}

A notation only becomes a shared language when the situations that keep
coming back are always drawn the same way. Running a portfolio of
hardware, firmware and software products, I found that a small number of
situations account for almost every planning and traceability question I
had to answer. This section walks through nine of them. Together they
exercise every rule R1--R7.

Each situation is told in three short steps: \emph{the situation} as a
product manager meets it, \emph{how HPB draws it}, and \emph{the decision}
the diagram supports. All products, items and dates are fictitious and
illustrative, and each figure stands on its own. Time runs left to right
in \emph{events}, not days (\cref{sec:notation-time}).
\Cref{tab:conditions} is the map of the section.

Read every diagram in the same order: the main line from left to right,
then above it (how the product gets better), then below it (what is wrong
and how it gets fixed), then the NOW marker, right of which everything is
plan. With several lines or streams, do this once per band. Then ask of
any branch the three questions HPB exists to answer: \emph{Where did it
start? What else was open at the same time? Which version closed it, or
why was it dropped?}

\begin{table}[t]
  \centering
  \small
  \caption{The nine canonical situations and the decision each diagram
    supports. All examples are illustrative.}
  \label{tab:conditions}
  \renewcommand{\arraystretch}{1.1}
  \begin{tabularx}{\linewidth}{@{}>{\raggedright\arraybackslash}p{3.6cm}
      >{\centering\arraybackslash}p{0.7cm}
      >{\raggedright\arraybackslash}X
      >{\raggedright\arraybackslash}p{1.2cm}@{}}
    \toprule
    \textbf{Situation} & \textbf{Fig.} & \textbf{Decision supported}
      & \textbf{Rules} \\
    \midrule
    A. Firmware or software line
      & \ref{fig:software-line}
      & Which version contains a change; in which version a defect was
        seen; release notes per version.
      & R1--R3 \\
    B. Release gate
      & \ref{fig:release-gate}
      & Can the planned version ship? Release, slip or descope.
      & R4 \\
    C. Hardware revision, new reference
      & \ref{fig:hardware-fork}
      & What goes into the new reference; how long the old one is
        maintained; which fixes must be carried over.
      & R5, R3 \\
    D. Discard and re-target
      & \ref{fig:discard}
      & Dropping work explicitly; moving a commitment without losing its
        history.
      & R2, R7 \\
    E. Many branches at once
      & \ref{fig:scale}
      & Which long-lived open work needs a decision; how much is in
        flight per type.
      & R6 \\
    F. Kanban board versus HPB
      & \ref{fig:same-data}
      & What a team board can and cannot tell you about the product.
      & R1, R2 \\
    G. Maintained versions, back-port
      & \ref{fig:backport}
      & Which maintained versions contain a fix; which fixes still have to
        be ported.
      & R2, R3 \\
    H. Field return (RMA)
      & \ref{fig:rma}
      & How a field failure was solved without a new design; whether the
        action worked; supplier-driven change.
      & R2, R7 \\
    I. Shared firmware, joint launch
      & \ref{fig:shared}
      & Which configurations a defect affects; whether a coordinated
        launch can go ahead.
      & R4, R1 \\
    \bottomrule
  \end{tabularx}
\end{table}

% =====================================================================
\subsection{A. A firmware or software line}
\label{sec:cond-software}

% F2 -- Condition A: a firmware/software product line (history only).
% Illustrative data. Requires hpb.sty v1.1. Canvas: document defaults.
\begin{figure}[t]
\centering
\hpbfit{%
\begin{hpbgraph}
  % lanes: improvements above, issues and tasks below the main line
  \hpblane{1}{imp}{Improvement (I)}{0}{16.7}
  \hpblane{-1}{iss}{Issue (X)}{0}{16.7}
  \hpblane{-2}{task}{Task (T)}{0}{16.7}
  % main line and baselines
  \hpbmainline{0}{16.7}
  \hpborigin{0}{Product A}
  \hpbrelease{1}{v1.0}
  \hpbrelease{5}{v1.0.1}
  \hpbrelease{10}{v1.1}
  \hpbrelease{16}{v1.2}
  % branches (fork col, lane, start col, end col, end, merge col)
  \hpbbranch[label={X-101 Boot loop}]{iss}{1}{-1}{2}{4}{merge}{5}
  \hpbbranch[forkoff=0.6, mergeoff=0.6, label={I-21 Energy-saving mode}]{imp}{1}{1}{3}{9}{merge}{10}
  \hpbbranch[forkoff=0.6, label={X-102 Log time zone}]{iss}{5}{-1}{6}{8}{merge}{10}
  \hpbbranch[forkoff=0.3, mergeoff=0.3, label={T-8 Renew signing certificate}]{task}{5}{-2}{7}{14}{merge}{16}
  \hpbbranch[forkoff=0.6, mergeoff=0.6, label={I-22 Remote diagnostics}]{imp}{10}{1}{11}{15}{merge}{16}
  \hpbbranch[mergeoff=0.6, label={X-103 Wi-Fi reconnect}]{iss}{10}{-1}{12}{13}{merge}{16}
  % ordinal time axis with calendar ticks (month of each baseline)
  \hpbaxis{0}{16.7}{-2.75}
  \hpbtick{1}{Jan 2026}{-2.75}
  \hpbtick{5}{Feb}{-2.75}
  \hpbtick{10}{Apr}{-2.75}
  \hpbtick{16}{Jul}{-2.75}
\end{hpbgraph}}
\caption{Situation~A: a firmware or software line, history
only (illustrative). Each branch shows the version it started from and the
version that closed it; X-101 is a hotfix closed by \texttt{v1.0.1}. Time is
ordinal: one column per event, with months only as ticks.}
\label{fig:software-line}
\end{figure}


\paragraph{Situation.}
One line, released in successive versions, with improvements, defects and
tasks running in parallel. A customer reports a boot loop: is it fixed, and
since when? Sales asks what went into \texttt{v1.1}. A backlog knows the
items, not the versions around them.

\paragraph{In HPB.}
\Cref{fig:software-line} is pure history: four releases of
\emph{Product~A}, no NOW marker. I-21 forks from \texttt{v1.0} and merges
into \texttt{v1.1}; its path passes over the patch \texttt{v1.0.1}, so the
patch visibly does not contain it. X-101 \emph{Boot loop} is seen in
\texttt{v1.0} and closed by the hotfix \texttt{v1.0.1}. T-8 stays open
across \texttt{v1.1} and closes with \texttt{v1.2}.

\paragraph{Decision.}
Every history question is one path. The boot loop is fixed in
\texttt{v1.0.1} and in everything to its right. Release notes are the
rises into each version; escaped defects are the drops below it. R1--R3
can be checked by eye.

% =====================================================================
\subsection{B. Planning with a release gate}
\label{sec:cond-gate}

% F3 -- Condition B: planning with a release gate.
% Illustrative data. Requires hpb.sty v1.1. Canvas: document defaults.
% Lanes follow Algorithm 1: closed branches take the innermost free lane,
% branches open at NOW the outermost one.
\begin{figure}[t]
\centering
\hpbfit{%
\begin{hpbgraph}
  % lanes: improvements above, issues below (three issue lanes: lane packing)
  \hpblane{1}{imp}{Improvement (I)}{0}{17.7}
  \hpblane{-1}{iss}{Issue (X)}{0}{17.7}
  \hpblane{-2}{iss}{Issue (X)}{0}{17.7}
  \hpblane{-3}{iss}{Issue (X)}{0}{17.7}
  % main line: history (solid) up to NOW, plan (dashed) afterwards
  \hpbmainline{0}{11}
  \hpbmainline*{11}{17.7}
  \hpborigin{0}{Product A}
  \hpbrelease{1}{v2.0}
  \hpbrelease{6}{v2.0.1}
  \hpbplanned{13}{v2.1}{May 2027}
  \hpbplanned{17}{v2.2}{Sep 2027}
  % history: a hotfix already merged (closed: innermost lane)
  \hpbbranch[forkoff=0.6, label={X-210 Sync error}]{iss}{1}{-1}{2}{5}{merge}{6}
  % open but untargeted issue (open at NOW: outermost lane); no gate
  \hpbbranch[forkoff=0.3, label={X-213 Display flicker (untargeted)}]{iss}{1}{-3}{4}{11}{open}{}
  % targeted branches: open at NOW, dashed continuation into the gate of v2.1
  % (nested rises: farther lane, smaller offset; offsets 0.2 column apart)
  \hpbbranch[forkoff=0.5, mergeoff=0.8, label={I-50 Multi-user access}]{imp}{1}{1}{3}{11}{target}{13}
  \hpbbranch[forkoff=0.3, mergeoff=0.4, label={X-211 Report totals}]{iss}{6}{-2}{7}{11}{target}{13}
  \hpbbranch[forkoff=0.6, mergeoff=0.6, label={X-212 CSV encoding}]{iss}{6}{-1}{8}{11}{target}{13}
  % steps: a done sub-task on I-50; X-212 verified in release candidate RC1
  \hpbstep{imp}{10}{1}{done}
  \hpbstep{iss}{9}{-1}{pass}
  % gate around the planned baseline v2.1 and NOW marker
  \hpbgate{13}{1}{-2}{gate v2.1}
  \hpbnow{11}{1}{-3}
  \hpbcallout{6.3,2.35}{v2.1 cannot be published until I-50 and X-211 are verified\\
    (as X-212 is), re-targeted or discarded (rule R4)}{12.6,1.62}
  % ordinal time axis with calendar ticks (history part)
  \hpbaxis{0}{17.7}{-3.75}
  \hpbtick{1}{Jan 2027}{-3.75}
  \hpbtick{6}{Feb}{-3.75}
  \hpbtick{11}{Apr}{-3.75}
\end{hpbgraph}}
\caption{Situation~B: a release gate (illustrative). Three
branches are committed to \texttt{v2.1}, and their dashed lines enter its
gate; only X-212 has a passed verification. By \cref{rule:gate},
\texttt{v2.1} cannot ship until I-50 and X-211 are verified, moved or
discarded. X-213 is open but committed to nothing.}
\label{fig:release-gate}
\end{figure}


\paragraph{Situation.}
Three things are promised for a release; two are not finished. Too often
this is discovered in release week, and the choice to ship, slip or cut is
made in a corridor and never written down.

\paragraph{In HPB.}
In \cref{fig:release-gate}, three open branches are targeted at the
planned \texttt{v2.1}: each stops at NOW and continues as a dashed line
into the shaded gate. Only X-212 carries a passed test on a release
candidate. X-213 is open with no target; its path just ends at NOW.

\paragraph{Decision.}
The diagram \emph{is} the agenda of the gate review. By \cref{rule:gate},
\texttt{v2.1} cannot ship while I-50 and X-211 are unverified; each one is
verified, moved to \texttt{v2.2} or discarded with a reason, and the
outcome (release, slip or descope) is recorded. X-213 goes to triage
before it becomes an eternal branch.

% =====================================================================
\subsection{C. A hardware revision that creates a new reference}
\label{sec:cond-hardware}

% F4 -- Condition C: hardware revision and line fork.
% Illustrative data. Requires hpb.sty v1.1. Canvas: document defaults.
% The hardware baseline Rev B is the FIRST node of the child line (drawn
% above its parent, in its own lane band); the lineage connector leaves the
% parent main line after v1.3, the baseline Rev B was derived from.
\begin{figure}[t]
\centering
\hpbfit{%
\begin{hpbgraph}
  % ---- child line DMS-4000 Rev B (y = 3.5), its own band and lane
  \hpblane{4.5}{imp}{Improvement (I)}{0}{19.7}
  \hpblane{3.5}{main}{DMS-4000\\Rev B}{0}{9}
  \hpbmainline{9}{19.7}[3.5]
  \hpbseparator{2.85}{-2.4}{19.7}
  % ---- parent line DMS-4000 Rev A (y = 0) and its lanes
  \hpblane{2}{lab}{Lab (L)}{0}{19.7}
  \hpblane{1}{imp}{Improvement (I)}{0}{19.7}
  \hpblane{-1}{iss}{Issue (X)}{0}{19.7}
  \hpbmainline{0}{19}
  \hpborigin{0}{DMS-4000\\Rev A}
  \hpbrelease{1}{v1.2}
  \hpbrelease{5}{v1.3}
  \hpbmilestone{10}{end of sale}
  \hpbrelease{11}{v1.4}
  \hpbrelease{16}{v1.4.1}
  \hpbeos{19}
  % ---- line fork: lineage connector from the parent main line (after v1.3)
  %      to the hardware baseline Rev B, first baseline of the child line
  \hpblinefork{9}{0}{3.5}
  \hpbhardware{9}{Rev B}[3.5]
  \hpbrelease{12}{v2.0}[3.5]
  \hpbrelease{18}{v2.1}[3.5]
  % lab evolution with test steps, and an improvement, both forked from v1.2
  % and merged into Rev B (cross-line merge; nested rises)
  \hpbbranch[forkoff=0.5, mergeoff=0.6, mergeline=3.5,
    label={L-5 Sensor board v2}]{lab}{1}{2}{2}{7}{merge}{9}
  \hpbbranch[forkoff=0.8, mergeoff=0.4, mergeline=3.5,
    label={I-30 Low-power mode}]{imp}{1}{1}{3}{8}{merge}{9}
  \hpbstep{lab}{4}{2}{fail}
  \hpbstep{lab}{6}{2}{pass}
  % work on the child line
  \hpbbranch[line=3.5, label={I-40 Remote unlock}]{imp}{12}{4.5}{13}{17}{merge}{18}
  % Rev A in maintenance: a defect observed in v1.4, fixed by firmware v1.4.1
  \hpbbranch[label={X-118 Relay chatter}]{iss}{11}{-1}{12}{14}{merge}{16}
  % ---- annotation
  \hpbcallout{12.6,1.6}{line fork (R5): parent(Rev B) = Rev A,\\
    Rev B derived from v1.3}{9,1.6}
  % ---- ordinal time axis with calendar ticks
  \hpbaxis{0}{19.7}{-1.85}
  \hpbtick{1}{Jan 2026}{-1.85}
  \hpbtick{5}{Mar}{-1.85}
  \hpbtick{9}{Jun}{-1.85}
  \hpbtick{16}{Sep}{-1.85}
  \hpbtick{19}{Dec}{-1.85}
\end{hpbgraph}}
\caption{Situation~C: a hardware revision that creates a new
reference (illustrative). L-5 and I-30 merge into the diamond Rev~B, which
starts the child line \emph{DMS-4000 Rev~B}; the lineage connector leaves the
parent after \texttt{v1.3}, the version Rev~B derives from
(\cref{rule:lineage}). The parent keeps its own life: end of sale, the fix
for X-118 in \texttt{v1.4.1}, and end of service (EoS).}
\label{fig:hardware-fork}
\end{figure}


\paragraph{Situation.}
A hardware change fails the interchangeability test, so manufacturing
needs a new reference. From then on there are two physical products in
the field and every question doubles. This is where the methods of
\cref{sec:gaps} leave the widest gap.

\paragraph{In HPB.}
In \cref{fig:hardware-fork}, the lab branch L-5 (one failed and one passed
test) and the improvement I-30 merge into the hardware diamond
\emph{Rev~B}. Rev~B starts a child line in its own band, linked to the
parent after \texttt{v1.3}. New work lives on the child; the parent passes
end of sale, ships the fix for X-118 in \texttt{v1.4.1} and ends with the
\textsf{EoS} bar.

\paragraph{Decision.}
What the new reference contains (L-5 and I-30, backed by the passed test),
how long the old one lives (its milestones and the work still forking
from it), and which fixes must be carried over: the fix for X-118 is
\emph{not} in Rev~B, so its status there must be recorded, and ported if
needed (situation~G).

% =====================================================================
\subsection{D. Discarded work and re-targeting}
\label{sec:cond-discard}

% F5 -- Condition D: discarded and not-viable work, re-targeting.
% Illustrative data. Requires hpb.sty v1.1. Canvas: document defaults.
% Lanes follow Algorithm 1 (closed: innermost free lane; open at NOW:
% outermost free lane).
\begin{figure}[t]
\centering
\hpbfit{%
\begin{hpbgraph}
  % lanes: improvements above; two issue lanes below (parallel defects)
  \hpblane{1}{imp}{Improvement (I)}{0}{18.7}
  \hpblane{-1}{iss}{Issue (X)}{0}{18.7}
  \hpblane{-2}{iss}{Issue (X)}{0}{18.7}
  % main line: history up to NOW (col 13), plan afterwards
  \hpbmainline{0}{13}
  \hpbmainline*{13}{18.7}
  \hpborigin{0}{Product A}
  \hpbrelease{1}{v3.0}
  \hpbrelease{8}{v3.0.1}
  \hpbplanned{15}{v3.1}{Mar 2027}
  \hpbplanned{18}{v3.2}{Jun 2027}
  % improvement judged not viable: discarded, never rejoins the line
  \hpbbranch[label={I-60 Solar panel}]{imp}{1}{1}{2}{5}{discard}{}
  \hpblabel[right]{5}{1}{\itshape not viable\\\itshape (unit cost)}
  % defect fixed by the patch baseline v3.0.1
  \hpbbranch[forkoff=0.3, label={X-52 Display flicker}]{iss}{1}{-2}{3}{7}{merge}{8}
  % same defect reported again: closed as duplicate
  \hpbbranch[forkoff=0.6, label={X-57 Screen flickers}]{iss}{1}{-1}{4}{5}{discard}{}
  \hpblabel[right]{5}{-1}{\itshape duplicate\\\itshape of X-52}
  % improvement targeted at v3.1
  \hpbbranch[label={I-61 Offline mode}]{imp}{8}{1}{9}{13}{target}{15}
  % defect re-targeted from v3.1 to v3.2: the decision is a recorded (done)
  % step on the branch with a label; former target faded, current target at
  % full strength
  \hpbbranch[label={X-63 OTA timeout}]{iss}{8}{-2}{10}{13}{target}{18}
  \hpbtarget[faded]{13}{-2}{15}
  \hpbstep{iss}{11}{-2}{done}
  \hpblabel[below center]{11}{-2}{re-targeted\\[-1pt]v3.1\,$\rightarrow$\,v3.2}
  % NOW marker
  \hpbnow{13}{1}{-2}
  % annotation of the former target
  \hpbcallout{16.0,-1.0}{former target\\(faded)}{14.5,-1.0}
  % ordinal time axis with calendar ticks
  \hpbaxis{0}{18.7}{-3.1}
  \hpbtick{1}{Jul 2026}{-3.1}
  \hpbtick{8}{Sep}{-3.1}
  \hpbtick{13}{Nov}{-3.1}
  \hpbtick{15}{Mar 2027}{-3.1}
  \hpbtick{18}{Jun}{-3.1}
\end{hpbgraph}}
\caption{Situation~D: discards and re-targeting (illustrative). I-60
and the duplicate X-57 end in gray crosses with their reasons and never
rejoin the line. X-63 was moved from \texttt{v3.1} to \texttt{v3.2}; the old
commitment stays visible as a faded dashed line.}
\label{fig:discard}
\end{figure}


\paragraph{Situation.}
A feature turns out too expensive, two reports describe the same defect,
an issue is pushed to the next release. In list tools all three end as
``closed'' or as an edited date, and six months later nobody knows why.

\paragraph{In HPB.}
In \cref{fig:discard}, I-60 is discarded: its path turns gray, never
returns to the line and carries the reason \emph{not viable (unit cost)}.
X-57 is discarded as \emph{duplicate of X-52}. X-63 was targeted at
\texttt{v3.1} and moved to \texttt{v3.2}; the old target remains as a
faded rise.

\paragraph{Decision.}
Decisions \emph{not} to deliver are records, not deletions
(\cref{rule:closure,rule:discussion}). Discards can be counted by reason,
and the faded rise shows which commitment moved, which is what an honest
gate-slip report needs (\cref{tab:metrics}).

% =====================================================================
\subsection{E. Many branches at once}
\label{sec:cond-scale}

% F6 -- Condition E: readability at scale (collapsed history, lane packing).
% Illustrative data. Requires hpb.sty v1.1. Canvas: document defaults.
% Lanes follow Algorithm 1 (closed: innermost free lane; open at NOW:
% outermost free lane).
\begin{figure}[t]
\centering
\hpbfit{%
\begin{hpbgraph}
  % two lanes per type; older open branches occupy the outer lanes
  \hpblane{4}{lab}{Lab (L)}{0}{19.7}
  \hpblane{3}{lab}{Lab (L)}{0}{19.7}
  \hpblane{2}{imp}{Improvement (I)}{0}{19.7}
  \hpblane{1}{imp}{Improvement (I)}{0}{19.7}
  \hpblane{-1}{iss}{Issue (X)}{0}{19.7}
  \hpblane{-2}{iss}{Issue (X)}{0}{19.7}
  \hpblane{-3}{task}{Task (T)}{0}{19.7}
  \hpblane{-4}{task}{Task (T)}{0}{19.7}
  % main line; history before v2.2 is collapsed into count badges
  \hpbmainline{0}{19}
  \hpbmainline*{19}{19.7}
  \hpborigin{0}{Product B}
  \hpbrelease{1}{v1.0}  \hpbcollapsed{1}{14}
  \hpbrelease{3}{v1.1}  \hpbcollapsed{3}{9}
  \hpbrelease{4}{v2.0}  \hpbcollapsed{4}{21}
  \hpbrelease{8}{v2.1}  \hpbcollapsed{8}{12}
  \hpbrelease{15}{v2.2}
  % long-lived open branches (outer lanes)
  \hpbbranch[label={T-14 EMC re-test}]{task}{1}{-4}{2}{19}{open}{}
  \hpbbranch[forkoff=0.6, label={L-5 Radio module}]{lab}{4}{4}{5}{19}{open}{}
  \hpbbranch[forkoff=0.4, label={T-20 Field logs}]{task}{4}{-3}{6}{19}{open}{}
  \hpbbranch[forkoff=0.7, label={X-88 Rare reboot}]{iss}{4}{-2}{7}{19}{open}{}
  % branches forked from v2.1 (nested rises: farther lane, smaller offset;
  % offsets at least 0.25 column apart and clear of the v2.1 tag)
  \hpbbranch[forkoff=0.55, label={L-7 Antenna tuning}]{lab}{8}{3}{9}{19}{open}{}
  \hpbbranch[forkoff=0.8, label={I-42 Fleet dashboard}]{imp}{8}{2}{10}{19}{open}{}
  \hpbbranch[forkoff=1.05, label={I-41 Quick setup}]{imp}{8}{1}{11}{13}{merge}{15}
  \hpbbranch[label={X-90 Clock drift}]{iss}{8}{-1}{12}{14}{merge}{15}
  % lanes reused after the merges into v2.2
  \hpbbranch[forkoff=0.6, label={I-45 Voice UI}]{imp}{15}{1}{16}{19}{open}{}
  \hpbbranch[label={X-95 Pairing}]{iss}{15}{-1}{17}{19}{open}{}
  % NOW marker (column 18 is left free so that labels clear the open heads)
  \hpbnow{19}{4}{-4}
  % callouts (boxes and leaders clear of every path and label)
  \hpbcallout{2.75,-2.3}{collapsed history:\\badge = merged\\branches hidden}{3,-0.62}
  \hpbcallout{2.35,3.2}{long-lived open\\branch, outer lane}{4.6,3.2}
  \hpbcallout{14.6,-4.85}{8 open branches at NOW on 8 lanes\\(= maximum overlap per type)}{19,-4.45}
\end{hpbgraph}}
\caption{Situation~E: many branches at once (illustrative). Older
releases are collapsed into count badges. Each type gets as many lanes as it
has overlapping branches, and freed lanes are reused: I-45 and X-95 take the
lanes of I-41 and X-90. Long-lived open work sits on the outer lanes; the
length of an open path counts the releases it has outlived, not days.}
\label{fig:scale}
\end{figure}


\paragraph{Situation.}
A mature product with years of history and many open branches. Draw
everything and the diagram becomes wallpaper. In a fast-rotating
portfolio this is the normal case.

\paragraph{In HPB.}
Two devices keep \cref{fig:scale} readable. Old releases are collapsed
into count badges, and lanes are reused: each type gets as many lanes as
it has branches open at the same time, stacked outward. Long-lived open
branches take the outer lanes, so newer forks reach the inner ones without
crossing them.

\paragraph{Decision.}
Which open work needs attention. Because time is counted in events, the
length of an open path shows how many releases it has outlived: T-14 has
survived four and goes on the next triage agenda. The lanes occupied at
NOW are the work in progress per type.

% =====================================================================
\subsection{F. The same items on a board and in HPB}
\label{sec:cond-board}

% F8 -- The same portfolio seen in a Kanban board and in HPB.
% Illustrative data. Requires hpb.sty v1.1. The board (a) is drawn with plain
% TikZ in explicit cm; the HPB view (b) uses the hpb macros and the document
% canvas. The panels are stacked so that both keep their natural size.
\begin{figure}[t]
\centering
\begin{tikzpicture}[
    kb/col/.style   = {draw=hpbDiscInk, line width=0.6pt, rounded corners=3pt,
                       fill=hpbGuide!30},
    kb/head/.style  = {font=\sffamily\bfseries\fontsize{7}{8.4}\selectfont,
                       text=hpbInk, anchor=north},
    kb/card/.style  = {draw=hpbInk, line width=0.8pt, rounded corners=2pt,
                       fill=white},
    kb/ref/.style   = {font=\sffamily\bfseries\fontsize{7}{8.4}\selectfont,
                       text=hpbInk, anchor=west, inner sep=0pt},
    kb/title/.style = {font=\sffamily\fontsize{6.5}{7.8}\selectfont,
                       text=hpbInk, anchor=west, inner sep=0pt},
    kb/panel/.style = {font=\sffamily\bfseries\fontsize{7.5}{9}\selectfont,
                       text=hpbInk, anchor=south west, inner sep=0pt},
  ]
  \node[kb/panel] at (0,1.95) {(a) Kanban board, Product A};
  % columns: x of left edge (cm), title
  \foreach \x/\t in {0/To do, 3.6/In progress, 7.2/Done}{
    \draw[kb/col] (\x,1.75) rectangle (\x+3.4,-1.0);
    \node[kb/head] at (\x+1.7,1.68) {\t};
    \draw[draw=hpbDiscInk, line width=0.6pt] (\x,1.3) -- (\x+3.4,1.3);
  }
  % cards: column left edge, top y (cm), type color, letter, ref, title
  \foreach \x/\y/\c/\l/\r/\t in {
      0/1.15/hpbIss/X/X-35/Drift,
      3.6/1.15/hpbLab/L/L-3/New sensor,
      3.6/0.35/hpbImp/I/I-14/Eco mode,
      7.2/1.15/hpbIss/X/X-31/Crash,
      7.2/0.35/hpbImp/I/I-12/Faster pairing,
      7.2/-0.45/hpbTask/T/T-7/Migrate logs}{
    \draw[kb/card] (\x+0.1,\y) rectangle (\x+3.3,\y-0.62);
    \node[hpb/event, fill=\c, minimum size=10pt, line width=1.1pt,
          font=\sffamily\bfseries\fontsize{6}{7}\selectfont]
      at (\x+0.38,\y-0.31) {\l};
    \node[kb/ref]   at (\x+0.65,\y-0.31) {\r};
    \node[kb/title] at (\x+1.35,\y-0.31) {\t};
  }
\end{tikzpicture}

\medskip
\hpbfit{%
\begin{hpbgraph}
  \node[font=\sffamily\bfseries\fontsize{7.5}{9}\selectfont, text=hpbInk,
        anchor=south west, inner sep=0pt] at (-2.75,3.0)
    {(b) HPB view of the same items};
  \hpblane{2}{lab}{Lab (L)}{0}{17.7}
  \hpblane{1}{imp}{Improvement (I)}{0}{17.7}
  \hpblane{-1}{iss}{Issue (X)}{0}{17.7}
  \hpblane{-2}{task}{Task (T)}{0}{17.7}
  \hpbmainline{0}{14}
  \hpbmainline*{14}{17.7}
  \hpborigin{0}{Product A}
  \hpbrelease{1}{v1.0}
  \hpbrelease{10}{v1.1}
  \hpbplanned{16}{v1.2}{Jun 2027}
  % history: two merges into v1.1, one discard
  \hpbbranch[forkoff=0.6, label={X-31 Crash}]{iss}{1}{-1}{2}{7}{merge}{10}
  \hpbbranch[forkoff=0.8, mergeoff=0.6, label={I-12 Faster pairing}]{imp}{1}{1}{3}{9}{merge}{10}
  \hpbbranch[forkoff=0.3, label={T-7 Migrate logs}]{task}{1}{-2}{4}{6}{discard}{}
  % open work at NOW
  \hpbbranch[forkoff=0.5, label={L-3 New sensor}]{lab}{1}{2}{5}{14}{open}{}
  \hpbstep{lab}{8}{2}{pass}
  \hpbbranch[forkoff=0.6, label={I-14 Eco mode}]{imp}{10}{1}{11}{14}{target}{16}
  \hpbbranch[label={X-35 Drift}]{iss}{10}{-1}{12}{14}{target}{16}
  \hpbnow{14}{2}{-2}
  % ordinal time axis with calendar ticks
  \hpbaxis{0}{17.7}{-2.75}
  \hpbtick{1}{Jan 2026}{-2.75}
  \hpbtick{10}{Jun}{-2.75}
  \hpbtick{14}{Nov}{-2.75}
\end{hpbgraph}}
\caption{Situation~F: the same six items on a Kanban board (a) and in
HPB (b) (illustrative). The board shows the workflow state of each item. HPB
adds what the board cannot show: the version each item started from, the
version that closed it, that T-7 was discarded rather than delivered, and
which items are committed to \texttt{v1.2} rather than merely open.}
\label{fig:same-data}
\end{figure}


\paragraph{Situation.}
Most teams already have a board, and they are right to. The first
question I am asked is: why do I need anything else?

\paragraph{In HPB.}
\Cref{fig:same-data} shows the same six items twice: on a three-column
Kanban board~\citep{anderson2010kanban} and as an HPB diagram. On the
board, X-31, I-12 and T-7 are all \emph{Done}. In HPB, X-31 and I-12 merge
into \texttt{v1.1} while T-7 was discarded; X-35 and I-14 are committed to
\texttt{v1.2}, and L-3 to nothing.

\paragraph{Decision.}
The board answers a team question, \emph{in which workflow state is each
item?}, and answers it well. HPB answers the product questions: what was
delivered and in which version, what is committed to the next gate, and
which hardware reference the work belongs to. Both views can be driven by
the same tracker items (\cref{subsec:method-coexistence}).

% =====================================================================
\subsection{G. Maintained versions and back-ports}
\label{sec:cond-backport}

% F11 -- Condition G: maintained release stream and back-port.
% Illustrative data. Requires hpb.sty v1.1. Canvas: document defaults.
\begin{figure}[t]
\centering
\hpbfit{%
\begin{hpbgraph}
  % ---- main stream of the line (y = 0) and its lanes
  \hpblane{1}{imp}{Improvement (I)}{0}{16.7}
  \hpblane{-1}{iss}{Issue (X)}{0}{16.7}
  \hpbmainline{0}{13}
  \hpbmainline*{13}{16.7}
  \hpborigin{0}{Gateway GW-2}
  \hpbrelease{1}{v1.4}
  \hpbrelease{6}{v2.0}
  \hpbrelease{12}{v2.0.1}
  % ---- maintenance stream 1.x (y = -2.2), forked at v1.4, own band
  \hpbseparator{-1.6}{-2.4}{16.7}
  \hpblane{-2.2}{main}{1.x stream\\(maintenance)}{0}{1.8}
  \hpblane{-3.2}{iss}{Issue (X)}{0}{16.7}
  \hpbstream[0.8]{1}{-2.2}{13}
  \hpbmainline*{13}{16.7}[-2.2]
  \hpbrelease{7}{v1.4.1}[-2.2]
  \hpbplanned{15}{v1.4.2}{Feb 2027}[-2.2]
  % ---- branches
  \hpbbranch[forkoff=0.6, label={I-70 Mesh networking}]{imp}{1}{1}{3}{5}{merge}{6}
  \hpbbranch[forkoff=0.3, mergeline=-2.2, label={X-298 Log rotation}]{iss}{1}{-3.2}{2}{4}{merge}{7}
  \hpbbranch[label={X-305 TLS parsing}]{iss}{6}{-1}{8}{11}{merge}{12}
  \hpbstep{iss}{9}{-1}{pass}
  \hpbbranch[line=-2.2, label={X-306 Port of X-305}]{iss}{7}{-3.2}{10}{13}{target}{15}
  % ---- propagation: X-306 carries the fix of X-305 to the 1.x stream (the
  %      connector reaches the start node from the upper left, clear of its label)
  \hpbpropagates{8,-1}{10,-3.2}{back-port}
  % ---- NOW and ordinal time axis
  \hpbnow{13}{1}{-3.2}
  \hpbaxis{0}{16.7}{-4.0}
  \hpbtick{1}{Mar 2026}{-4.0}
  \hpbtick{6}{Jul}{-4.0}
  \hpbtick{7}{Aug}{-4.0}
  \hpbtick{13}{Nov}{-4.0}
  \hpbtick{15}{Feb 2027}{-4.0}
\end{hpbgraph}}
\caption{Situation~G: a maintained stream and a back-port
(illustrative). The 1.x stream leaves \texttt{v1.4}; its patch
\texttt{v1.4.1} ships after \texttt{v2.0} but is not part of
\texttt{v2.0.1}. X-305 is fixed in \texttt{v2.0.1} and ported to 1.x as a
separate branch, X-306, linked by a propagation connector and targeted at
\texttt{v1.4.2}.}
\label{fig:backport}
\end{figure}


\paragraph{Situation.}
Some customers cannot move to the new major version, so the old one is
maintained in parallel. Every fix now raises a second question: does it
also have to go into the old version, and did anyone do it?

\paragraph{In HPB.}
In \cref{fig:backport}, a 1.x maintenance stream leaves \texttt{v1.4} in
its own band. Its patch \texttt{v1.4.1} ships after \texttt{v2.0}, but
precedence follows streams, not dates: no path connects them. X-305 is
fixed in \texttt{v2.0.1}; since 1.x is affected too, X-306 is opened on
the stream as a port, linked to X-305 by a propagation connector.

\paragraph{Decision.}
Which maintained versions contain a fix (follow the rises and the stream)
and which fixes have not been ported yet (merged issues known to affect
another stream, with no port and no recorded decision). The cost of
keeping 1.x alive is the work that forks from it.

% =====================================================================
\subsection{H. From field return to corrective action}
\label{sec:cond-rma}

% F12 -- Condition H: field return to corrective action.
% Illustrative data. Requires hpb.sty v1.1. Canvas: document defaults.
% No improvement lane in this diagram, so the lab lane is the first lane
% above the main line (lanes are stacked outward without gaps).
\begin{figure}[t]
\centering
\hpbfit{%
\begin{hpbgraph}
  \hpblane{1}{lab}{Lab (L)}{0}{17.7}
  \hpblane{-1}{iss}{Issue (X)}{0}{17.7}
  \hpblane{-2}{iss}{Issue (X)}{0}{17.7}
  \hpbmainline{0}{13}
  \hpbmainline*{13}{17.7}
  \hpborigin{0}{DMS-4000\\Rev B}
  \hpbrelease{1}{v2.0}
  \hpbprocess[]{8}{PC-7}
  \hpbplanned[interchangeable]{15}{Rev B.1}{Mar 2027}
  % field return closed without a design change: contained (rationale class
  % "resolved without a product baseline")
  \hpbbranch[forkoff=0.6, label={X-152 Housing damage}]{iss}{1}{-1}{3}{4}{discard}{}
  \hpblabel[right]{4}{-1}{\itshape contained:\\\itshape stock reworked}
  % field return with failure analysis, closed by a process baseline, then
  % checked for effectiveness after closure
  \hpbbranch[forkoff=0.3, label={X-150 Relay failure (RMA)}]{iss}{1}{-2}{2}{7}{merge}{8}
  \hpbstep{iss}{5}{-2}{done}
  \hpbstep{iss}{6}{-2}{pass}
  \hpbpoststep{iss}{7.5}{10}{-2}{pass}
  \hpblabel{10}{-2}{\itshape effectiveness verified}
  % supplier end-of-life notice: qualification of an alternative relay,
  % targeted at an interchangeable hardware baseline
  \hpbbranch[forkoff=0.6, label={L-9 PCN-221 relay EOL}]{lab}{8}{1}{9}{13}{target}{15}
  \hpbstep{lab}{11}{1}{pass}
  \hpbnow{13}{1}{-2}
  % ordinal time axis
  \hpbaxis{0}{17.7}{-3.0}
  \hpbtick{1}{Apr 2026}{-3.0}
  \hpbtick{8}{Jul}{-3.0}
  \hpbtick{13}{Nov}{-3.0}
  \hpbtick{15}{Mar 2027}{-3.0}
\end{hpbgraph}}
\caption{Situation~H: from field return to corrective action
(illustrative). X-150 is solved by changing how the product is built, so it
merges into the process baseline PC-7 (square), and a step after closure
records that the action worked. X-152 is closed as contained, without a
product baseline. The supplier notice PCN-221 opens L-9, targeted at the
interchangeable hardware baseline \texttt{Rev~B.1}.}
\label{fig:rma}
\end{figure}


\paragraph{Situation.}
Units come back from the field. Some failures are fixed by changing how
the product is built, not what it is; some are not product problems at
all. Quality wants to know what was done, on which evidence, and whether
it worked.

\paragraph{In HPB.}
In \cref{fig:rma}, X-150 \emph{Relay failure} carries the failure analysis
and the verification on a pilot lot as steps, and merges into the process
baseline PC-7 (a square). A dotted continuation records that the action
was effective. X-152 closes as \emph{resolved without a product baseline}
(transport damage, stock reworked). A supplier end-of-life notice opens
the lab branch L-9, targeted at an interchangeable hardware baseline.

\paragraph{Decision.}
The corrective-action loop is recorded without inventing product
versions. Containment is its own closure class and does not inflate the
abandoned-work ratio; effectiveness is a step after closure.

% =====================================================================
\subsection{I. Shared firmware and a joint launch}
\label{sec:cond-shared}

% F13 -- Condition I: firmware shared by two hardware references, with a
% dependency and a coordinated launch.  Illustrative data.  Requires hpb.sty
% v1.1.  Canvas: document defaults.  Three lines, each in its own band:
% hardware Rev B (top), the shared firmware line (middle), hardware Rev A.
\begin{figure}[t]
\centering
\hpbfit{%
\begin{hpbgraph}
  % ---- hardware line DMS-4000 Rev B (y = 2.7)
  \hpbmainline{0}{10}[2.7]
  \hpbmainline*{10}{17.7}[2.7]
  \hpborigin{0}{DMS-4000\\Rev B}[2.7]
  \hpbrelease{2}{B.1}[2.7]
  \hpbplanned{13}{B.2}{Apr 2027}[2.7]
  \hpbseparator{1.85}{-2.4}{17.7}
  % ---- shared firmware line (y = 0) with its lanes
  \hpblane{1}{imp}{Improvement (I)}{0}{17.7}
  \hpblane{-1}{iss}{Issue (X)}{0}{17.7}
  \hpbmainline{0}{10}
  \hpbmainline*{10}{17.7}
  \hpborigin{0}{DMS firmware}
  \hpbrelease{3}{v2.0}
  \begin{pgfonlayer}{hpbfront}% applies-to set, left of the tag (right side
    % is reserved for upward forks)
    \node[font=\sffamily\itshape\fontsize{6.5}{7.5}\selectfont, text=hpbMuted,
          anchor=east, inner sep=1pt] at ([xshift=-1.5pt]hpb@t.west)
      {applies to A, B};
  \end{pgfonlayer}
  \hpbplanned[applies to A, B]{14}{v2.1}{Apr 2027}
  \hpbseparator{-1.85}{-2.4}{17.7}
  % ---- hardware line DMS-4000 Rev A (y = -2.7)
  \hpbmainline{0}{10}[-2.7]
  \hpbmainline*{10}{17.7}[-2.7]
  \hpborigin{0}{DMS-4000\\Rev A}[-2.7]
  \hpbrelease{1}{A.3}[-2.7]
  \hpbmilestone{7}{end of sale}[-2.7]
  % ---- branches of the firmware line
  \hpbbranch[forkoff=0.5, label={X-140 Relay chatter (Rev A + v2.0)}]{iss}{3}{-1}{4}{10}{target}{14}
  \hpbstep{iss}{8}{-1}{pass}
  \hpbbranch[forkoff=0.6, label={I-77 Low-power mode}]{imp}{3}{1}{5}{10}{target}{14}
  \hpbstep{imp}{6}{1}{done}
  % ---- I-77 depends on the hardware baseline B.2; launch bundle
  \hpbdepends{10,1}{13,2.7}{depends on}
  \hpbbundle{12.1,-0.55}{17.4,3.85}{launch bundle K-1 (all or nothing)}
  \hpbnow{10}{2.7}{-2.7}
  % ---- ordinal time axis
  \hpbaxis{0}{17.7}{-3.75}
  \hpbtick{1}{Feb 2026}{-3.75}
  \hpbtick{3}{May}{-3.75}
  \hpbtick{7}{Sep}{-3.75}
  \hpbtick{10}{Nov}{-3.75}
  \hpbtick{13}{Apr 2027}{-3.75}
\end{hpbgraph}}
\caption{Situation~I: one firmware shared by two hardware references
(illustrative). Each firmware version declares the hardware it applies to;
X-140 is drawn on the firmware line, with its observed configuration in the
label. I-77 depends on the hardware baseline B.2, so \texttt{v2.1} and B.2
form the launch bundle K-1: published together or not at all
(\cref{rule:gate}).}
\label{fig:shared}
\end{figure}


\paragraph{Situation.}
One firmware build runs on several hardware references, and a new feature
needs both a firmware release and a hardware change. A report says
``relay chatter'', but the real question is: on which combinations?

\paragraph{In HPB.}
In \cref{fig:shared}, the firmware is a line of its own between the
hardware lines Rev~A and Rev~B, and each firmware version declares the
hardware it applies to. X-140 is drawn on the firmware line with its
observed configuration (Rev~A, \texttt{v2.0}). I-77 depends on the
hardware baseline B.2, so \texttt{v2.1} and B.2 form the launch bundle
K-1.

\paragraph{Decision.}
Who is affected (the applies-to sets and the observed configuration) and
whether the launch can go ahead: by \cref{rule:gate}, K-1 ships only when
everything targeted at either baseline is verified. A slip of B.2 visibly
delays \texttt{v2.1}, unless I-77 is moved out.

% ---------------------------------------------------------------------
%  07-operating.tex -- Running HPB: the operating guide
%  Practices, intake and typing, cadence, roles, coexistence with
%  existing frameworks, adoption path, anti-patterns.
% ---------------------------------------------------------------------
\section{Running HPB: The Operating Guide}
\label{sec:operating}

The model says what a correct HPB graph is, and the notation says how to
draw it. Neither says who changes the graph, when, and on what evidence.
This section does. The operating guide is deliberately small: a handful
of practices, one intake convention, four recurring meetings, roles that
already exist, and an adoption path. It is meant to be started by one
product manager on one product line and to scale by repetition, not by
adding ceremony. It prescribes \emph{what must be recorded} at each
decision; how engineering organizes its own work stays with whatever
team process is already in place (\cref{subsec:method-coexistence}).

\subsection{Four principles}
\label{subsec:method-principles}

\begin{enumerate}[label=\textbf{M\arabic*},leftmargin=*,itemsep=2pt]
  \item \textbf{Record at the decision, not after it.} Every change to the
    graph matches a decision that is already being taken: accepting a
    request, committing it to a release, publishing a version, dropping
    work. HPB attaches the record to that moment, so the graph is a
    by-product of deciding, not an extra reporting duty.
  \item \textbf{One accountable owner per node.} Every line, baseline and
    branch has exactly one accountable role. Many people contribute; one
    answers for it.
  \item \textbf{Closure is explicit and evidenced.} Nothing disappears.
    A branch either merges into a named baseline or is discarded with a
    written reason (\cref{rule:closure}). This is the principle that
    produces traceability, and the one list-based trackers break most
    easily, because there an item is simply set to ``closed''.
  \item \textbf{History is fixed, the plan is not.} Everything left of NOW
    is a record and is corrected only through audited events. Everything
    right of NOW is a plan and is expected to change. Plan changes should
    be cheap; history changes should be visible.
\end{enumerate}

\subsection{Practices}
\label{subsec:method-practices}

\Cref{tab:practices} lists the practices. Each one has a trigger, the
roles involved, the rule it enforces and what it leaves in the graph.
The point of the table is that the rules are not checked in a separate
compliance step: each rule is enforced at the one moment in the workflow
where it can be broken.

\begin{table}[tbp]
  \centering
  \small
  \caption{HPB practices. Each practice runs at a clear trigger and enforces
    at least one rule (R1--R7). Roles: A accountable, R responsible (if
    different), C consulted or approving. PM product manager; RM release or
    program manager; CCB change control board; Eng R\&D/engineering; QA
    quality assurance or laboratory; QR quality or regulatory function.}
  \label{tab:practices}
  \begin{tabularx}{\linewidth}{@{}>{\raggedright\arraybackslash}p{2.0cm}
                                  >{\raggedright\arraybackslash}X
                                  >{\raggedright\arraybackslash}p{2.2cm}
                                  >{\raggedright\arraybackslash}p{1.0cm}
                                  >{\raggedright\arraybackslash}X@{}}
    \toprule
    \textbf{Practice} & \textbf{Trigger} & \textbf{Roles} & \textbf{Rules} & \textbf{Leaves in the graph} \\
    \midrule
    Fork at intake
      & A request, defect report, task, lab evaluation or supplier notice is accepted at triage.
      & A: PM; R: support, Eng, QA
      & R1, R6
      & Typed branch with origin baseline (and configuration), opening date, owner. \\
    \addlinespace
    Target at planning
      & Release planning for a planned baseline; a priority change.
      & A: PM; C: RM
      & R4, R7
      & Target baseline; dashed continuation; earlier targets kept as history. \\
    \addlinespace
    Verify
      & A candidate of the target (build, prototype, pilot lot) is tested.
      & A: QA; R: Eng
      & R4, R7
      & Pass/fail step citing the candidate; a pass makes the branch verified. \\
    \addlinespace
    Gate review
      & Freeze or target date of a planned baseline or launch bundle.
      & A: RM; C: PM, Eng, QA
      & R4
      & Decision on the baseline: release, slip with new date, or move/discard the unverified members. \\
    \addlinespace
    Merge at release
      & The baseline is published (build released, or change order approved with production effectivity).
      & A: RM; R: Eng
      & R2, R3
      & Merge of every verified member; release notes derived. \\
    \addlinespace
    Discard with reason
      & Not viable, duplicate, not reproducible, superseded, out of scope; or resolved without a product baseline.
      & A: PM; R: owner; C: QR if safety or regulatory
      & R2, R7
      & Discard terminal with reason class, category and text. \\
    \addlinespace
    Reopen (audited)
      & A closure was recorded in error.
      & A: PM; C: second approver
      & R2, R7
      & Reopen event with reason and actors; original terminal kept. \\
    \addlinespace
    Line fork
      & A hardware change fails the interchangeability test.
      & A: CCB; R: PM, hardware Eng
      & R5
      & Diamond baseline, child line with parent link, change-order reference, parent life-cycle milestones. \\
    \addlinespace
    Record step
      & A sub-task completes, a test runs, an effectiveness check is made.
      & A: Eng or QA
      & R6, R7
      & Step (done/pending, pass/fail) with evidence attached. \\
    \bottomrule
  \end{tabularx}
\end{table}

\paragraph{Fork at intake.}
A request becomes a branch only when it is accepted at triage. At that
moment three facts are fixed: its type, its origin baseline and its
opening date (\cref{rule:origin}). For an issue, the origin is the
version in which the defect was \emph{seen}: the firmware version read
from the device, the software build, the hardware revision on the label.
It is mandatory. If nobody knows it, clarify before accepting, or accept
it with the origin flagged as \emph{inferred}; when the defect depends on
a combination of hardware and firmware, record the configuration too.
Triage is also where a request can be rejected before it ever becomes a
branch, which keeps the graph free of noise.

\paragraph{Granularity and typing.}
Diagrams are only comparable between product managers if the same kinds
of request become branches with the same types. \Cref{tab:intake} is the
default convention. A branch is a \emph{release-note-worthy unit of
change}: an epic or feature, a customer-visible defect, a lab campaign.
User stories, sub-tickets and commits stay on team boards; they may
appear on a branch as steps (\cref{rule:typed}). Security, regulatory
relevance, safety or a specific customer are \emph{attributes} of a
branch, not extra types, so the four lanes stay stable.

\begin{table}[tbp]
  \centering
  \small
  \caption{Default intake and typing convention. A request becomes one
    branch of the listed type; attributes are set in addition.
    Organizations may refine the table, but should publish it with their
    diagrams.}
  \label{tab:intake}
  \begin{tabularx}{\linewidth}{@{}>{\raggedright\arraybackslash}X
                                  >{\centering\arraybackslash}p{1.0cm}
                                  >{\raggedright\arraybackslash}p{5.6cm}@{}}
    \toprule
    \textbf{Request} & \textbf{Type} & \textbf{Origin baseline; attributes} \\
    \midrule
    New or changed customer-visible capability (epic, feature) & $\mathsf{I}$ & latest release; customer if requested by one \\
    Defect seen in a released version & $\mathsf{X}$ & version where it was seen (mandatory); configuration if relevant \\
    Security vulnerability & $\mathsf{X}$ & version where it was seen; attribute \emph{security} \\
    Regulatory or certification change & $\mathsf{T}$ & latest release; attribute \emph{regulatory}; $\mathsf{I}$ if it adds a capability \\
    Customer-specific request & $\mathsf{I}$ or $\mathsf{X}$ & as above; attribute \emph{customer}, with the commitment made \\
    Cost-down redesign, component substitution & $\mathsf{L}$ & latest release; line fork only if not interchangeable \\
    Supplier change or end-of-life notice & $\mathsf{L}$ or $\mathsf{T}$ & latest release; $\mathsf{L}$ to qualify an alternative, $\mathsf{T}$ for a last-time buy \\
    Field return with failure analysis & $\mathsf{X}$ & version of the returned unit; RMA evidence as steps \\
    Internal work (certificates, tooling, documentation) & $\mathsf{T}$ & latest release \\
    User story, sub-ticket, single test run & -- & a step of its branch, not a branch \\
    \bottomrule
  \end{tabularx}
\end{table}

\paragraph{Target at planning.}
Targeting commits an open branch to a planned version. It is the only
practice that is expected to repeat for the same branch: when priorities
change, the branch is re-targeted and the old target stays visible as a
faded dashed line (\cref{fig:discard}). My rule of thumb: target a branch
only when its owner can state what ``done'' means for the merge. A target
without an acceptance condition is a wish, and wishes inflate the gate.
Before committing the targeted set, compare its size with what recent
releases of the same kind actually contained (\cref{tab:metrics}), and check
that every dependency of a member is targeted at the same or an earlier
version.

\paragraph{Gate review.}
At the freeze date (or the target date, for teams without a freeze) the
release manager, with product, engineering and quality, goes through the
targeted set. Each member is verified, not yet verified, or no longer
wanted; each dependency is satisfied or not. \Cref{rule:gate} forbids
publishing while any member is unverified or depends on unreleased work,
so the review must end in one of three recorded outcomes:
\emph{release}, if everything is verified; \emph{slip}, with a new date
and the reason written on the baseline; or \emph{descope}, moving or
discarding the unverified members. The rule does not choose for you. It
forces the choice to be made and written down. \Cref{fig:release-gate}
shows the situation the review resolves.

\paragraph{Discard with reason.}
Discarding is a normal outcome, not a failure. The engineering-change
literature treats the rejection of proposed changes as part of the change
process \citep{jarratt2011ec}; HPB makes those decisions visible instead of
hiding them in a generic ``closed''. The reason is picked from the two
closed lists of \cref{sec:model-closing}, which makes discards countable,
plus free text and, for a duplicate, a link to the surviving branch, which
makes them auditable. Discarding a branch marked \emph{safety} or
\emph{regulatory} also needs the approval of the quality or regulatory
function.

\paragraph{The other practices.}
\emph{Verify}: whoever ran the test records the passed step with the
candidate it ran on (a release-candidate build, a prototype or a pilot
lot); that record replaces the familiar ``yes, it's in the release''.
\emph{Merge at release}: publishing a version merges all its verified
targeted branches at once, and a hardware baseline cites its change
order. \emph{Reopen} needs a second approver and keeps the original
terminal; a fix that regresses later is a \emph{new} issue from the
version where the regression is seen. \emph{Line fork}: when the change
control board approves a new reference, the parent gets dated life-cycle
milestones at the same time; leaving them open is the surest way to
accumulate maintenance nobody planned for. \emph{Record step}: sub-tasks,
test runs and effectiveness checks are steps, so a lab branch becomes an
evidence trail (\cref{fig:hardware-fork,fig:rma}).

\subsection{Cadence}
\label{subsec:method-cadence}

The practices are event-driven, but most events land in four recurring
meetings. Adapt the frequencies to the release rhythm of your product;
keep the order.

\begin{description}[leftmargin=0pt,itemsep=3pt,font=\normalfont\bfseries]
  \item[Weekly triage (30--45 minutes).] New requests from customers,
    support, sales, engineering and the lab. Practices: fork at intake;
    discard obvious duplicates. Read the diagram below the line first (new
    issues and their origins often reveal a bad release early), then above
    it. End by reviewing open, untargeted branches older than an agreed
    threshold (age of open work, \cref{tab:metrics}).
  \item[Release planning (per planned version).] Open branches, capacity
    from the team process, dependencies, customer commitments. Practice:
    target at planning, with the size and dependency checks. Output: the
    targeted set of the next version or launch bundle and, optionally,
    provisional targets for the one after. Planning only targets work that
    already exists as a branch, so every planned change has a recorded
    origin.
  \item[Gate review (at freeze or target date).] Gate review, then merge
    at release if the decision is to ship. The diagram keeps it short: the
    dashed lines entering the gate without a passed step are, by
    construction, exactly the open questions.
  \item[Line review (monthly or after each release).] Release
    contents, escaped defects, installed base and rollout where available,
    unported fixes, and the metrics of \cref{sec:metrics}. Decisions:
    life-cycle milestones of maintained lines and streams, whether a patch
    is needed, which fixes must be ported, and whether recent discard
    reasons point to a problem in intake quality.
\end{description}

\subsection{Roles}
\label{subsec:method-roles}

HPB introduces no new job titles. It gives accountability for graph
elements to roles that product organizations already have
(\cref{tab:practices}).

\begin{itemize}[leftmargin=*,itemsep=2pt]
  \item \textbf{Product manager.} Owns the product lines and the content of
    planned versions: triage, typing, targeting, discards. Can delegate,
    as long as the actor of every transition is recorded.
  \item \textbf{Release or program manager.} Owns dates and gates: runs the
    gate review and publishes versions and launch bundles.
  \item \textbf{Change control board.} Where one exists, approves change
    orders, and therefore hardware baselines and line forks; the HPB
    baseline references the approved order.
  \item \textbf{R\&D / engineering.} Owns improvement and task branches
    once accepted, and records their steps.
  \item \textbf{Quality / laboratory.} Owns lab branches, verification
    steps and their evidence; is consulted at every gate. The quality or
    regulatory function approves discards of safety- or
    regulation-relevant issues and the closure of corrective actions.
  \item \textbf{Support and field service.} The main source of issues.
    Owns the quality of the reported version and configuration, which is
    what escaped-defect counts depend on, and the installed-base data
    where it is kept.
\end{itemize}

\noindent In a small organization one person wears several hats. That is
fine, as long as the role under which each action is taken is visible in
the audit trail. Two segregations of duty stay even then: reopens and
safety- or regulation-relevant discards need a second approver.

\subsection{Living with Scrum, Kanban, Stage-Gate and PLM}
\label{subsec:method-coexistence}

HPB is a \emph{product-level} view. It does not schedule engineers, set
team WIP limits or replace the boards teams use every day. It sits above
them and answers what they are not designed to answer: from which product
state work started, which version closed it, and what else was in flight
on the same product at the same time. \Cref{tab:coexistence} gives the
mapping; \cref{fig:same-data} shows the same items in both views.

\begin{table}[tbp]
  \centering
  \small
  \caption{How HPB maps onto the units of common frameworks. HPB adds the
    product-state anchors (origin and closing baseline) without replacing
    the framework's own planning unit.}
  \label{tab:coexistence}
  \begin{tabularx}{\linewidth}{@{}>{\raggedright\arraybackslash}p{2.0cm}
                                  >{\raggedright\arraybackslash}X
                                  >{\raggedright\arraybackslash}X@{}}
    \toprule
    \textbf{Framework} & \textbf{Its unit} & \textbf{In HPB} \\
    \midrule
    Scrum
      & Product backlog item; sprint; increment \citep{schwaber2020scrum}.
      & A backlog item is (part of) a branch. A sprint is a team time-box,
        \emph{not} a baseline; an increment becomes a baseline only when it
        is released to customers. \\
    \addlinespace
    Kanban
      & Card; column; WIP limit \citep{anderson2010kanban}.
      & A card is a branch or a step of one. Columns refine the HPB states
        open and targeted; ``done'' becomes a merge only once the closing
        version is published. \\
    \addlinespace
    Stage-Gate
      & Project; stage; go/kill gate \citep{cooper1990stagegate}.
      & A new-reference project is lab and improvement work ending in a
        hardware baseline and a line fork. A \emph{kill} is a discard with
        reason. Stage-Gate gates decide whether to keep investing; HPB
        gates decide whether a release is complete. \\
    \addlinespace
    PLM / CM
      & Change order; configuration baseline \citep{stark2015plm,iso10007}.
      & A released HPB baseline references the change order or
        configuration baseline that formally defines it. \\
    \bottomrule
  \end{tabularx}
\end{table}

The most common confusion is the sprint: it is a team time-box, not a
baseline, because firmware and hardware ship less often than teams finish
sprints. In practice the product owner usually holds the HPB product
manager role, and board and diagram share the same tracker items, so
there is no double entry.

\paragraph{Regulated contexts.}
HPB does not replace configuration management as defined in IEEE~828 and
ISO~10007 \citep{ieee828,iso10007}, nor change processes in a PLM system
\citep{stark2015plm}; it adds the planning view and the link between the
intent of a change (the branch) and the configuration that realized it
(the baseline). Used as evidence, it needs at least append-only history,
identified actors and approvers on every closing decision, links to the
requirement, risk and test records, and a per-version export. Whether
that satisfies a particular regulation is outside the scope of this paper.

\subsection{Start on Monday: an adoption path}
\label{subsec:method-adoption}

HPB needs no migration and no new tool on day one. This is the sequence I
recommend.

\paragraph{Week 1: one line, drawn by hand.}
Pick \emph{one} product line with parallel work and a release coming up,
not the whole portfolio. Write down its last three to five released
versions and where the line came from. Go through the open items in your
trackers and record, for each release-note-worthy one, its type and its
best-known origin, marking guesses as \emph{inferred}. Create the next
planned version and target what is already promised for it: the gate is
now visible. Draw it, on a whiteboard, with \texttt{hpb.sty}
(\cref{app:package}) or in a few lines of HPB Markup (\cref{sec:tooling}),
and show it at your next triage.

\paragraph{Month 1: run the cadence.}
Apply \emph{fork at intake} to every new request: no branch without a type
and an origin. Run weekly triage on the diagram, issues first. Run the
first gate review at the next freeze and record its outcome. When the
version ships, merge its verified branches; from then on, release notes
are generated, not written. Publish your intake convention
(\cref{tab:intake}) with your changes.

\paragraph{Quarter 1: extend and measure.}
After a full release cycle, add the next lines, with their hardware
revisions, line forks and maintained streams. Set life-cycle milestones
for every old reference still maintained. Hold the first line review with
the metrics of \cref{sec:metrics}, and connect HPB to your tracker so the
diagram is generated from the items, not maintained by hand.

\subsection{Anti-patterns}
\label{subsec:method-antipatterns}

These are easy to spot in the diagram, and each one tells you which
practice is being skipped.

\begin{itemize}[leftmargin=*,itemsep=1pt]
  \item \emph{Silent closure.} A tracker item is closed but its branch has
    no merge or discard terminal (breaks \cref{rule:closure}).
  \item \emph{Eternal branch.} An open, untargeted branch that keeps
    outliving releases: usually a discard nobody dared to decide.
  \item \emph{Gate stuffing.} A targeted set much larger than recent
    releases, which turns every gate review into a slip.
  \item \emph{Floating evidence.} Test reports and decisions kept outside
    the graph (breaks \cref{rule:discussion}).
  \item \emph{In-place reference edit.} A non-interchangeable hardware
    change without a new line (breaks \cref{rule:lineage}).
  \item \emph{Unported fix.} A merged issue known to affect another
    maintained stream, with no port and no decision not to port.
  \item \emph{Shipping ahead of a dependency.} A dependency connector that
    points to the right of the gate (blocked by \cref{rule:gate}).
\end{itemize}

% ---------------------------------------------------------------------
%  08-metrics.tex -- Metrics that fall out of the graph
%  The formulas (labels eq:leadtime ... eq:unprop) live in
%  appendix-a-formal.tex, subsection app:metrics. % F10 -- Metrics derived from the graph, computed from the Condition A example
% (figure "software-line"). Illustrative data. Requires hpb.sty v1.0 and
% pgfplots (compat=1.18). Ordinal columns are those of fig:software-line.
\begin{figure}[t]
\centering
\resizebox{\linewidth}{!}{%
\begin{tikzpicture}[
    font=\sffamily\fontsize{7}{8.4}\selectfont,
    mx/bar/.style  = {draw=hpbInk, line width=0.8pt, rounded corners=2pt},
    mx/in/.style   = {font=\sffamily\fontsize{6.5}{7.8}\selectfont, text=hpbInk,
                      inner sep=1pt},
    mx/ann/.style  = {hpb/callout, align=left},
  ]
  \pgfplotsset{
    mx/common/.style = {
      width=13.2cm, scale only axis, xmin=0, xmax=17,
      axis line style={draw=hpbDiscInk, line width=0.6pt},
      tick style={draw=hpbDiscInk, line width=0.6pt},
      tick label style={font=\sffamily\fontsize{7}{8.4}\selectfont, text=hpbInk},
      label style={font=\sffamily\fontsize{7}{8.4}\selectfont, text=hpbInk},
      xtick={1,5,10,16},
      xmajorgrids, major x grid style={draw=hpbDiscInk, line width=0.5pt,
        dash pattern=on 2pt off 2pt},
      clip=false, axis on top=false,
    },
  }
  % ------------------------------------------------------------------ (a)
  \begin{axis}[mx/common, name=wip, height=3.3cm,
      ymin=0, ymax=4.3, ytick={0,1,2,3}, xticklabels={},
      ymajorgrids, major y grid style={draw=hpbGuide, line width=0.4pt},
      ylabel={open branches $|W_\theta(t)|$},
      stack plots=y, const plot mark left,
      every axis plot/.append style={draw=hpbInk, line width=0.8pt, mark=none},
      area legend,
      reverse legend,
      legend style={draw=none, fill=white, fill opacity=0.85, text opacity=1,
        at={(0.008,0.975)}, anchor=north west, inner sep=1pt,
        legend cell align=left, row sep=-1pt,
        font=\sffamily\fontsize{7}{8.4}\selectfont, text=hpbInk},
    ]
    % stacked bottom-to-top as in the diagram: T (lane -2), X (lane -1), I (lane +1)
    \addplot[fill=hpbTask] coordinates {
      (0,0)(1,0)(2,0)(3,0)(4,0)(5,0)(6,0)(7,1)(8,1)(9,1)(10,1)(11,1)(12,1)
      (13,1)(14,1)(15,1)(16,0)(17,0)} \closedcycle;
    \addplot[fill=hpbIss] coordinates {
      (0,0)(1,0)(2,1)(3,1)(4,1)(5,0)(6,1)(7,1)(8,1)(9,1)(10,0)(11,0)(12,1)
      (13,1)(14,1)(15,1)(16,0)(17,0)} \closedcycle;
    \addplot[fill=hpbImp] coordinates {
      (0,0)(1,0)(2,0)(3,1)(4,1)(5,1)(6,1)(7,1)(8,1)(9,1)(10,0)(11,1)(12,1)
      (13,1)(14,1)(15,1)(16,0)(17,0)} \closedcycle;
    \legend{Task (T), Issue (X), Improvement (I)}
    % type letters inside the bands: the stack is readable without color
    \foreach \c/\y/\l in {12.5/0.5/T, 3/0.5/X, 12.5/1.5/X,
                          4/1.5/I, 12.5/2.5/I}{
      \edef\mxtmp{\noexpand\node[font=\noexpand\sffamily\noexpand\bfseries
        \noexpand\fontsize{7}{8.4}\noexpand\selectfont, text=hpbInk,
        inner sep=0pt] at (axis cs:\c,\y) {\l};}\mxtmp
    }
    % baseline tags above the panel
    \foreach \c/\v in {1/v1.0, 5/v1.0.1, 10/v1.1, 16/v1.2}{
      \edef\mxtmp{\noexpand\node[hpb/tag, anchor=south, yshift=3pt]
        at (axis cs:\c,4.3) {\v};}\mxtmp
    }
    % annotations
    \draw[hpb/leader] (axis cs:9.7,2.2) |- (axis cs:8.6,3.72)
      node[mx/ann, anchor=east]
      {release of v1.1: I-21 and X-102\\close together ($|C|=2$)};
    \draw[hpb/leader] (axis cs:15.7,2.2) |- (axis cs:15.1,3.72)
      node[mx/ann, anchor=east]
      {batch merge at v1.2:\\three branches close ($|C|=3$)};
  \end{axis}
  % ------------------------------------------------------------------ (b)
  \begin{axis}[mx/common, at={(wip.south west)}, anchor=north west,
      yshift=-0.55cm, height=3.6cm,
      ymin=0.4, ymax=6.6, ytick={1,...,6}, y dir=reverse,
      yticklabels={I-21, I-22, X-101, X-102, X-103, T-8},
      y tick label style={font=\sffamily\bfseries\fontsize{7}{8.4}\selectfont},
      major y tick style={draw=none},
      xticklabels={Jan, Feb, Apr, Jul},
      xlabel={ordinal time (one column per event; calendar month of each baseline)},
      ylabel={branch},
    ]
    % one bar per branch: from opening t_o(b) (start node) to closing t_c(b)
    % (release of the merge baseline);  #1 row, #2 type color, #3 letter,
    % #4 t_o column, #5 t_c column, #6 fork -> merge baseline, #7 D(b)
    \foreach \r/\tc/\tl/\o/\e/\fm/\dd in {
        1/hpbImp/I/3/10/{v1.0 $\to$ v1.1}/2,
        2/hpbImp/I/11/16/{v1.1 $\to$ v1.2}/1,
        3/hpbIss/X/2/5/{v1.0 $\to$ v1.0.1}/1,
        4/hpbIss/X/6/10/{v1.0.1 $\to$ v1.1}/1,
        5/hpbIss/X/12/16/{v1.1 $\to$ v1.2}/1,
        6/hpbTask/T/7/16/{v1.0.1 $\to$ v1.2}/2}{
      \edef\mxtmp{%
        \noexpand\draw[mx/bar, fill=\tc!55]
          (axis cs:\o,\r-0.27) rectangle (axis cs:\e,\r+0.27);
        \noexpand\node[hpb/event, fill=\tc] at (axis cs:\o,\r) {\tl};
        \noexpand\node[mx/in, anchor=west] at (axis cs:\o,\r) [xshift=9pt]
          {\unexpanded\expandafter{\fm}};
        \noexpand\node[mx/in, font=\noexpand\sffamily\noexpand\bfseries
          \noexpand\fontsize{7}{8.4}\noexpand\selectfont]
          at (axis cs:17.55,\r) {\dd};}%
      \mxtmp
    }
    \node[mx/in, font=\sffamily\bfseries\fontsize{7}{8.4}\selectfont]
      at (axis cs:17.55,0.4) [anchor=south, yshift=1pt] {$D(b)$};
  \end{axis}
  % panel letters
  \node[font=\sffamily\bfseries\fontsize{7.5}{9}\selectfont, anchor=north east]
    at ([xshift=-1.15cm]wip.north west) {(a)};
  \node[font=\sffamily\bfseries\fontsize{7.5}{9}\selectfont, anchor=north east]
    at ([xshift=-1.15cm,yshift=-0.55cm]wip.south west) {(b)};
\end{tikzpicture}}
\caption{Flow metrics read straight from the graph, for the data of
situation~A (\cref{fig:software-line}) on the same ordinal columns. (a)~Work
in progress per type: each opening raises a step, each closure lowers it, and
a release drops it by its batch size $|C(v)|$. (b)~One bar per branch, from
its opening to the release that closed it, with its fork and merge versions
and its merge distance $D(b)$. Lead times in calendar time come from the
event dates, not from the length of the bars.}
\label{fig:metrics}
\end{figure}

%  appears here exactly once.
% ---------------------------------------------------------------------
\section{Metrics That Fall Out of the Graph}
\label{sec:metrics}

A useful side effect of the rules is that measurement comes for free.
Every branch has a type, an origin baseline and an opening date and, once
closed, exactly one terminal with a closing date. That is enough to compute
the flow and quality metrics of \cref{tab:metrics} straight from the
graph, with no extra data collection and no spreadsheet kept on the side.
The exact definitions are in \cref{app:metrics}.\footnote{Two details keep
the counts honest. Sub-tasks are steps, not branches (\cref{rule:typed}),
so every count counts a unit of change once. And a reopened branch is rare
but possible: each period during which a branch is open or targeted is a
\emph{visit}, lead times are measured per visit, and reopens are reported
separately.}

All of them are \emph{descriptive}. They exist to bring good questions to
the meetings of \cref{subsec:method-cadence}, not to rank teams or people
(\cref{subsec:metrics-goodhart}).

\begin{table}[tbp]
  \centering
  \small
  \caption{The HPB metrics: the question each one answers, how to read it,
    and the meeting it feeds (\cref{subsec:method-cadence}). Definitions in
    \cref{app:metrics}.}
  \label{tab:metrics}
  \renewcommand{\arraystretch}{1.1}
  \begin{tabularx}{\linewidth}{@{}>{\raggedright\arraybackslash}p{2.3cm}
      >{\raggedright\arraybackslash}X
      >{\raggedright\arraybackslash}X
      >{\raggedright\arraybackslash}p{1.9cm}@{}}
    \toprule
    \textbf{Metric} & \textbf{Question it answers} & \textbf{How to read it}
      & \textbf{Feeds} \\
    \midrule
    Lead time
      & How long from accepting a request to having it in a published
        version (or deciding not to)?
      & Merged and discarded apart, by type. Ends at publication: the
        time customers feel.
      & Line review \\
    Age of open work
      & What has been open too long?
      & Older than the usual lead time of its type: on the triage agenda.
      & Triage \\
    Work in progress
      & How much is open now, per type?
      & Open issues rising and not brought down by releases: an early sign
        of a quality problem.
      & Triage, planning \\
    Throughput
      & How much do we deliver per period?
      & Use windows that span several releases; read it with lead time.
      & Line review \\
    Merge and discard ratios
      & Of what we closed, how much did we deliver?
      & Only by discard reason. Many duplicates: intake problem. Many
        \emph{not viable} lab branches: often healthy.
      & Line review \\
    Merge distance
      & How many releases did a change wait?
      & Large values: work pushed back again and again; its specification
        may be stale.
      & Planning \\
    Release batch size
      & How much does a version contain?
      & Compare a planned set with recent releases before committing it.
      & Planning \\
    Escaped defects
      & Which issues were seen in each version, per change shipped?
      & Compare only at similar exposure; ``seen in'' is not ``introduced
        by''.
      & Line review \\
    Gate slip
      & How late did the version ship?
      & Never alone: with the branches moved out and discarded from it.
      & Gate review \\
    Maintenance load
      & How much new work still lands on an old reference?
      & Sets life-cycle dates with numbers instead of opinions.
      & Line review \\
    Unported fixes
      & Which fixes affect another maintained stream and were not ported?
      & Should be zero after every line review.
      & Line review \\
    \bottomrule
  \end{tabularx}
\end{table}

\paragraph{Flow.}
Lead time ends when the version is \emph{published}, not when engineering
is done: it measures how long until the change is available in the
product, the same idea as lead time in software delivery
\citep{forsgren2018accelerate}, applied to firmware and hardware too. Its
leading counterpart is the age of what is still open. Work in progress is
a staircase over time, one step up at each fork, one down at each terminal
and a big drop at each release; \cref{fig:metrics} shows it for
situation~A. Little's law \citep{little1961} ties work in progress,
throughput and lead time together. I use it as a sanity check (if the
three numbers do not roughly agree, the window is too short or the system
is not stable) and as a planning lever, in the spirit of
\citet{reinertsen2009flow}: at a given delivery rate, the only way to
shorten lead time is to have less open work.

% F10 -- Metrics derived from the graph, computed from the Condition A example
% (figure "software-line"). Illustrative data. Requires hpb.sty v1.0 and
% pgfplots (compat=1.18). Ordinal columns are those of fig:software-line.
\begin{figure}[t]
\centering
\resizebox{\linewidth}{!}{%
\begin{tikzpicture}[
    font=\sffamily\fontsize{7}{8.4}\selectfont,
    mx/bar/.style  = {draw=hpbInk, line width=0.8pt, rounded corners=2pt},
    mx/in/.style   = {font=\sffamily\fontsize{6.5}{7.8}\selectfont, text=hpbInk,
                      inner sep=1pt},
    mx/ann/.style  = {hpb/callout, align=left},
  ]
  \pgfplotsset{
    mx/common/.style = {
      width=13.2cm, scale only axis, xmin=0, xmax=17,
      axis line style={draw=hpbDiscInk, line width=0.6pt},
      tick style={draw=hpbDiscInk, line width=0.6pt},
      tick label style={font=\sffamily\fontsize{7}{8.4}\selectfont, text=hpbInk},
      label style={font=\sffamily\fontsize{7}{8.4}\selectfont, text=hpbInk},
      xtick={1,5,10,16},
      xmajorgrids, major x grid style={draw=hpbDiscInk, line width=0.5pt,
        dash pattern=on 2pt off 2pt},
      clip=false, axis on top=false,
    },
  }
  % ------------------------------------------------------------------ (a)
  \begin{axis}[mx/common, name=wip, height=3.3cm,
      ymin=0, ymax=4.3, ytick={0,1,2,3}, xticklabels={},
      ymajorgrids, major y grid style={draw=hpbGuide, line width=0.4pt},
      ylabel={open branches $|W_\theta(t)|$},
      stack plots=y, const plot mark left,
      every axis plot/.append style={draw=hpbInk, line width=0.8pt, mark=none},
      area legend,
      reverse legend,
      legend style={draw=none, fill=white, fill opacity=0.85, text opacity=1,
        at={(0.008,0.975)}, anchor=north west, inner sep=1pt,
        legend cell align=left, row sep=-1pt,
        font=\sffamily\fontsize{7}{8.4}\selectfont, text=hpbInk},
    ]
    % stacked bottom-to-top as in the diagram: T (lane -2), X (lane -1), I (lane +1)
    \addplot[fill=hpbTask] coordinates {
      (0,0)(1,0)(2,0)(3,0)(4,0)(5,0)(6,0)(7,1)(8,1)(9,1)(10,1)(11,1)(12,1)
      (13,1)(14,1)(15,1)(16,0)(17,0)} \closedcycle;
    \addplot[fill=hpbIss] coordinates {
      (0,0)(1,0)(2,1)(3,1)(4,1)(5,0)(6,1)(7,1)(8,1)(9,1)(10,0)(11,0)(12,1)
      (13,1)(14,1)(15,1)(16,0)(17,0)} \closedcycle;
    \addplot[fill=hpbImp] coordinates {
      (0,0)(1,0)(2,0)(3,1)(4,1)(5,1)(6,1)(7,1)(8,1)(9,1)(10,0)(11,1)(12,1)
      (13,1)(14,1)(15,1)(16,0)(17,0)} \closedcycle;
    \legend{Task (T), Issue (X), Improvement (I)}
    % type letters inside the bands: the stack is readable without color
    \foreach \c/\y/\l in {12.5/0.5/T, 3/0.5/X, 12.5/1.5/X,
                          4/1.5/I, 12.5/2.5/I}{
      \edef\mxtmp{\noexpand\node[font=\noexpand\sffamily\noexpand\bfseries
        \noexpand\fontsize{7}{8.4}\noexpand\selectfont, text=hpbInk,
        inner sep=0pt] at (axis cs:\c,\y) {\l};}\mxtmp
    }
    % baseline tags above the panel
    \foreach \c/\v in {1/v1.0, 5/v1.0.1, 10/v1.1, 16/v1.2}{
      \edef\mxtmp{\noexpand\node[hpb/tag, anchor=south, yshift=3pt]
        at (axis cs:\c,4.3) {\v};}\mxtmp
    }
    % annotations
    \draw[hpb/leader] (axis cs:9.7,2.2) |- (axis cs:8.6,3.72)
      node[mx/ann, anchor=east]
      {release of v1.1: I-21 and X-102\\close together ($|C|=2$)};
    \draw[hpb/leader] (axis cs:15.7,2.2) |- (axis cs:15.1,3.72)
      node[mx/ann, anchor=east]
      {batch merge at v1.2:\\three branches close ($|C|=3$)};
  \end{axis}
  % ------------------------------------------------------------------ (b)
  \begin{axis}[mx/common, at={(wip.south west)}, anchor=north west,
      yshift=-0.55cm, height=3.6cm,
      ymin=0.4, ymax=6.6, ytick={1,...,6}, y dir=reverse,
      yticklabels={I-21, I-22, X-101, X-102, X-103, T-8},
      y tick label style={font=\sffamily\bfseries\fontsize{7}{8.4}\selectfont},
      major y tick style={draw=none},
      xticklabels={Jan, Feb, Apr, Jul},
      xlabel={ordinal time (one column per event; calendar month of each baseline)},
      ylabel={branch},
    ]
    % one bar per branch: from opening t_o(b) (start node) to closing t_c(b)
    % (release of the merge baseline);  #1 row, #2 type color, #3 letter,
    % #4 t_o column, #5 t_c column, #6 fork -> merge baseline, #7 D(b)
    \foreach \r/\tc/\tl/\o/\e/\fm/\dd in {
        1/hpbImp/I/3/10/{v1.0 $\to$ v1.1}/2,
        2/hpbImp/I/11/16/{v1.1 $\to$ v1.2}/1,
        3/hpbIss/X/2/5/{v1.0 $\to$ v1.0.1}/1,
        4/hpbIss/X/6/10/{v1.0.1 $\to$ v1.1}/1,
        5/hpbIss/X/12/16/{v1.1 $\to$ v1.2}/1,
        6/hpbTask/T/7/16/{v1.0.1 $\to$ v1.2}/2}{
      \edef\mxtmp{%
        \noexpand\draw[mx/bar, fill=\tc!55]
          (axis cs:\o,\r-0.27) rectangle (axis cs:\e,\r+0.27);
        \noexpand\node[hpb/event, fill=\tc] at (axis cs:\o,\r) {\tl};
        \noexpand\node[mx/in, anchor=west] at (axis cs:\o,\r) [xshift=9pt]
          {\unexpanded\expandafter{\fm}};
        \noexpand\node[mx/in, font=\noexpand\sffamily\noexpand\bfseries
          \noexpand\fontsize{7}{8.4}\noexpand\selectfont]
          at (axis cs:17.55,\r) {\dd};}%
      \mxtmp
    }
    \node[mx/in, font=\sffamily\bfseries\fontsize{7}{8.4}\selectfont]
      at (axis cs:17.55,0.4) [anchor=south, yshift=1pt] {$D(b)$};
  \end{axis}
  % panel letters
  \node[font=\sffamily\bfseries\fontsize{7.5}{9}\selectfont, anchor=north east]
    at ([xshift=-1.15cm]wip.north west) {(a)};
  \node[font=\sffamily\bfseries\fontsize{7.5}{9}\selectfont, anchor=north east]
    at ([xshift=-1.15cm,yshift=-0.55cm]wip.south west) {(b)};
\end{tikzpicture}}
\caption{Flow metrics read straight from the graph, for the data of
situation~A (\cref{fig:software-line}) on the same ordinal columns. (a)~Work
in progress per type: each opening raises a step, each closure lowers it, and
a release drops it by its batch size $|C(v)|$. (b)~One bar per branch, from
its opening to the release that closed it, with its fork and merge versions
and its merge distance $D(b)$. Lead times in calendar time come from the
event dates, not from the length of the bars.}
\label{fig:metrics}
\end{figure}


\paragraph{Quality and maintenance.}
The release batch is exactly the release-note list, and smaller batches go
with faster feedback \citep{reinertsen2009flow}. Escaped defects need three
cautions: \emph{seen in} is not \emph{introduced by}; counts grow with how
widely a version is deployed, so normalize by installed base where you
have it; and all of it assumes the fork baseline of every issue is the
version where it was really seen, which is why intake makes it mandatory
(\cref{subsec:method-practices}). Gate slip only works if the original
target date is kept when the date changes. Maintenance load and unported
fixes keep old references and maintenance streams honest.

\subsection{Reading metrics responsibly}
\label{subsec:metrics-goodhart}

Every metric above can be improved without improving the product. Lead
time drops if branches are opened late or split artificially. The discard
ratio drops if hopeless work is kept open. Escaped defects drop if issues
are filed against the wrong version. Gate slip disappears if scope is cut
silently. This is Goodhart's law: when a measure becomes a target, it stops
being a good measure. Three safeguards:

\begin{enumerate}[leftmargin=*,itemsep=2pt]
  \item \emph{Report in pairs} that are hard to game together: lead time
    with throughput, gate slip with descoped count, merge ratio with age
    of open work.
  \item \emph{Always show the diagram next to the numbers}, so every
    anomaly can be traced to the branches that cause it.
  \item \emph{Do not use the metrics to compare or reward people or
    teams.} Rules R1--R7 are what make the numbers trustworthy; a regime
    that rewards breaking them destroys its own data.
\end{enumerate}

\noindent The metrics describe how change flows through a product. They do
not measure customer value, effort or the quality of decisions; like
value-stream mapping \citep{rother1999lts}, they make the flow visible in
order to discuss it, not to replace the discussion.

% ---------------------------------------------------------------------
%  09-tooling.tex -- HPB Markup, the open-source renderer and the
%  prototype in a product management portal. The layout algorithm, the
%  data model and the implementation status are in appendix-a-formal.tex
%  (app:layout, app:impl).
% ---------------------------------------------------------------------
\section{Tooling: HPB Markup and an Open-Source Renderer}
\label{sec:tooling}

% plain listings language for code snippets (HTML, TSX, shell)
\lstdefinelanguage{hpbplain}{}

A notation only spreads if drawing it costs nothing. Every diagram in this
paper was drawn by hand with \texttt{hpb.sty} (\cref{app:package}). That is
fine for a paper. It is not fine for a product line that gets a new branch
every few days: nobody redraws a diagram by hand every week, and a diagram
that is not redrawn is soon wrong.

So HPB comes with tooling. \emph{HPB Markup} is a small text language in
the spirit of Mermaid's \texttt{gitGraph}~\citep{mermaid}: you write a few
lines of text and get an HPB diagram. A set of open-source packages renders
that text as SVG in a plain HTML page, in React, Next.js or Vite, in
Markdown documentation and from the command line. The toolkit is released
under the MIT license; its code, documentation and the canonical
situations of \cref{sec:practice} as test fixtures are published at
\hpbrepo.\footnote{Package names in this section are those of version~1.0
of the toolkit. If a name has to change when the packages are published,
the repository lists the published names.} At the end of this section I
describe the prototype I built inside the product management portal I
work with, which draws the same diagrams directly from the trackers.

% =====================================================================
\subsection{Why a text language}
\label{sec:tooling-why}

Product documentation increasingly lives as text: Markdown pages, README
files, wikis, pull requests. Text has three properties that pictures do not.
It can be \emph{diffed}, so a reviewer sees that \texttt{X-212} was moved
from \texttt{v2.1} to \texttt{v2.2} instead of comparing two images. It can
be \emph{generated}, so a script can turn an export from any tracker into an
HPB diagram. And it can be \emph{checked}, so the rules of the method are
applied every time the diagram is drawn, not only when someone remembers
them.

The last point matters most to me. A generic drawing tool will happily
show a branch that never closes, or a release that ships with open
committed work. An HPB Markup document that does either produces a
diagnostic. The language is not only a faster way to draw: it makes
\cref{rule:origin,rule:closure,rule:time,rule:gate} hard to break by
accident.

% =====================================================================
\subsection{Design principles}
\label{sec:tooling-principles}

\begin{itemize}[leftmargin=1.5em,itemsep=2pt,topsep=3pt]
  \item \textbf{Statement order is time order.} Every statement that
        creates an event gets the next column, like the ordinal columns of
        the notation (\cref{sec:notation}). Dates such as \texttt{@2026-03}
        are optional: they only feed the calendar ticks and are checked to
        be non-decreasing.
  \item \textbf{A sensible default fork, an explicit one for issues.}
        Without \texttt{from}, a branch forks from the latest released
        baseline of its line, which is the default of \cref{def:branch} for
        improvements, tasks and lab work. For an issue, \cref{rule:origin}
        requires the baseline in which it was observed: write it with
        \texttt{from v1.4} unless it is the latest release.
  \item \textbf{One terminal per branch.} A branch ends with a merge or a
        discard, or stays open and is drawn up to NOW
        (\cref{rule:closure}).
  \item \textbf{The rules are checked, not assumed.} Merging into a
        baseline that does not succeed the fork baseline, including the
        fork itself, is an error (\cref{rule:time}). Releasing a planned
        baseline that still has open targeted branches is an error
        (\cref{rule:gate}), unless the release is marked \texttt{force}: the
        violation is then recorded and stays visible as a warning badge.
  \item \textbf{Forgiving to write.} Keywords are case-insensitive and have
        aliases (\texttt{bug} for \texttt{issue}, \texttt{feature} for
        \texttt{improvement}); identifiers can be any version or reference
        such as \texttt{v1.0.1} or \texttt{Rev-B}; \texttt{\#} and
        \texttt{\%\%} start comments; indentation is ignored.
  \item \textbf{Clear errors.} Every diagnostic has a stable code, a line,
        a column and a plain message. By default the renderer draws an
        error box instead of throwing, so a broken diagram never breaks the
        page around it.
  \item \textbf{One visual contract.} The renderer draws the glyphs and
        palette of \cref{sec:notation}, the same ones \texttt{hpb.sty}
        draws, so a diagram on a web page and one in a printed report read
        the same way.
\end{itemize}

% =====================================================================
\subsection{A complete example}
\label{sec:tooling-example}

\Cref{fig:markup-example} shows a firmware line written in HPB Markup and
the diagram it produces. It is the same situation as
\cref{fig:software-line}, with different illustrative items. Line by line:

\begin{itemize}[leftmargin=1.5em,itemsep=1pt,topsep=3pt]
  \item \texttt{hpb} gives the diagram a title, and \texttt{line} declares
        the product line \emph{Product A}.
  \item \texttt{release v1.0} creates the first released baseline.
  \item \texttt{improvement I-21} and \texttt{issue X-101} open two
        branches. Neither says where it starts, so both fork from
        \texttt{v1.0}, the latest release. The improvement goes above the
        main line, the issue below it. X-101 has no title: its short life
        leaves no room for a label (layout rule~\ref{lay:labels}).
  \item \texttt{release v1.0.1 merges X-101} publishes a patch that closes
        the defect. The next two branches fork from \texttt{v1.0.1}, now
        the latest release.
  \item \texttt{release v1.1} and \texttt{release v1.2} close the remaining
        work. T-8 runs across two releases on the task lane, below the
        issue lane.
\end{itemize}

\noindent Two layout conventions of the renderer are worth stating. First,
the merge terminals of all branches closed by one
\texttt{release \dots{} merges} statement share one column, just before
the baseline; they sit on different lanes, which is the exception allowed
by layout rule~\ref{lay:ordinal}. Second, NOW is drawn only when the
diagram has something in the future (planned baselines, open branches, or
merges awaiting release); a pure history such as this one has no NOW
marker. The renderer stacks extra lanes per type, outward from the main
line, and packs them first-fit; the preference of
rule~\ref{lay:onelane} for putting still-open branches on the outer lanes
is not implemented yet (the full algorithm is in \cref{app:layout}).

% F-M -- HPB Markup: the canonical example A (a firmware line) as text and
% as the diagram it produces. Illustrative data. Requires hpb.sty v1.1 and
% the listings language "hpb" (defined in main.tex).
% Column positions are those computed by @hpb/core 1.0 for this source
% (x positions divided by the column width): one column per event in
% statement order; the merge terminals of the branches closed by one
% "release ... merges" statement share one column just before the baseline;
% a dated baseline after a shared terminal column gets a small spacer. The
% document is pure history, so the renderer draws no NOW marker.
%   1 v1.0 | 2 I-21 | 3 X-101 | 4 terminal X-101 | 5 v1.0.1 | 6 X-102 |
%   7 T-8 | 8 terminals I-21, X-102 | 9.4 v1.1 | 10.4 I-22 | 11.4 X-103 |
%   12.4 terminals T-8, I-22, X-103 | 14.4 v1.2
\begin{figure}[tbp]
\centering
\begin{minipage}{0.66\linewidth}
\begin{lstlisting}[style=hpb,basicstyle=\ttfamily\footnotesize,numbers=left,numberstyle=\tiny\color{black!45},numbersep=6pt,xleftmargin=16pt]
hpb "Product A: firmware line"
line A "Product A"
release v1.0 @2026-01
improvement I-21 "Faster pairing"
issue X-101
release v1.0.1 merges X-101
issue X-102 "Clock drift"
task T-8 "Migrate logs"
release v1.1 @2026-03 merges I-21, X-102
improvement I-22 "Low-power mode"
issue X-103 "Menu translation"
release v1.2 @2026-05 merges T-8, I-22, X-103
\end{lstlisting}
\end{minipage}

\medskip
\hpbfit{%
\begin{hpbgraph}
  % lanes: improvements above, issues and tasks below the main line
  \hpblane{1}{imp}{Improvement (I)}{0}{15}
  \hpblane{-1}{iss}{Issue (X)}{0}{15}
  \hpblane{-2}{task}{Task (T)}{0}{15}
  % main line: pure history, no NOW marker
  \hpbmainline{0}{15}
  \hpborigin{0}{Product A}
  % released baselines
  \hpbrelease{1}{v1.0}
  \hpbrelease{5}{v1.0.1}
  \hpbrelease{9.4}{v1.1}
  \hpbrelease{14.4}{v1.2}
  % branches (fork col, lane, start col, end col, end, merge col)
  \hpbbranch[label={I-21 Faster pairing}]{imp}{1}{1}{2}{8}{merge}{9.4}
  \hpbbranch[label={X-101}]{iss}{1}{-1}{3}{4}{merge}{5}
  \hpbbranch[forkoff=0.6, label={X-102 Clock drift}]{iss}{5}{-1}{6}{8}{merge}{9.4}
  \hpbbranch[forkoff=0.3, mergeoff=0.3, label={T-8 Migrate logs}]{task}{5}{-2}{7}{12.4}{merge}{14.4}
  \hpbbranch[label={I-22 Low-power mode}]{imp}{9.4}{1}{10.4}{12.4}{merge}{14.4}
  \hpbbranch[mergeoff=0.6, label={X-103 Menu translation}]{iss}{9.4}{-1}{11.4}{12.4}{merge}{14.4}
  % ordinal time axis with the calendar ticks given by the @dates
  \hpbaxis{0}{15}{-2.75}
  \hpbtick{1}{Jan 2026}{-2.75}
  \hpbtick{9.4}{Mar}{-2.75}
  \hpbtick{14.4}{May}{-2.75}
\end{hpbgraph}}
\caption{HPB Markup in practice (illustrative). Top: twelve lines of text
that describe a firmware line. Bottom: the diagram they produce, drawn here
with \texttt{hpb.sty} at the column positions the renderer computes for
this source. Nobody types a coordinate: statement order gives the columns,
the branch type gives the lane, and each \texttt{release \dots{} merges}
closes the listed branches (\cref{rule:origin,rule:closure}). The terminals
closed by one release share a column; the history has no NOW marker.}
\label{fig:markup-example}
\end{figure}


The other canonical situations are equally short. A release gate, a
discard with its rationale, and a hardware revision that forks a new
reference are written like this (the full language is in
\cref{app:markup}):
\begin{lstlisting}[style=hpb]
now
plan v2.1 @2026-12 targets I-50, X-211, X-212
discard I-60 "not viable: BOM cost"
release RevB hardware merges L-5, I-30
line RB "DMS-4000 Rev B" from RevB
\end{lstlisting}

\paragraph{What the Markup checks, and what it does not.}
The Markup checks a structural version of the gate rule: no branch
targeted at a baseline may still be open when that baseline is released.
It does not model verification evidence, so ``verified'' in the sense of
\cref{rule:gate} remains a human check at the gate review. A
\texttt{release \dots{} merges} statement plays the role of publication
(\cref{def:gate}), with the listed branches as its targeted set. A
\texttt{force} release records a deliberate violation: the resulting model
is not well formed under \cref{rule:gate}, and the diagram says so with a
\texttt{W-R4-FORCED} warning and a ``!'' badge.

% =====================================================================
\subsection{Packages}
\label{sec:tooling-packages}

The renderer is a small set of TypeScript packages
(\cref{tab:packages}). The core has no runtime dependencies and produces a
plain SVG string, so it runs the same way in a browser, in Node.js and in
server-side rendering.

\begin{table}[tbp]
  \centering
  \small
  \caption{HPB Markup toolkit, version 1.0 (MIT license).}
  \label{tab:packages}
  \renewcommand{\arraystretch}{1.12}
  \begin{tabularx}{\linewidth}{@{}>{\raggedright\arraybackslash}p{2.6cm}
      >{\raggedright\arraybackslash}X
      >{\raggedright\arraybackslash}p{3.5cm}@{}}
    \toprule
    \textbf{Package} & \textbf{What it does} & \textbf{Typical host} \\
    \midrule
    \texttt{@hpb/core}
      & \texttt{parse}, \texttt{validate}, \texttt{buildModel} (checks
        R1--R4, returns diagnostics), \texttt{layout}, \texttt{renderSVG}
        (pure string, no DOM).
      & Any JavaScript runtime, SSR, build scripts. \\
    \addlinespace
    \texttt{@hpb/web}
      & Mermaid-style auto-render of \texttt{pre.hpb} and
        \texttt{code.language-hpb} blocks; custom element
        \texttt{<hpb-diagram>}; API \texttt{initialize}, \texttt{run},
        \texttt{render}; \texttt{dist/hpb.min.js} exposes
        \texttt{window.HPB} for a \texttt{<script>} tag.
      & Plain HTML pages, wikis, static sites. \\
    \addlinespace
    \texttt{@hpb/react}
      & \texttt{<HpbDiagram code=\dots{} />}, SSR-safe; optional hover
        highlighting, \texttt{onNodeClick} and \texttt{onError}.
      & React, Next.js (App Router), Vite. \\
    \addlinespace
    \texttt{@hpb/cli}
      & \texttt{hpb render} to SVG and \texttt{hpb validate} for
        continuous integration.
      & Terminals, CI pipelines, document builds. \\
    \addlinespace
    \texttt{remark-hpb}
      & remark (unified) and markdown-it plugins that render
        \texttt{hpb} fenced code blocks as inline SVG.
      & Markdown and MDX, documentation sites, static generators. \\
    \bottomrule
  \end{tabularx}
\end{table}

\paragraph{In a plain HTML page.}
Load one script and write the diagram as text. Every matching block is
replaced by its SVG when the page loads.
\begin{lstlisting}[language=hpbplain]
<script src="https://unpkg.com/@hpb/web@1.0.0/dist/hpb.min.js"></script>
<pre class="hpb">
line A "Product A"
release v2.0
issue X-211 "Battery drain"
release v2.0.1 merges X-211
</pre>
\end{lstlisting}

\paragraph{In React, Next.js or Vite.}
The component renders the SVG string during render and never touches
\texttt{window}, so it works in server-side rendering. In the Next.js App
Router, a server component can call \texttt{renderSVG} directly; the
interactive component is a client component.
\begin{lstlisting}[language=hpbplain]
import { HpbDiagram } from "@hpb/react";

<HpbDiagram code={source} theme="default"
            onNodeClick={(n) => openItem(n.id)} />
\end{lstlisting}

\paragraph{From the command line.}
\texttt{hpb validate} returns a non-zero exit code when a document breaks
a rule, which makes it usable as a check in a documentation pipeline.
\begin{lstlisting}[language=hpbplain]
hpb render product-a.hpb -o product-a.svg --theme mono
hpb validate product-a.hpb
\end{lstlisting}

\paragraph{Accessibility and print.}
The SVG carries a title and a description, every node has its own
\texttt{<title>}, and no text is smaller than 10\,px at scale one. The
\texttt{mono} theme renders in grayscale with letters and line styles only,
which is the redundant encoding of \cref{sec:notation} put to work.

% =====================================================================
\subsection{Prototype implementation in a product management portal}
\label{sec:implementation}

A method that needs a new tool, and a migration of all existing work into
it, rarely gets adopted. I did not want HPB to be one more system to keep
in sync, so I built it as a thin layer on top of the trackers in which
improvements, defects, tasks and laboratory evaluations were already
managed, as a module of the product management portal I use daily.

\paragraph{An overlay, not a migration.}
The trackers stay exactly as they are and remain the system of record for
the work itself. HPB adds a few tables that store only what the trackers
do not know: the product lines and their family tree, the baselines, the
attachment of each tracker item to a line as a branch with its target and
terminal, and comments attached to exactly one node
(\cref{app:impl}). Because tracker items are referenced and never copied,
adopting HPB needs no migration, and removing it leaves the trackers
untouched.

\paragraph{Derived mode: a picture on day one.}
The same property gives a read-only \emph{derived mode} for products
whose overlay is still empty. Baselines are rebuilt from existing release
records, branches from the tracker items linked to the product, and
terminals from tracker statuses. A completed item with no link to a
release is drawn as closed without a baseline: a violation of
\cref{rule:closure} that the diagram shows instead of hiding. On the first
day, derived mode gives a retrospective picture of the product and a list
of the gaps the method then closes.

\paragraph{The diagram as the place where decisions happen.}
In the portal the diagram is not just a view. One action on a line
creates a typed item in its tracker and attaches it as a branch, so it has
a type and an origin from the first second
(\cref{rule:origin,rule:typed}). Targeting offers only planned baselines;
a discard asks for a rationale; publishing a planned baseline lists
everything targeted at it and merges the set in one operation; a hardware
baseline that changes the reference records the new line it creates
(\cref{rule:lineage}); and every line, baseline and branch has its own
comment thread and a link back to the item in its tracker
(\cref{rule:discussion}). Several rule checks are designed but not yet
enforced in the prototype; \cref{app:impl} lists them honestly. Until they
are, the violations are at least visible in the diagram, and HPB Markup
already enforces R1--R4 for anyone who draws with it.

% ---------------------------------------------------------------------
%  10-limitations.tex -- limits of the method and how to evaluate it
% ---------------------------------------------------------------------
\section{Limitations and How to Evaluate HPB}
\label{sec:limitations}

HPB comes out of years of running real products, but this paper does not
prove that it works better than what you use today. It offers a method, a
notation and tools, and explains why I believe they help. This section says
where HPB fits, where it does not, what can go wrong, and how to find out
whether it pays off in your own organization.

% =====================================================================
\subsection{Where HPB fits, and where it does not}
\label{sec:limitations-fit}

HPB earns its keep when three things are true at once: the product has
\emph{identifiable versions} that customers or devices actually run;
\emph{several pieces of work} are in flight on the same product; and people
need to know \emph{which version} contains a given change. That is the daily
reality of products that combine hardware, firmware and software, of
software shipped in discrete versions with long support, and of regulated
contexts, where HPB supports formal configuration management
\citep{ieee828,iso10007} without replacing it.

HPB adds little in four situations. With \emph{continuous deployment
without versions} \citep{humble2010cd}, every merge is its own release and
baselines lose their meaning; a reduced timeline with deployments as
baselines is possible, but I do not claim it beats a deployment log. When
a product has \emph{one change at a time}, the diagram is a straight line
and a changelog is enough. \emph{Before the first release} there is
nothing to fork from; early exploration can live as lab branches of a
provisional line. And \emph{portfolio decisions}, how to spread people and
money across many products, need aggregated views the method does not
define yet.

% =====================================================================
\subsection{Known limitations}
\label{sec:limitations-known}

\paragraph{No empirical evidence yet.}
I report no data showing that HPB improves planning accuracy, traceability,
quality or decision speed, and none should be read into this paper. The
benefits I describe are hypotheses that follow from the structure of the
method; \cref{sec:limitations-evaluate} explains how to test them.

\paragraph{Designer bias.}
I designed HPB from my own practice and built the prototype in that same
context. The method may reflect the needs and vocabulary of that context
more than I can see, and the comparison in \cref{sec:gaps} rates HPB by
design, not by observed use.

\paragraph{Scale.}
The notation is designed for one product line with tens of open branches.
Collapsed history and lane packing (\cref{fig:scale}) stretch that range,
but screen width and human attention are still limits. A line with
hundreds of open branches needs filters, and a filtered diagram no longer
shows everything at a glance.

\paragraph{Garbage in, garbage out.}
Every metric and every audit answer depends on three fields being right:
the fork baseline, the closing baseline and the type of each branch. The
fork baseline of issues is the weak spot: when nobody records the version
in which a defect was observed, the default quietly blames recent
releases. HPB reduces this risk by making someone accountable for these
fields at intake and by flagging inferred values; it does not remove it.

\paragraph{Culture.}
HPB makes decisions visible, including the decision not to do something
and the decision to ship late. Where that feels threatening, people comply
on the surface: they never discard, or they re-target silently. No method
fixes that by itself; HPB can only make the symptoms show in the diagram.

\paragraph{Deliberate simplifications.}
A branch has one fork baseline and closes into one baseline
(\cref{rule:origin,rule:closure}), so a fix that must reach several
maintained versions is one branch per version, linked as propagations,
much like cherry-picking \citep{chacon2014progit}: change logs stay exact,
but widely back-ported fixes multiply branches. A line has one parent, so
merging references or a family of variants on a shared platform is
recorded as annotations and sibling lines. Security, regulatory or
customer-specific work are attributes, not new lanes. And for hardware the
moment of ``release'' is not crisp; HPB fixes a default per baseline kind,
and an organization that uses another must say so, or its lead times will
not be comparable.

% =====================================================================
\subsection{How to evaluate HPB in your organization}
\label{sec:limitations-evaluate}

You do not need a research project to find out whether HPB helps you. You
need one product line, two release cycles and the discipline to compare
before and after honestly.
\begin{enumerate}[leftmargin=1.8em,itemsep=1pt,topsep=3pt]
  \item \textbf{Take a baseline.} Rebuild the last year of one line in
        derived mode (\cref{sec:implementation}) or in HPB Markup, flag
        every fork baseline you had to guess, and compute the metrics you
        can.
  \item \textbf{Ask ten questions first.} Write down ten questions your
        team gets asked (which version first contained this fix, what is
        still open for the next release), time the answers with your
        current tools and check them against your release records.
  \item \textbf{Run HPB for two release cycles} with the practices of
        \cref{sec:operating}, without changing your team process at the
        same time.
  \item \textbf{Ask an equivalent set of questions again} with the
        diagram, built from the same data, so HPB does not win just because
        its data is cleaner.
  \item \textbf{Look at outcomes and cost:} rule violations, gate slip with
        what was moved or discarded, age of open work, escaped defects, and
        the time spent keeping HPB up to date.
\end{enumerate}

\noindent Four hypotheses are worth testing, each with the result that
would count against HPB:
\begin{description}[leftmargin=0pt,itemsep=1pt,font=\normalfont\bfseries]
  \item[H1 Comprehension.] Questions about origin, parallel work and
        closure are answered more accurately, and not more slowly, with the
        diagram. \emph{Against:} no gain, or a gain paid with much longer
        answers.
  \item[H2 Discipline.] Violations of R1--R4 fall to a low, stable rate.
        \emph{Against:} silent closures and releases with unverified
        committed work that stay frequent.
  \item[H3 Decisions.] More slips and descopes are recorded as decisions,
        and fewer branches stay open forever. \emph{Against:} unchanged age
        profiles, descopes still found out after the release.
  \item[H4 Cost.] Keeping HPB up to date costs little compared with the
        meetings it uses. \emph{Against:} data entry the team judges out of
        proportion, or quiet abandonment after the trial.
\end{description}

\noindent With few releases a year, differences between baselines are
mostly noise, so look at trends; everything looks better in the first
month, so compare the first cycle with the last; and never use these
numbers to judge individuals. If your evaluation shows that HPB does
\emph{not} help in your context, I would especially like to hear about it.

% =====================================================================
\subsection{What comes next}
\label{sec:limitations-next}

Four directions follow from these limits: a \emph{portfolio view} that
aggregates many lines sharing platforms or teams without breaking
\cref{rule:origin,rule:closure}; \emph{automatic links} to version control
and build systems, so that release tags, build numbers and hardware
revision records create baselines and merges by themselves, with semantic
versioning \citep{semver} giving the size of a software baseline; an
\emph{interchange format}, for which HPB Markup is the first candidate; and
a \emph{versioned specification} of the notation, changed only on evidence
about readability, not on taste. The conclusion says how to contribute.

% ---------------------------------------------------------------------
%  11-conclusion.tex
% ---------------------------------------------------------------------
\section{Conclusion}
\label{sec:conclusion}

Software teams learned long ago that parallel change can only be managed,
and audited, if every line of work has a known starting point and an
explicit moment at which it rejoins the main line. Product management never
got the same discipline. Our tools tell us what is being done and in which
workflow state. They do not tell us from which version of the product the
work started, what else was in flight on the same product, or which
firmware release or hardware revision finally closed it. After more than
seven years of managing hardware, firmware and software products, that is
the gap I kept running into, and HPB is my answer to it.

HPB takes the branch-and-merge model from source code and applies it to the
product. It gives you:
\begin{itemize}[leftmargin=1.5em,itemsep=1pt,topsep=3pt]
  \item a \emph{model} of product lines, baselines, typed branches, merges,
        discards, line forks and release gates, governed by seven rules
        (\cref{sec:onepage,sec:model,app:formal});
  \item a \emph{standard notation}, so that everybody draws the same
        situation the same way (\cref{sec:notation,app:card});
  \item an \emph{operating guide} and \emph{metrics} that fall out of the
        graph (\cref{sec:operating,sec:metrics});
  \item \emph{HPB Markup} and an open-source renderer, so a diagram is a
        few lines of text in any web stack or document
        (\cref{sec:tooling,app:markup}); and
  \item a \emph{prototype} in a product management portal that sits on
        top of existing trackers instead of replacing them
        (\cref{sec:implementation}).
\end{itemize}

HPB does not replace Scrum, Kanban, Stage-Gate or formal configuration
management. It sits above them and makes one thing visible that none of
them shows: the product itself, as it evolves through released and planned
versions, with every piece of work anchored to the version it started from
and the version that closed it.

Whether that visibility improves planning and quality in your organization
is something you can test in two release cycles
(\cref{sec:limitations-evaluate}). The notation, the \LaTeX{} package and
the renderer are open, so the method can be adopted, criticized and
improved in a common visual language.

\paragraph{An invitation.}
Try HPB on one product line for two release cycles. Then share what you
found, the diagram that worked or the counter-example that broke the
method, through the issue tracker of \hpbrepo{} or by email to the
address on the first page. Changes to the method and the notation go
through versioned releases (HPB~1.x), so a diagram drawn today stays
readable tomorrow. Start with one line, draw it honestly, and let the
diagram show you what your current tools have been hiding.

% ---------------------------------------------------------------------
%  12-about-author.tex -- unnumbered "About the author"
% ---------------------------------------------------------------------
\section*{About the Author}
\phantomsection
\addcontentsline{toc}{section}{About the Author}
\label{sec:author}

\noindent\textbf{Achraf Bouchtarouif El Khiat} is a senior product manager
with more than seven years of experience managing complete products end to
end: hardware, firmware and software. His work covers a large portfolio
with high rotation of references, continuous improvement requests,
competitor analysis and a constant, fast-moving evolution of the products
in the market. HPB grew out of the gaps he found in today's methods while
running that portfolio; he built a working prototype of it in the
product management portal he uses, and wrote the open-source HPB Markup
toolkit.

\medskip
\noindent Contact:
\href{mailto:achrafbouchtarouif@gmail.com}{\texttt{achrafbouchtarouif@gmail.com}}.
HPB Markup: \hpbrepo.


% ========================================================== bibliography
\bibliographystyle{plainnat}
\bibliography{references}

% ============================================================ appendices
\clearpage
\appendix
% ---------------------------------------------------------------------
%  appendix-a-formal.tex -- the formal HPB model
%  Single place where Definitions, Rules R1--R7 (labels rule:*) and
%  Properties P1--P4 (labels prop:*) are stated. The main text
%  (sections 03, 04, 05) explains them in plain language.
% ---------------------------------------------------------------------
\section{Formal Model}
\label{app:formal}

This appendix states HPB precisely, for readers who want to check the
method or build a tool that enforces it. \Cref{sec:model} explains the same
concepts in plain language; the rules R1--R7 are introduced, with their
practical motivation, in \cref{sec:onepage}. The last three subsections
give the exact definitions of the metrics, the layout algorithm, and the
data model and status of the prototype. Time is taken from a totally
ordered domain $\mathbb{T}$ (calendar time); how time is \emph{drawn} is a
separate matter (\cref{sec:notation,def:ordinal}). Concepts that only some
organizations need (configurations, dependencies, launch bundles and the
deployment overlay) can be left out without affecting the rest of the
model.

\subsection{Lines and baselines}
\label{app:formal-lines}

\begin{definition}[Product line, line forest]
  \label{def:line}
  A \emph{product line} $P$ is an identified product reference, such as a
  model together with its manufacturer reference, or a firmware or
  software product. A portfolio has a finite set of lines $\mathcal{P}$
  and a partial function
  $\mathrm{parent}\colon \mathcal{P} \rightharpoonup \mathcal{P}$, where
  $\mathrm{parent}(P') = P$ states that $P'$ was created by a hardware
  revision of $P$ (\cref{def:linefork}). The graph with an edge from
  $\mathrm{parent}(P')$ to $P'$ is a forest, the \emph{line forest}.
\end{definition}

\begin{definition}[Baseline]
  \label{def:baseline}
  A \emph{baseline} $v$ is a named, identifiable state of a product line
  $\ell(v) \in \mathcal{P}$. It carries a name, unique within its line
  (e.g., \texttt{v2.3.1} or \texttt{Rev~B}), a kind $k(v)$, one of
  \textsf{firmware}, \textsf{software}, \textsf{hardware} or
  \textsf{other}, a status $s(v)$, either \textsf{planned} or
  \textsf{released}, and a date $t(v) \in \mathbb{T}$, which is the
  release date if $v$ is released and the target date if it is planned.
  We write $\mathcal{V}$ for the set of all baselines and
  $\mathcal{V}(P) = \{v \in \mathcal{V} : \ell(v) = P\}$. A status may
  change only from \textsf{planned} to \textsf{released}; once released,
  a baseline is immutable.
\end{definition}

The kind \textsf{other} is used for process baselines. The normative
release event per kind, and the form--fit--function criterion that decides
whether a hardware change is a baseline of an existing line or a line
fork, are given in \cref{sec:model-lines}.

\begin{definition}[Release stream, main line, precedence]
  \label{def:mainline}
  The baselines of a line $P$ are partitioned into one or more
  \emph{release streams}. The \emph{main stream} starts at the
  \emph{first baseline} $v_0(P)$ of $P$, i.e., its first release or, for a
  line created by a line fork, the hardware baseline that created it
  (\cref{def:linefork}). Every other stream $R$, a \emph{maintenance
  stream}, is created at a released \emph{stream-fork baseline}
  $\beta(R)$ of another stream of $P$. Within a stream, released baselines
  are ordered by date and followed by the planned baselines of the stream,
  ordered by target date; the sequence of the main stream,
  \[
    M(P) = (v_0, v_1, \dots, v_n, u_1, \dots, u_r),
  \]
  is the \emph{main line} of $P$. \emph{Precedence} $\prec$ is the smallest
  transitive relation on $\mathcal{V}$ such that $x \prec y$ whenever
  (i)~$y$ immediately follows $x$ in a stream, (ii)~$y$ is the first
  baseline of a maintenance stream $R$ and $x = \beta(R)$, or (iii)~$y$ is
  the first baseline of a child line $P'$ and $x = \delta(P')$, the
  baseline of the parent line from which $P'$ was derived
  (\cref{def:linefork}). We write $\preceq$ for its reflexive closure.
\end{definition}

By construction every baseline has at most one immediate predecessor: its
predecessor in its stream, the fork baseline of its stream, or the parent
baseline of its line. Precedence is therefore a forest order, and in
particular a strict partial order. Release dates order baselines only
within a stream.

\begin{definition}[Applies-to relation, configuration]
  \label{def:config}
  A line is a \emph{hardware} line if its reference is a physical unit and
  a \emph{software} line if it is a firmware, application or service
  product; a \emph{combined} line carries the hardware revisions of one
  reference together with the firmware released for it. For every
  baseline $v$ of a software line, the \emph{applies-to} set
  $\mathrm{app}(v) \subseteq \mathcal{P}$ lists the hardware references on
  which $v$ is released to run (its effectivity). A \emph{configuration}
  $\kappa$ is a set of baselines with at most one baseline per line in
  which every software baseline applies to the hardware reference of
  $\kappa$; for example (\texttt{Rev~A}, firmware \texttt{v2.0}, app
  \texttt{5.1}). The relation $\mathrm{app}$ is the compatibility matrix
  of the portfolio.
\end{definition}

\subsection{Branches, states and closure}
\label{app:formal-branches}

\begin{definition}[Branch]
  \label{def:branch}
  A \emph{branch} $b$ is a unit of change intent on a product. We write
  $\mathcal{B}$ for the set of all branches. Every branch has
  \begin{itemize}[leftmargin=1.5em,itemsep=1pt,topsep=2pt]
    \item a \emph{type} $\tau(b) \in \{\mathsf{I}, \mathsf{X}, \mathsf{T},
          \mathsf{L}\}$: improvement, issue (defect), task, or laboratory
          evolution;
    \item an \emph{opening time} $t_o(b) \in \mathbb{T}$, the time at which
          the request is accepted at triage;
    \item a \emph{fork baseline} $f(b) \in \mathcal{V}$ with
          $s(f(b)) = \textsf{released}$ and $t(f(b)) \le t_o(b)$, fixed at
          intake: for an issue, the baseline in which the defect was
          \emph{observed}, as reported; for other types, the baseline
          relative to which the change is specified, by default the latest
          released baseline of the main stream at $t_o(b)$;
    \item for an issue, optionally, the \emph{observed configuration}
          $\Phi(b) \ni f(b)$ when the defect was reported for a combination
          of baselines of several lines; $f(b)$ is then its \emph{primary}
          element, on whose line the branch is drawn and attributed;
    \item for an issue, two attributes set by later analysis: the set
          $\mathrm{aff}(b) \subseteq \mathcal{V}$ of further baselines in
          which the defect is confirmed (\emph{also affects}) and, when root
          cause is known, the baseline $\iota(b)$ that \emph{introduced} it;
    \item a finite, time-ordered sequence of \emph{steps}
          $S(b) = \big((t_1, \kappa_1, e_1), \dots, (t_q, \kappa_q, e_q)\big)$
          with $t_o(b) \le t_1$, kinds
          $\kappa_i \in \{\textsf{done}, \textsf{todo}, \textsf{pass},
          \textsf{fail}\}$ and optional evidence $e_i$ (a report, a
          measurement, the identifier of a build).
  \end{itemize}
\end{definition}

The fork baseline is never revised by later knowledge: confirmation in
other baselines extends $\mathrm{aff}(b)$ and root-cause analysis sets
$\iota(b)$. An error in $f(b)$ itself is corrected only through an audited
transition, and an inferred fork baseline is flagged as such.

\begin{definition}[Branch state, verification]
  \label{def:state}
  The \emph{state} of a branch is a function
  $\sigma\colon \mathcal{B} \times \mathbb{T} \to \{\textsf{open},
  \textsf{targeted}, \textsf{merged}, \textsf{discarded}\}$, defined for
  $t \ge t_o(b)$. A \emph{targeted} branch is open and committed to a
  \emph{target baseline} $g(b)$, which must be planned when the target is
  set. States change only through the transitions of \cref{fig:lifecycle}:
  \emph{target} and \emph{re-target} (set or change $g(b)$),
  \emph{un-target} (clear $g(b)$), \emph{merge}, \emph{discard} and, as an
  exceptional and audited operation, \emph{reopen}. The states
  \textsf{merged} and \textsf{discarded} are \emph{terminal}; the time of
  the most recent terminal transition is the \emph{closing time} $t_c(b)$,
  and $t_c(b) = \infty$ while $b$ is open or targeted. A targeted branch is
  \emph{verified} at time $t$ if $S(b)$ contains a \textsf{pass} step at a
  time $\le t$ whose evidence identifies a candidate of $g(b)$ (a release
  candidate build, a prototype or a pilot lot); the time of the latest such
  step is its \emph{verification time} $t_r(b)$.
\end{definition}

\begin{definition}[Merge and discard]
  \label{def:merge}
  \emph{Merging} a targeted branch $b$ sets its \emph{merge baseline}
  $m(b) = g(b)$ and establishes the relation $\mathrm{closes}(b, m(b))$:
  baseline $m(b)$ contains the resolution of $b$. Merges happen only when
  $m(b)$ is published (\cref{def:gate}), so $t_c(b) = t(m(b))$ for every
  merged branch. \emph{Discarding} a branch closes it without a baseline
  and records a non-empty \emph{rationale} $\rho(b)$ from one of two
  classes: \emph{abandoned} (not viable, duplicate of $b'$, cannot
  reproduce, superseded, out of scope, will not fix) or \emph{resolved
  without a product baseline} (already resolved in $v$, contained without
  design change, resolved by the supplier).
\end{definition}

\subsection{Line forks, life cycle, gates and relations}
\label{app:formal-forks}

\begin{definition}[Line fork]
  \label{def:linefork}
  A \emph{line fork} is a baseline $h$ with $k(h) = \textsf{hardware}$ that
  creates a new product reference $P'$, with $\mathrm{parent}(P') = P$.
  The baseline $h$ belongs to the child line, $\ell(h) = P'$, and is its
  first baseline, $v_0(P') = h$. The \emph{derived-from baseline}
  $\delta(P') \in \mathcal{V}(P)$ is the latest released baseline of the
  main stream of $P$ before $t(h)$; it is the only immediate predecessor of
  $h$ (\cref{def:mainline}).
\end{definition}

Branches that produce the new reference fork from baselines of $P$ and
merge into $h$, which is admitted by the precedence edge
$\delta(P') \prec h$. Because $h$ is not a baseline of the parent line,
later baselines of $P$ do not contain that work.

\begin{definition}[Line life cycle]
  \label{def:lifecycle}
  Every line $P$ has a sequence of dated \emph{milestones}
  \[
    t_{\mathrm{sale}}(P) \le t_{\mathrm{prod}}(P) \le t_{\mathrm{sec}}(P)
    \le t_{\mathrm{eos}}(P)
  \]
  (end of sale, end of production, start of security-only maintenance,
  end of service), each possibly undefined while it has not been decided.
  They divide the life of the line into the phases \emph{active},
  \emph{no longer sold}, \emph{no longer produced}, \emph{security only}
  and \emph{retired}. \Cref{tab:phases} states which branches may fork from
  the baselines of a line, and which baselines it may receive, in each
  phase.
\end{definition}

\begin{table}[tbp]
  \centering
  \small
  \caption{What a line still accepts in each phase of its life
    (\cref{def:lifecycle}). These are defaults: an organization may
    tighten them; relaxing them is an explicit decision recorded on the
    line.}
  \label{tab:phases}
  \begin{tabularx}{\linewidth}{@{}>{\raggedright\arraybackslash}p{3.2cm}
      >{\raggedright\arraybackslash}X
      >{\raggedright\arraybackslash}X@{}}
    \toprule
    \textbf{Phase (from milestone)} & \textbf{Branches that may fork}
      & \textbf{Baselines that may be added} \\
    \midrule
    Active & $\mathsf{I}$, $\mathsf{L}$, $\mathsf{X}$, $\mathsf{T}$
      & all kinds \\
    No longer sold (end of sale)
      & $\mathsf{X}$, $\mathsf{T}$; $\mathsf{I}$ only for recorded customer
        commitments; $\mathsf{L}$ only for supplier-driven changes
      & firmware and software; hardware and process only for
        supplier-driven changes \\
    No longer produced (end of production)
      & $\mathsf{X}$, $\mathsf{T}$
      & firmware and software, preferably on a maintenance stream \\
    Security only
      & $\mathsf{X}$ with the security attribute
      & security patches \\
    Retired (end of service) & none & none \\
    \bottomrule
  \end{tabularx}
\end{table}

\begin{definition}[Gate, publication]
  \label{def:gate}
  Every planned baseline $g$ acts as a \emph{gate}. Its \emph{targeted
  set} at time $t$ is
  \[
    \mathcal{B}_t(g) = \{\, b \in \mathcal{B} : \sigma(b,t) =
    \textsf{targeted} \ \wedge\ g(b) = g \,\},
  \]
  written $\mathcal{B}(g)$ when the time is clear from context.
  \emph{Publishing} $g$ at time $t$ is an atomic operation that sets
  $s(g) = \textsf{released}$ and $t(g) = t$ and merges every branch of
  $\mathcal{B}_t(g)$ into $g$. Before publication, every member of
  $\mathcal{B}_t(g)$ must be verified (\cref{def:state}); members that are
  not are re-targeted to another planned baseline or discarded first
  (\cref{rule:gate}).
\end{definition}

\begin{definition}[Dependency, propagation, launch bundle]
  \label{def:relations}
  The \emph{dependency} relation
  $\mathrm{dep} \subseteq \mathcal{B} \times (\mathcal{B} \cup
  \mathcal{V})$ states that a branch can be delivered only together with
  or after another branch, or only once a baseline (possibly of another
  line) is released. The \emph{propagation} relation
  $\mathrm{prop} \subseteq \mathcal{B} \times \mathcal{B}$ states that
  $b'$ carries the change of $b$ to another stream or line (a back-port or
  forward-port); then $\tau(b') = \tau(b)$. A \emph{launch bundle} is a set
  $\mathcal{G}$ of planned baselines of distinct lines that are published
  together, in one atomic operation, or not at all.
\end{definition}

\begin{definition}[Deployment overlay]
  \label{def:deployment}
  An HPB model may carry an optional \emph{deployment overlay}: for every
  released baseline $v$ and time $t$, a \emph{rollout status}
  $\mathrm{ro}(v,t) \in \{\textsf{pilot}, \textsf{staged}, \textsf{general}\}$
  with, for staged rollouts, the share of the fleet reached; an
  \emph{installed base} $N(v,t) \in \mathbb{N}$ (units or installations
  running $v$); for each customer $c$ of a finite set $\mathcal{K}$, its
  deployed configuration $d(c,t)$; and, for branches, \emph{commitments}
  $\mathrm{cm}(b) \subseteq \mathcal{K} \times (\mathcal{V} \cup
  \mathbb{T})$, i.e., the customers to which $b$ was promised and the
  baseline or date promised.
\end{definition}

For an issue $b$, let $\mathrm{Aff}(b) = \{f(b)\} \cup \mathrm{aff}(b)$
be the baselines in which it is known to occur. The customers exposed to
$b$ at time $t$ are
\[
  \mathcal{K}(b,t) = \{\, c \in \mathcal{K} :
  d(c,t) \cap \mathrm{Aff}(b) \neq \emptyset \,\}.
\]
Customers who have not upgraded to a baseline that contains a fix, and
commitments whose promised baseline has slipped, are obtained in the same
way.

\begin{definition}[HPB model and graph]
  \label{def:model}
  An \emph{HPB model} is a structure
  \[
    H = (\mathcal{P}, \mathrm{parent}, \mathcal{V}, \ell, k, s, t, \beta,
    \delta, \mathrm{app}, \mathcal{B}, \tau, t_o, f, \Phi, \mathrm{aff},
    \iota, S, \sigma, g, m, \rho, \mathrm{dep}, \mathrm{prop},
    \mathbb{G}, A, \alpha),
  \]
  where $\beta$ also determines the streams of every line, $\mathbb{G}$ is
  the set of launch bundles, $A$ is a set of annotations (comments,
  decisions, evidence and approvals) and
  $\alpha\colon A \to \mathcal{P} \cup \mathcal{V} \cup \mathcal{B}$
  attaches each annotation to a node. The relation $\mathrm{closes}$ is
  derived from $m$, and the deployment overlay, when present, is attached
  to $H$. Its \emph{graph} has the baselines and branches as nodes and
  three kinds of structural edges: \emph{precedence} edges between
  immediately preceding baselines (\cref{def:mainline}), \emph{fork} edges
  $f(b) \to b$ and \emph{merge} edges $b \to m(b)$. Dependencies and
  propagations are drawn as annotation edges.
\end{definition}

\subsection{Well-formedness rules}
\label{sec:model-rules}

An HPB model is \emph{well formed} if it satisfies the following rules at
all times. Each rule is followed by the failure it prevents.

\begin{hpbrule}[Single origin]
  \label{rule:origin}
  Every branch has exactly one fork baseline $f(b)$, which is released,
  satisfies $t(f(b)) \le t_o(b)$, is fixed at intake, and is changed only
  by an audited correction.
\end{hpbrule}
\noindent Without a unique origin, ``relative to which product state?''
has no answer and defects cannot be attributed to versions
(\cref{prop:escaped}). Later confirmations extend $\mathrm{aff}(b)$; they
do not create additional forks.

\begin{hpbrule}[Explicit closure]
  \label{rule:closure}
  Every closed branch has exactly one terminal: either a merge into exactly
  one baseline $m(b)$, or a discard with a non-empty rationale $\rho(b)$.
  No branch is closed silently.
\end{hpbrule}
\noindent ``Closed'' without a closing version or a reason is the most
common loss of traceability in practice. A change that must reach several
streams or lines is modeled as one branch per target, linked by
$\mathrm{prop}$, so that a missing propagation is detectable
(\cref{sec:operating}).

\begin{hpbrule}[Time moves forward]
  \label{rule:time}
  For every merged branch, $f(b) \prec m(b)$ and $t(m(b)) \ge t_o(b)$.
\end{hpbrule}
\noindent A baseline cannot contain work requested after it was fixed, and
a branch cannot be resolved by a state that precedes its origin. Since
$\prec$ follows streams and the line forest, merges are confined to
successors of the fork: later baselines of the same stream, of streams
forked after it, or of descendant lines. A defect already resolved in an
existing baseline is discarded with rationale \emph{already resolved in
$v$}, so that change logs only contain work actually incorporated.

\begin{hpbrule}[Gate integrity]
  \label{rule:gate}
  A planned baseline $g$ may be published at time $t$ only if every
  $b \in \mathcal{B}_t(g)$ is verified at $t$, and every baseline or branch
  on which a member depends is released, respectively merged, at $t$ or is
  published or merged in the same operation. If $g$ belongs to a launch
  bundle $\mathcal{G}$, the condition applies to the union of the targeted
  sets of $\mathcal{G}$, and all baselines of $\mathcal{G}$ are published
  together.
\end{hpbrule}
\noindent Every planned baseline becomes a decision point in the sense of
Stage-Gate~\citep{cooper1990stagegate}, applied to releases rather than
projects: committed work is delivered and verified, explicitly deferred or
explicitly abandoned, never left pending behind a shipped version and
never shipped ahead of what it depends on. The rule has content because
publication is the only way to merge; it is checked against the targeted
set \emph{before} publication.

\begin{hpbrule}[Lineage and life cycle]
  \label{rule:lineage}
  A hardware change that fails the interchangeability test creates a new
  line with an explicit parent link (\cref{def:linefork}); release streams
  of any kind do not. References are never edited in place, and parent
  links are never changed once created. New branches and baselines respect
  the life-cycle phase of their line (\cref{tab:phases}).
\end{hpbrule}
\noindent Editing a reference in place would rewrite the history of units
already in the field; the explicit link preserves the ancestry needed for
maintenance and audit (\cref{prop:lineage}), and the phase table keeps a
retired reference from silently accumulating work.

\begin{hpbrule}[Typed work]
  \label{rule:typed}
  Every branch has exactly one type $\tau(b)$. Sub-tasks are steps of the
  branch they belong to and are not branches.
\end{hpbrule}
\noindent Types determine lanes, colors and metrics; a branch with two
types, or a sub-task counted next to its parent, would be counted twice.
Task hierarchies belong to team boards, not to a product-level view.

\begin{hpbrule}[Traceable discussion]
  \label{rule:discussion}
  Every annotation in $A$ is attached by $\alpha$ to exactly one node of the
  model: a line, a baseline or a branch. Annotations and transitions are
  append-only.
\end{hpbrule}
\noindent Decisions and evidence that float in chat threads or meeting
notes cannot be audited; anchoring them to the node they concern puts the
rationale where the diagram is read, and append-only keeps the record
trustworthy.

\begin{remark}
  \label{rem:dag}
  The graph of a well-formed model is acyclic. Precedence edges follow
  $\prec$, which is a forest order; a fork edge followed by a merge edge
  goes from $f(b)$ to $m(b)$ with $f(b) \prec m(b)$ by \cref{rule:time};
  and no edge enters a branch except its fork edge, since sub-tasks are
  steps (\cref{rule:typed}). HPB thus keeps the structural property that
  makes version-control histories readable: every path runs forward in
  precedence. The relation $\mathrm{dep}$ must itself be acyclic, since a
  cycle of dependencies could never be released under \cref{rule:gate}.
\end{remark}

\subsection{Derived properties}
\label{sec:model-properties}

The rules are chosen so that release notes, quality attribution, flow
figures and audit answers are \emph{functions} of the graph rather than
separate records maintained by hand.

\begin{property}[Change log]
  \label{prop:changelog}
  For every baseline $v$, the set
  $C(v) = \{\, b \in \mathcal{B} : \mathrm{closes}(b, v) \,\}$ is the
  content of $v$. The non-empty sets among $\{C(v)\}_{v \in \mathcal{V}}$
  partition the set of merged branches, and $C(v)$ does not change after
  $v$ is released except through audited reopen transitions.
\end{property}
\begin{proof}[Proof sketch]
  By \cref{rule:closure}, each merged branch has exactly one merge baseline,
  hence belongs to exactly one $C(v)$. Merges happen only at publication
  (\cref{def:merge,def:gate}), so no ordinary transition adds a branch to
  $C(v)$ once $v$ is released. Computing all change logs is a single pass
  over $\mathcal{B}$; grouping $C(v)$ by $\tau$ yields release notes
  without manual compilation.
\end{proof}

\begin{property}[Escaped defects]
  \label{prop:escaped}
  For every baseline $v$, the set
  $E(v) = \{\, b \in \mathcal{B} : \tau(b) = \mathsf{X} \ \wedge\ f(b) = v
  \,\}$ contains the issues observed in $v$. The non-empty sets among
  $\{E(v)\}_{v \in \mathcal{V}}$ partition the set of issues, and every
  merged $b \in E(v)$ is resolved in a successor of $v$, i.e.,
  $v \prec m(b)$.
\end{property}
\begin{proof}[Proof sketch]
  By \cref{rule:origin,rule:typed}, each issue has exactly one type and
  exactly one fork baseline, hence belongs to exactly one $E(v)$. The
  second claim is \cref{rule:time} with $f(b) = v$.
\end{proof}
\noindent $E(v)$ attributes \emph{observed} quality to versions; the
baseline $\iota(b)$ in which a defect was \emph{introduced} may be earlier
and requires root-cause analysis. The property holds only if $f(b)$ is the
observed baseline, which is why \cref{rule:origin} fixes it at intake and
why an implementation must store it explicitly (\cref{sec:impl-datamodel};
HPB Markup stores it through \texttt{from}).

\begin{property}[Work in progress]
  \label{prop:wip}
  For every time $t$, the open set
  $W(t) = \{\, b \in \mathcal{B} : \sigma(b,t) \in \{\textsf{open},
  \textsf{targeted}\} \,\}$ is well defined and partitions by type into
  $W_{\mathsf{I}}(t), W_{\mathsf{X}}(t), W_{\mathsf{T}}(t),
  W_{\mathsf{L}}(t)$. Count every maximal interval during which a branch is
  open or targeted as one \emph{visit}. If, over a long observation window,
  the arrival rate $\lambda$ of visits and their mean duration
  $\overline{LT}$ have finite limits and the number of open branches does
  not grow without bound, then the time-average of $|W(t)|$ equals
  $\lambda \cdot \overline{LT}$.
\end{property}
\begin{proof}[Proof sketch]
  $\sigma$ is a total function on the branches opened by time $t$
  (\cref{def:state}), so $W(t)$ is determined, and \cref{rule:typed} makes
  the type partition exact. With visits as customers of a queueing system,
  $|W(t)|$ is the number of customers in the system and the visit length is
  the sojourn time. Little's law~\citep{little1961} then applies under the
  stated limit conditions; it requires no assumption on the distribution of
  arrivals or on the order of service. Merges and discards both count as
  departures, so discarded work must be included when visit lengths are
  averaged. (The arrival rate $\lambda$ here is unrelated to the lane map
  of \cref{eq:lanemap}.)
\end{proof}

\begin{property}[Lineage]
  \label{prop:lineage}
  For every line $P$, the ancestry $P, \mathrm{parent}(P),
  \mathrm{parent}^2(P), \dots$ is finite and is the unique path from a root
  of the line forest to $P$. For every baseline $v$, its predecessors
  $\{u : u \preceq v\}$ form a chain, so that the \emph{cumulative content}
  $\widehat{C}(v) = \bigcup_{u \preceq v} C(u)$ is well defined and follows
  the ancestry of $v$ through streams and lines.
\end{property}
\begin{proof}[Proof sketch]
  By \cref{rule:lineage}, each line has at most one parent, assigned when
  the line is created and never changed afterwards; since a parent must
  exist before its child, ancestry strictly decreases in creation order,
  cannot cycle, and terminates at a root of the forest. For the second
  claim, every baseline has at most one immediate predecessor
  (\cref{def:mainline}), and $\prec$ is well founded because each
  immediate predecessor was created earlier; following immediate
  predecessors from $v$ therefore enumerates $\{u : u \prec v\}$ as a
  single finite chain.
\end{proof}

\noindent As a corollary, a fix merged into a baseline of a maintenance
stream, or into a later baseline of a parent line after a line fork, is
\emph{not} in $\widehat{C}(w)$ for any baseline $w$ of the main stream or
of the child line that does not descend from it. By \cref{rule:closure},
the decision to carry it over is recorded as a second branch linked by
$\mathrm{prop}$, or as an explicit decision on the first one
(\cref{fig:backport}).

\subsection{Ordinal layout}
\label{app:formal-ordinal}

Let $\mathcal{E}_{\le}$ be the finite set of \emph{historical events}
drawn in a diagram, i.e., those with a time $t(e) \le t_{\mathrm{now}}$:
released baselines, branch starts, steps and terminals, with times taken
from the model. For a discarded branch the terminal time is $t_c(b)$. For
a merged branch it is its verification time $t_r(b)$ (\cref{def:state}),
which does not exceed $t_c(b) = t(m(b))$: the check mark is drawn when the
resolution has been verified, and the rise into the baseline shows the
merge at publication. Let $\mathcal{E}_{>}$ be the planned baselines.

\begin{definition}[Ordinal layout]
  \label{def:ordinal}
  An \emph{ordinal layout} assigns a \emph{column} $c(e)$ to every event.
  The historical events receive the columns $1, \dots, |\mathcal{E}_{\le}|$,
  injectively and such that $t(e) < t(e')$ implies $c(e) < c(e')$; events
  with equal times are ordered by a fixed tie-breaking rule (terminals
  before the baseline that closes them, baselines before all other branch
  events, then by identifier). The NOW marker receives the next column,
  $c_{\mathrm{now}} = |\mathcal{E}_{\le}| + 1$, and the planned baselines
  of $\mathcal{E}_{>}$ follow in order of target date. Spacer columns that
  carry no event may be inserted (rule~\ref{lay:tags}). Column $j$ is drawn
  at horizontal position $x = j \cdot w$ for a constant column width $w$.
  Calendar time is shown only as tick labels on a time axis below the
  diagram, at the column of the first event of a calendar period, written
  ``Mon YYYY'' at the first tick and after every change of year.
\end{definition}

A \emph{linear} layout instead uses $x(e) = (t(e) - t_{\min}) / \Delta$
for a fixed time scale $\Delta$. It preserves durations but requires
collision handling for tags and labels and cannot guarantee a minimum
spacing between events (\cref{sec:notation-time}).

% =====================================================================
\subsection{Metric definitions}
\label{app:metrics}

This subsection defines the metrics of \cref{tab:metrics}. $\mathcal{B}$
is the set of branches of a product line (or of several lines, when
stated) and $\mathcal{V}$ its set of baselines. Closed branches split into
merged, $\mathcal{B}^{\mathrm{m}}$, and discarded,
$\mathcal{B}^{\mathrm{d}}$; $\mathcal{B}_\theta$ is the set of branches of
type $\theta$. The closing time $t_c(b)$ is the publication time of $m(b)$
for a merged branch and the time of the discard otherwise
(\cref{def:state}). Windows are intervals $(t_1,t_2]$ of length
$\Delta=t_2-t_1$, and a subscript $(t_1,t_2]$ keeps the branches with
$t_c(b)$ in the window. Lead times are measured per \emph{visit} (each
period during which a branch is open or targeted).

\paragraph{Flow.}
\begin{gather}
  \mathrm{LT}(b) = t_c(b) - t_o(b), \quad b \in \mathcal{B}^{\mathrm{m}} \cup \mathcal{B}^{\mathrm{d}};
  \label{eq:leadtime}\\
  \mathrm{Age}(b,t) = t - t_o(b), \quad b \in W(t);
  \label{eq:age}\\
  W(t) = \{\, b \in \mathcal{B} : \sigma(b,t) \in \{\textsf{open}, \textsf{targeted}\} \,\},
  \quad W_\theta(t) = W(t) \cap \mathcal{B}_\theta;
  \label{eq:wip}\\
  \mathrm{TH}(t_1,t_2) = \bigl|\{\, b \in \mathcal{B}^{\mathrm{m}} : t_1 < t_c(b) \le t_2 \,\}\bigr| \,/\, \Delta .
  \label{eq:throughput}
\end{gather}
Little's law \citep{little1961} states that, in the long run and provided
the averages exist,
\begin{equation}
  \overline{|W|} = \lambda \cdot \overline{\mathrm{LT}},
  \label{eq:little}
\end{equation}
where $\overline{|W|}$ is the time average of $|W(t)|$, $\lambda$ the rate
at which visits start (equal, in a stable system, to the rate at which
they end) and $\overline{\mathrm{LT}}$ the mean lead time per visit. It
makes no assumption about the distribution of arrivals, but it needs a
window long enough and no growing population of branches that never close
(the \emph{eternal branches} of \cref{subsec:method-antipatterns}).

\paragraph{Closure.}
\begin{gather}
  \mathrm{MR} = \frac{|\mathcal{B}^{\mathrm{m}}_{(t_1,t_2]}|}
    {|\mathcal{B}^{\mathrm{m}}_{(t_1,t_2]}| + |\mathcal{B}^{\mathrm{d}}_{(t_1,t_2]}|},
  \qquad \mathrm{DR} = 1-\mathrm{MR};
  \label{eq:mergeratio}\\
  D(b) = \bigl|\{\, v \in \mathcal{V} : s(v)=\textsf{released},\ f(b) \prec v \preceq m(b) \,\}\bigr| ,
  \label{eq:mergedistance}
\end{gather}
where $\prec$ is the precedence of \cref{def:mainline}, crossing the
lineage connector when a branch merges into a child line. Discards
\emph{resolved without a product baseline} are reported apart from
abandoned work.

\paragraph{Releases and quality.}
\begin{gather}
  |C(v)| = \bigl|\{\, b : \mathrm{closes}(b,v) \,\}\bigr| ;
  \label{eq:batch}\\
  E(v) = \{\, b \in \mathcal{B}_{\mathsf{X}} : f(b) = v \,\}, \qquad
  \epsilon(v) = \frac{|E(v)|}{\max(1,|C(v)|)} ;
  \label{eq:escaped}\\
  \eta(v;T) = \frac{\bigl|\{\, b \in E(v) : t_o(b) \le t(v) + T \,\}\bigr|}
    {\int_{t(v)}^{t(v)+T} N(v,t)\,\mathrm{d}t} ;
  \label{eq:exposure}\\
  \mathrm{GS}(g) = t(g) - \hat t(g) .
  \label{eq:gateslip}
\end{gather}
$C(v)$ is the release-note list of $v$. $E(v)$ counts issues \emph{seen}
in $v$; when root-cause analysis finds the baseline $\iota(b)$ that
\emph{introduced} a defect, count by $\iota$ as well. $N(v,t)$ is the
installed base (units running $v$ at time $t$), so $\eta$ is the number of
issues seen within a period $T$ after release per unit-time of operation;
without installed-base data, counting the issues opened within a fixed
period after $t(v)$ is a weaker substitute. In the gate slip, $\hat t(g)$
is the target date of $g$ when its content was first committed; it is
reported with the number of branches moved out of $g$ and discarded from
it.

\paragraph{Maintenance.}
\begin{gather}
  \mathrm{ML}(P;t_1,t_2) =
  \frac{\bigl|\{\, b : f(b)\in \mathcal{V}(P),\ t_1 < t_o(b) \le t_2 \,\}\bigr|}
       {\bigl|\{\, b : f(b)\in \mathcal{V}(P)\cup\mathcal{V}(\mathrm{desc}(P)),\ t_1 < t_o(b) \le t_2 \,\}\bigr|} ;
  \label{eq:maintload}\\
  U = \{\, b \in \mathcal{B}^{\mathrm{m}} \cap \mathcal{B}_{\mathsf{X}} :
    \exists\, v \in \mathrm{aff}(b),\ v \npreceq m(b),\
    \nexists\, b' : (b,b') \in \mathrm{prop} \,\} ,
  \label{eq:unprop}
\end{gather}
where $\mathcal{V}(P)$ is the set of baselines of line $P$,
$\mathrm{desc}(P)$ its descendant lines, $\mathrm{aff}(b)$ the baselines
known to be affected and $\mathrm{prop}$ the ``is a port of'' relation
(\cref{def:relations}); branches with a recorded decision not to port are
removed from $U$. The same maintenance ratio works for a maintenance
stream relative to the main stream of its line.

% =====================================================================
\subsection{Layout algorithm}
\label{app:layout}
\label{sec:impl-layout}

Two layout rules fix the geometry of a diagram completely: ordinal columns
(\ref{lay:ordinal}) fix the horizontal position of every event, and ``one
branch per lane at a time'' (\ref{lay:onelane}) constrains the vertical
position of every branch. \Cref{alg:layout} computes both for one product
line. $f(b)$, $m(b)$ and $g(b)$ are the fork, merge and target baselines
of branch $b$, $\tau(b)$ its type, and $c(\cdot)$ the column of an event
(\cref{def:ordinal}).

Each branch occupies a half-open column interval, from its fork baseline
to its end (merge baseline, discard terminal, target baseline or NOW), so a
branch that merges into a baseline and one that forks from it can share a
lane. Assigning lanes is then the classic \emph{interval partitioning}
problem~\citep{kleinberg2006algo}, solved separately for each type. The
greedy solution uses exactly as many lanes as the maximum number of
branches of that type open at the same moment, which is the minimum
possible. Among the free lanes, a branch already closed by NOW takes the
innermost one and a branch still open or targeted the outermost one, so
long-lived work moves outward and later forks reach their inner lanes
without crossing it (\cref{fig:scale}). The HPB Markup renderer~1.0 packs
lanes first-fit per type and does not yet apply this preference.

\begin{algorithm}[tbp]
  \caption{HPB layout of one product line: ordinal columns and per-type
    lane partitioning.}
  \label{alg:layout}
  \small
  \begin{algorithmic}[1]
    \Require line $P$; baselines $\mathcal{V}(P)$; branches $\mathcal{B}$
      of $P$; present time $t_{\mathrm{now}}$
    \Ensure column $c(e)$ of every event; lane $y(b)$ of every branch
    \State $\mathcal{E} \gets$ released baselines, branch starts, steps and
      terminals with time $\le t_{\mathrm{now}}$
    \State sort $\mathcal{E}$ by (time, tie-breaking rank, identifier)
    \State $c(e_j) \gets j$ for the $j$-th event $e_j$ of $\mathcal{E}$;
      \quad $c_{\mathrm{now}} \gets |\mathcal{E}| + 1$
    \State assign columns $c_{\mathrm{now}}+1, c_{\mathrm{now}}+2, \dots$ to
      the planned baselines in order of target date
    \ForAll{$b \in \mathcal{B}$} \Comment{lane interval $[s_b, e_b)$}
      \State $s_b \gets c(f(b))$
      \State $e_b \gets c(m(b))$ if merged; $c(\text{terminal})$ if
        discarded; $c(g(b))$ if targeted; $c_{\mathrm{now}}$ otherwise
    \EndFor
    \ForAll{$\theta \in \{\mathsf{I}, \mathsf{L}, \mathsf{X}, \mathsf{T}\}$}
      \State $Q \gets$ branches of type $\theta$ sorted by $(s_b, e_b)$
      \State $k_\theta \gets$ maximum number of intervals of $Q$ sharing a
        point \Comment{sweep over sorted endpoints}
      \State $F \gets \{1, \dots, k_\theta\}$ \Comment{free lanes,
        ordered set}
      \State $H \gets \emptyset$ \Comment{busy lanes, min-heap on $e_b$}
      \ForAll{$b \in Q$ in order}
        \While{$H \neq \emptyset$ \textbf{and} $\min_{H} e \le s_b$}
          \State move the lane of the minimum of $H$ back to $F$
        \EndWhile
        \If{$b$ is open or targeted at $t_{\mathrm{now}}$}
          \State $\ell_b \gets \max F$ \Comment{outermost free lane}
        \Else
          \State $\ell_b \gets \min F$ \Comment{innermost free lane}
        \EndIf
        \State remove $\ell_b$ from $F$
        \State insert $(e_b, \ell_b)$ into $H$
      \EndFor
    \EndFor
    \ForAll{$b \in \mathcal{B}$} \Comment{stack lanes outward from the
      main line}
      \State $y(b) \gets
        \begin{cases}
          +\ell_b & \tau(b) = \mathsf{I}\\
          +(k_{\mathsf{I}} + \ell_b) & \tau(b) = \mathsf{L}\\
          -\ell_b & \tau(b) = \mathsf{X}\\
          -(k_{\mathsf{X}} + \ell_b) & \tau(b) = \mathsf{T}
        \end{cases}$
    \EndFor
  \end{algorithmic}
\end{algorithm}

With $n$ events and branches, the layout runs in $O(n \log n)$ time and
$O(n)$ space; drawing glyphs, tags and labels afterwards is linear. Every
sort ends with the identifier, so the same model at the same present time
always gives the same diagram. Two variants keep long histories readable:
\emph{collapsed history} hides the merged branches of all but the most
recent released baselines and replaces them by count badges
(\cref{fig:scale}), and \emph{type filters} remove the branches, and the
lanes, of deselected types. An optional \emph{compaction} that lets
consecutive events of one kind within a day, week or month share a column
bends rule~\ref{lay:ordinal}, and a view that uses it must say so.

% =====================================================================
\subsection{Prototype implementation: data model and status}
\label{app:impl}
\label{sec:impl-datamodel}

\Cref{fig:datamodel} shows the overlay schema of the prototype described
in \cref{sec:implementation}. The existing trackers stay as they are.
\texttt{products} holds the product lines, and \texttt{parent\_product\_id}
the family tree of references. \texttt{product\_versions} holds the
baselines with kind, status and date; a hardware baseline that creates a
new reference stores it in \texttt{result\_product\_id}, which is the line
fork of \cref{rule:lineage}. \texttt{product\_plan\_items} attaches one
tracker item, by (\texttt{item\_type}, \texttt{item\_id}), to one line:
the pair gives the branch type, \texttt{version\_id} holds the target
while the branch is open and the merge baseline once it is merged, and
\texttt{outcome}, \texttt{closed\_at} and \texttt{note} record the terminal
and the rationale of a discard. \texttt{product\_plan\_comments} attaches
each comment to exactly one node, so \cref{rule:discussion} is enforced by
the schema.

% F9 -- Data model of the reference implementation: an overlay on existing
% trackers. Requires hpb.sty v1.0 (palette, styles, layers of hpbgraph).
% Units: hpbgraph with col = lane = 1cm, i.e. plain centimetres.
\begin{figure}[t]
\centering
\resizebox{\linewidth}{!}{%
\begin{hpbgraph}[col=1cm, lane=1cm]
  \tikzset{
    dm/head/.style  = {anchor=north west, inner xsep=4pt, inner ysep=3pt,
                       align=left, text=hpbInk,
                       font=\sffamily\fontsize{7}{8.6}\selectfont},
    dm/row/.style   = {anchor=north west, inner xsep=4pt, inner ysep=0pt,
                       minimum height=0.3cm, text height=5pt, text depth=1.6pt,
                       align=left, text=hpbInk,
                       font=\fontsize{6.5}{7.8}\selectfont},
    dm/frame/.style = {draw=hpbInk, line width=1pt, rounded corners=4pt},
    dm/fk/.style    = {draw=hpbInk, line width=0.9pt, rounded corners=5pt,
                       -{Stealth[length=5pt,width=4.5pt]}},
    dm/ref/.style   = {dm/fk, dash pattern=on 4pt off 2.5pt},
    dm/lbl/.style   = {font=\sffamily\fontsize{6.5}{7.8}\selectfont, text=hpbInk,
                       fill=white, inner sep=1.2pt, align=center},
    dm/mult/.style  = {font=\sffamily\bfseries\fontsize{6.5}{7.8}\selectfont,
                       text=hpbMuted, inner sep=1pt},
  }
  % ---- entity macro (local to this figure) --------------------------------
  % \dment{<id>}{<x,y top-left>}{<width cm>}{<fill>}{<table>}{<model element>}
  %       {<key/field, key/field, ...>}   key in {PK, FK, (empty)}
  % Defines nodes <id>-h (header), <id>-1, <id>-2, ... (one per field) and
  % the corners <id>-nw, <id>-se.
  \def\dmPK{PK}%
  \newcommand\dmkey[1]{\edef\dmk{#1}\makebox[0.42cm][l]{\sffamily\bfseries
    \ifx\dmk\dmPK\dmk\else\textcolor{hpbMuted}{\dmk}\fi}}
  \newcommand\dment[7]{%
    \pgfmathsetlengthmacro\dmw{#3cm-8pt}%
    \node[dm/head, text width=\dmw] (#1-h) at (#2)
      {\textbf{\ttfamily\fontsize{7.5}{9}\selectfont #5}\\[-0.5pt]{\itshape #6}};
    \foreach \k/\f [count=\i] in {#7}{%
      \node[dm/row, text width=\dmw] (#1-\i)
        at ([yshift={-0.08cm-(\i-1)*0.3cm}]#1-h.south west)
        {\dmkey{\k}\texttt{\f}};
      \xdef\dmlast{\i}}%
    \coordinate (#1-se) at ([yshift=-0.1cm]#1-\dmlast.south west -| #1-h.east);
    \coordinate (#1-nw) at (#1-h.north west);
    \begin{pgfonlayer}{background}
      \begin{scope}
        \clip[rounded corners=4pt] (#1-nw) rectangle (#1-se);
        \fill[white] (#1-nw) rectangle (#1-se);
        \fill[#4] (#1-nw) rectangle (#1-h.south east);
      \end{scope}
    \end{pgfonlayer}
    \draw[draw=hpbInk, line width=0.6pt] (#1-h.south west) -- (#1-h.south east);
    \draw[dm/frame] (#1-nw) rectangle (#1-se);}

  % ---- HPB tables -------------------------------------------------------------
  \dment{pl}{0,0}{3.7}{hpbMain}{products}{product line $P$}
    {PK/id, /model, /manufacturer\_ref, FK/parent\_product\_id}
  \dment{bl}{0,-4.25}{3.7}{hpbMain}{product\_versions}{baseline $v$}
    {PK/id, FK/product\_id, /name, /kind, /status,
     /target\_date, /released\_at, FK/result\_product\_id}
  \dment{pi}{5.9,-0.45}{3.9}{hpbHw!45}{product\_plan\_items}
    {state of branch $b$ on line $P$}
    {PK/id, FK/product\_id, /item\_type, /item\_id, FK/version\_id,
     /outcome, /closed\_at, /note}
  \dment{cm}{5.9,-4.95}{3.9}{hpbDisc}{product\_plan\_comments}
    {discussion on a node (R7)}
    {PK/id, FK/product\_id, /target\_type, /target\_id, /body,
     {/created\_by, created\_at}}

  % ---- existing trackers (referenced, not migrated) -------------------------
  \draw[draw=hpbMuted, line width=0.8pt, rounded corners=5pt,
        dash pattern=on 3pt off 2pt, fill=hpbGuide!25]
    (11.75,-0.1) rectangle (15.95,-5.55);
  \node[font=\sffamily\bfseries\fontsize{7}{8.4}\selectfont, text=hpbInk,
        anchor=north west] at (11.85,-0.15) {existing trackers};
  \node[font=\sffamily\itshape\fontsize{6.5}{7.8}\selectfont, text=hpbMuted,
        anchor=north west, align=left] at (11.85,-0.5)
    {branch $b$ with type $\tau(b)$};
  \foreach \t/\c/\l/\y in {improvements/hpbImp/I/-1.05,
                           issues/hpbIss/X/-2.15,
                           tasks/hpbTask/T/-3.25,
                           lab\_items/hpbLab/L/-4.35}{
    \begin{scope}
      \clip[rounded corners=4pt] (12.0,\y) rectangle (15.7,\y-0.9);
      \fill[white] (12.0,\y) rectangle (15.7,\y-0.9);
      \fill[\c] (12.0,\y) rectangle (15.7,\y-0.42);
    \end{scope}
    \draw[dm/frame] (12.0,\y) rectangle (15.7,\y-0.9);
    \draw[draw=hpbInk, line width=0.6pt] (12.0,\y-0.42) -- (15.7,\y-0.42);
    \node[hpb/event, fill=\c, minimum size=9.5pt, line width=1.1pt,
          font=\sffamily\bfseries\fontsize{6.5}{7.8}\selectfont]
      at (12.3,\y-0.21) {\l};
    \node[font=\ttfamily\bfseries\fontsize{7.5}{9}\selectfont, text=hpbInk,
          anchor=west, inner sep=0pt] at (12.62,\y-0.21) {\t};
    \node[font=\ttfamily\fontsize{6.5}{7.8}\selectfont, text=hpbInk,
          anchor=west, inner sep=0pt] at (12.12,\y-0.66)
      {id, reference, title, status};
  }
  \coordinate (tr-w)  at (11.75,-2.15);
  \coordinate (tr-s)  at (13.85,-5.55);

  % ---- foreign keys (solid): referencing -> referenced table ---------------
  % products.parent_product_id -> products.id   (line forest)
  \draw[dm/fk] (pl-4.west) -- ++(-0.35,0) |- (pl-1.west);
  \node[dm/lbl, rotate=90] at ([xshift=-0.35cm]pl-1.west |- 0,-1.42) {parent};
  % product_versions.result_product_id -> products   (line fork: new reference)
  \draw[dm/fk] (bl-8.west) -- ++(-0.75,0) |- (pl-h.west);
  \node[dm/lbl, rotate=90] at (-0.75,-4.6) {line fork (new reference)};
  \node[dm/mult, anchor=north east] at ([xshift=-1pt]bl-8.west) {*};
  \node[dm/mult, anchor=south east] at ([xshift=-2pt,yshift=2.5pt]pl-h.west) {0..1};
  % product_versions.product_id -> products   (baselines of a line)
  \draw[dm/fk] (1.85,-4.25) -- (1.85,0 |- pl-se);
  \node[dm/lbl, anchor=west] at (1.9,-3.15) {product\_id};
  \node[dm/mult, anchor=north west] at ([xshift=3pt,yshift=-1pt]1.85,0 |- pl-se) {1};
  \node[dm/mult, anchor=south west] at ([xshift=1pt]1.85,-4.25) {*};
  % product_plan_items.product_id -> products
  \draw[dm/fk] (pi-2.west) -- (4.75,0 |- pi-2.west) |- (pl-2.east);
  \node[dm/mult, anchor=south east] at ([xshift=-1pt]pi-2.west) {*};
  \node[dm/mult, anchor=south west] at ([xshift=2pt,yshift=2.5pt]pl-2.east) {1};
  % product_plan_items.version_id -> product_versions   (g(b) or m(b))
  \draw[dm/fk] (pi-5.west) -- (4.75,0 |- pi-5.west) |- (bl-3.east);
  \node[dm/mult, anchor=south east] at ([xshift=-1pt]pi-5.west) {*};
  \node[dm/mult, anchor=south west] at ([xshift=2pt,yshift=2.5pt]bl-3.east) {0..1};
  \node[dm/lbl, anchor=west, align=left] at (4.8,-4.4)
    {target $g(b)$ or\\closing $m(b)$};

  % ---- polymorphic references (dashed): (type, id) pairs ---------------------
  \draw[dm/ref] (pi-4.east) -- (tr-w |- pi-4.east);
  \node[dm/lbl, anchor=south] at ([yshift=1pt]10.78,0 |- pi-4.east)
    {item\_type,\\item\_id};
  \node[dm/lbl, anchor=north, font=\sffamily\itshape\fontsize{6.5}{7.8}\selectfont,
        text=hpbMuted] at ([yshift=-3pt]10.78,0 |- pi-4.east)
    {referenced,\\not migrated};
  % comments -> branch item, baseline (or product)
  \draw[dm/ref] (cm-4.east) -| (tr-s);
  \draw[dm/ref] (cm-nw |- bl-7.east) -- (bl-7.east);
  \node[dm/lbl, anchor=north west, align=left] at ([yshift=-2pt]cm-se -| cm-nw)
    {target\_type $\in$ \{product, version, improvement,\\
     issue, task, lab\}, identified by target\_id};
\end{hpbgraph}}
\caption{Data model of the prototype implementation. HPB is an
overlay: improvements, issues, tasks and lab items stay in their trackers and
are referenced, not migrated. Solid arrows are foreign keys; dashed arrows are
references by (type, identifier).}
\label{fig:datamodel}
\end{figure}


Some attributes are \emph{derived} today instead of stored: the opening
time of a branch is the creation time of its tracker item, and the fork
baseline is the latest release at that time. That default is right for
improvements, tasks and lab work, but wrong for issues, whose fork
baseline is the version in which the defect was observed
(\cref{prop:escaped}). The design therefore adds a stored
\texttt{fork\_version\_id}, mandatory for issues and with an
\emph{inferred} flag, and an append-only log of transitions with actor,
approver and time, which makes audited reopens possible and keeps the
original target date from which gate slip is measured. Streams,
dependencies, propagations, launch bundles, life-cycle milestones and the
deployment overlay need further tables.

\Cref{tab:impl-status} separates what the design specifies from what the
prototype does today. The prototype still accepts a discard without a
rationale, offers released baselines when an item is closed as merged,
merges the whole targeted set without checking verification, derives the
fork baseline instead of storing it, and packs lanes per side of the main
line instead of per type. It therefore does not yet enforce
\cref{rule:closure,rule:time,rule:gate}; until it does, those violations
are at least visible in the diagram.

\begin{table}[tbp]
  \centering
  \small
  \caption{Prototype in a product management portal: designed versus
    built. \checkmark: built; $\sim$: partly built; --: designed only.}
  \label{tab:impl-status}
  \begin{tabularx}{\linewidth}{@{}>{\raggedright\arraybackslash}X
      >{\centering\arraybackslash}p{1.4cm}
      >{\raggedright\arraybackslash}p{6.2cm}@{}}
    \toprule
    \textbf{Element of the design} & \textbf{Status} & \textbf{Remark} \\
    \midrule
    Overlay tables, references to tracker items & \checkmark & no migration of trackers \\
    Derived mode & \checkmark & read-only reconstruction \\
    Line family tree and line forks & \checkmark & \texttt{parent\_product\_id}, \texttt{result\_product\_id} \\
    Stored observed fork baseline, inferred flag & -- & needed for escaped defects per baseline \\
    Release streams, configurations, applies-to & -- & \\
    Dependencies, propagations, launch bundles & -- & \\
    Life-cycle milestones, deployment overlay & -- & \\
    Append-only transition log with approvers & -- & state tables only \\
    R1 single origin & $\sim$ & by derivation, not at intake \\
    R2 explicit closure & $\sim$ & rationale not yet mandatory \\
    R3 time moves forward & $\sim$ & released baselines still offered \\
    R4 gate integrity & $\sim$ & batch merge without verification \\
    R5 lineage and life cycle & \checkmark & life-cycle phases not checked \\
    R6 typed work & \checkmark & type from the source tracker \\
    R7 traceable discussion & $\sim$ & comments on nodes; not append-only \\
    Ordinal columns & \checkmark & \\
    Lane partitioning per type (\cref{alg:layout}) & -- & today: per side of the main line \\
    \bottomrule
  \end{tabularx}
\end{table}

% flush this appendix's floats before the next appendix
\clearpage

% ---------------------------------------------------------------------
%  appendix-a-card.tex -- HPB notation quick-reference card
%  Glyph column drawn with hpb.sty v1.0 (normative glyph sizes).
% ---------------------------------------------------------------------
\section{Notation Quick-Reference Card}
\label{app:card}

Print this card and keep it next to the diagram the first few times you
read one. \Cref{tab:card-types} gives the four branch types with their
lane, letter and color; \cref{tab:card,tab:card-branches} list every
glyph, what it means and the \texttt{hpb.sty} macro that draws it
(\cref{app:package}); the ten reading rules at the end are all you need to
read any HPB diagram. The notation itself is defined in
\cref{sec:notation}, the model behind it in \cref{sec:model} and
\cref{app:formal}.

\providecommand{\HPBmac}[1]{\texttt{\textbackslash #1}}
% \HPBglyph[<hpbgraph keys>]{<hpb macros>}: one glyph sample with 3pt of
% vertical padding (glyph sizes are normative; only spacing is reduced).
% Same definition as in sections/05-notation.tex.
\providecommand{\HPBglyph}[2][]{%
  \begin{hpbgraph}[col=0.55cm,lane=0.9cm,legend=1.3cm,#1]#2%
    \path ([yshift=3pt]current bounding box.north)
          ([yshift=-3pt]current bounding box.south);%
  \end{hpbgraph}}

\begin{table}[htbp]
  \centering
  \small
  \caption{Branch types: legend box, start node, lane, color and intent.}
  \label{tab:card-types}
  \renewcommand{\arraystretch}{1.15}
  \begin{tabular}{@{}>{\centering\arraybackslash}m{3.3cm}
      >{\centering\arraybackslash}m{0.8cm}
      >{\centering\arraybackslash}m{0.8cm}
      >{\raggedright\arraybackslash}m{2.6cm}
      >{\raggedright\arraybackslash}m{6.0cm}@{}}
    \toprule
    Legend box \& start node & Type & Lane & Color & Intent \\
    \midrule
    \HPBglyph[legend=1.9cm]{\hpblane{0}{imp}{Improvement (I)}{0}{0.6}
      \hpbstart{imp}{0.9}{0}{I}}
    & $\mathsf{I}$ & $+1$ & violet \texttt{\#C4B5FD}
    & Proactive: a request to make the product better. \\
    \HPBglyph[legend=1.9cm]{\hpblane{0}{lab}{Lab (L)}{0}{0.6}
      \hpbstart{lab}{0.9}{0}{L}}
    & $\mathsf{L}$ & $+2$ & teal \texttt{\#5EEAD4}
    & Proactive: laboratory evolution, evaluation or qualification. \\
    \HPBglyph[legend=1.9cm]{\hpblane{0}{iss}{Issue (X)}{0}{0.6}
      \hpbstart{iss}{0.9}{0}{X}}
    & $\mathsf{X}$ & $-1$ & coral \texttt{\#FB8A6E}
    & Reactive: a defect observed in the fork baseline. \\
    \HPBglyph[legend=1.9cm]{\hpblane{0}{task}{Task (T)}{0}{0.6}
      \hpbstart{task}{0.9}{0}{T}}
    & $\mathsf{T}$ & $-2$ & green \texttt{\#6EE7B7}
    & Reactive or internal: engineering work needed to keep the product
      operable. \\
    \bottomrule
  \end{tabular}
\end{table}

\begin{table}[htbp]
  \centering
  \small
  \caption{HPB quick-reference card, part~1: lines, baselines and time.
    Glyphs at normative size, drawn with the listed macro (arguments in
    \cref{tab:api}).}
  \label{tab:card}
  \renewcommand{\arraystretch}{1.0}
  \renewcommand{\tabularxcolumn}[1]{m{#1}}
  \begin{tabularx}{\linewidth}{@{}>{\centering\arraybackslash}m{3.3cm}
      >{\raggedright\arraybackslash}X
      >{\raggedright\arraybackslash}m{3.75cm}@{}}
    \toprule
    Glyph & Element and how to read it & Macro \\
    \midrule
    \HPBglyph{\hpbmainline{0}{2.2}\hpbmainline*{2.2}{4.2}}
    & \textbf{Main line}: the product; solid = history, dashed = plan.
    & \HPBmac{hpbmainline}, \HPBmac{hpbmainline*} \\
    \HPBglyph{\hpbmainline{0}{2}\hpborigin{0}{Product A}}
    & \textbf{Origin}: creation of the reference (not a baseline) and the
      product reference in the key column.
    & \HPBmac{hpborigin} \\
    \HPBglyph{\hpbmainline{-1.2}{1.2}\hpbrelease{0}{v1.1}}
    & \textbf{Released baseline}: immutable version with its name tag.
    & \HPBmac{hpbrelease} \\
    \HPBglyph{\hpbmainline*{-2.4}{2.6}\hpbplanned*{-1.6}{v1.2}{Jan}
      \hpbplanned{1.6}{v1.3}{Mar}}
    & \textbf{Planned baseline}: dashed node and tag, target date; a gate.
      ``!'' badge, red tag and date: \emph{overdue}, the target date has
      passed.
    & \HPBmac{hpbplanned}, \HPBmac{hpbplanned*} \\
    \HPBglyph{\hpbmainline{-1.8}{1.2}\hpbhardware{-0.7}{Rev B}}
    & \textbf{Hardware baseline}: diamond; with \emph{new ref.}, the first
      baseline of a new line (line fork).
    & \HPBmac{hpbhardware} \\
    \HPBglyph{\hpbmainline{-1.8}{1.2}\hpbprocess{-0.7}{PC-7}}
    & \textbf{Process baseline}: change to how the product is built (kind
      \textsf{other}).
    & \HPBmac{hpbprocess} \\
    \HPBglyph[lane=0.7cm]{\hpbmainline{-0.6}{3.4}\node[hpb/baseline] at (0,0) {};
      \hpbstream{0}{-1}{3.4}}
    & \textbf{Release stream}: maintenance stream leaving a released
      baseline, with its own band.
    & \HPBmac{hpbstream} \\
    \HPBglyph[lane=0.7cm]{\hpbmainline{-1.4}{1.4}\hpbmilestone{0}{end of sale}}
    & \textbf{Life-cycle milestone}: end of sale, end of production,
      security only.
    & \HPBmac{hpbmilestone} \\
    \HPBglyph[lane=0.7cm]{\hpbmainline{-1.2}{1.2}\node[hpb/baseline] at (0,0) {};
      \hpbcollapsed{0}{12}}
    & \textbf{Collapsed history}: number of merged branches not drawn.
    & \HPBmac{hpbcollapsed} \\
    \HPBglyph[lane=0.7cm]{\hpbmainline{-1.6}{0}\hpbeos{0}\path (1.2,0);}
    & \textbf{End of support}: the line is retired; no new branches.
    & \HPBmac{hpbeos} \\
    \HPBglyph[lane=0.7cm]{\hpbmainline{-1.6}{1.6}\hpbmainline{0}{1.6}[1.2]
      \hpblinefork{0}{0}{1.2}\node[hpb/hardware] at (0,1.2) {};}
    & \textbf{Line fork}: lineage connector from the parent main line to the
      hardware diamond, the first baseline of the child line $P'$
      (own band), with $\mathrm{parent}(P')=P$.
    & \HPBmac{hpblinefork} \\
    \midrule
    \HPBglyph[lane=0.7cm]{\hpbmainline{-1.6}{0}\hpbmainline*{0}{1.6}\hpbnow{0}{0}{0.2}}
    & \textbf{NOW}: history to the left, plan to the right.
    & \HPBmac{hpbnow} \\
    \HPBglyph[lane=0.7cm]{\hpbmainline*{-0.4}{1.6}[0.4]\hpbmainline*{-0.4}{1.6}[-0.6]
      \node[hpb/planned] at (0.6,0.4) {};\node[hpb/planned] at (0.6,-0.6) {};
      \hpbbundle{-0.2,-1.0}{1.4,0.8}{}}
    & \textbf{Launch bundle}: planned baselines of several lines published
      together or not at all.
    & \HPBmac{hpbbundle} \\
    \HPBglyph{\hpbaxis{0}{3.4}{0}\hpbtick{0.5}{Jan}{0}\hpbtick{2.3}{Apr}{0}}
    & \textbf{Ordinal time axis}: calendar periods as ticks only.
    & \HPBmac{hpbaxis}, \HPBmac{hpbtick} \\
    \bottomrule
  \end{tabularx}
\end{table}

\begin{table}[htbp]
  \centering
  \small
  \caption{HPB quick-reference card, part~2: branches, terminals, steps and markers
    (\HPBmac{hpbbranch} draws the start node, the path, the label and the
    terminal in one call).}
  \label{tab:card-branches}
  \renewcommand{\arraystretch}{1.0}
  \renewcommand{\tabularxcolumn}[1]{m{#1}}
  \begin{tabularx}{\linewidth}{@{}>{\centering\arraybackslash}m{3.3cm}
      >{\raggedright\arraybackslash}X
      >{\raggedright\arraybackslash}m{3.75cm}@{}}
    \toprule
    Glyph & Element and how to read it & Macro \\
    \midrule
    \HPBglyph[lane=0.75cm]{\hpbmainline{-0.6}{3.8}
      \node[hpb/baseline] at (0,0) {};
      \hpbbranch[label={X-31 Crash}]{iss}{0}{-1}{1.3}{3.8}{none}{}}
    & \textbf{Branch}: forks from a baseline and runs on its type lane;
      start node with type letter, label with reference and title.
    & \HPBmac{hpbbranch}, \HPBmac{hpbstart}, \HPBmac{hpblabel} \\
    \HPBglyph[lane=0.75cm]{\hpbmainline{-0.6}{4.6}
      \node[hpb/baseline] at (0,0) {};
      \node[hpb/baseline] at (4,0) {};
      \hpbbranch{imp}{0}{1}{1.2}{2.7}{merge}{4}}
    & \textbf{Merged}: check mark; the path rejoins the closing baseline.
    & \HPBmac{hpbbranch} (\texttt{merge}) \\
    \HPBglyph[lane=0.75cm]{\hpbmainline{-0.6}{3.8}
      \node[hpb/baseline] at (0,0) {};
      \hpbbranch{task}{0}{-1}{1.3}{3.3}{discard}{}}
    & \textbf{Discarded}: gray path and cross; read the rationale.
    & \HPBmac{hpbbranch} (\texttt{discard}) \\
    \HPBglyph[lane=0.75cm]{\hpbmainline{-0.6}{3.8}
      \node[hpb/baseline] at (0,0) {};
      \hpbbranch{lab}{0}{1}{1.3}{3.3}{open}{}
      \draw[hpb/now] (3.3,1.6) -- (3.3,0.4);}
    & \textbf{Open}: hollow head on the NOW marker.
    & \HPBmac{hpbbranch} (\texttt{open}) \\
    \HPBglyph[lane=0.75cm]{\hpbmainline{-0.6}{2.7}\hpbmainline*{2.7}{5}
      \node[hpb/baseline] at (0,0) {};
      \node[hpb/planned] at (4.4,0) {};
      \hpbbranch{iss}{0}{-1}{1}{2.7}{target}{4.4}}
    & \textbf{Targeted}: dashed continuation to the planned baseline.
    & \HPBmac{hpbbranch} (\texttt{target}) \\
    \HPBglyph[col=0.5cm,lane=0.75cm]{\hpbmainline*{-0.6}{4.2}
      \node[hpb/planned] at (2,0) {};
      \node[hpb/planned] at (3.6,0) {};
      \draw[hpb/path] (-2.2,-1) -- (0,-1);
      \hpbstep{iss}{-1.5}{-1}{done}
      \hpbterminal{target}{iss}{0}{-1}
      \hpbtarget[faded]{0}{-1}{2}
      \hpbtarget{0}{-1}{3.6}}
    & \textbf{Re-targeted}: a labeled done step records the decision; the
      former (nearer) target remains as a faded dashed continuation, the
      new target is drawn at full strength.
    & \HPBmac{hpbtarget} (\texttt{faded}), \HPBmac{hpbstep} \\
    \HPBglyph[col=0.6cm,lane=0.7cm]{\draw[hpb/path] (-0.4,0) -- (4.4,0);
      \hpbstep{imp}{0}{0}{done}\hpbstep{imp}{1.33}{0}{todo}
      \hpbstep{lab}{2.67}{0}{pass}\hpbstep{lab}{4}{0}{fail}}
    & \textbf{Steps}: sub-task done, pending; test or verification passed,
      failed. A pass citing a candidate of the target verifies the branch.
    & \HPBmac{hpbstep} \\
    \HPBglyph[col=0.6cm,lane=0.7cm]{\draw[hpb/path] (-0.6,0) -- (0.5,0)
      -- (0.5,0.6);\hpbterminal{merge}{iss}{0}{0}
      \hpbpoststep{iss}{0.5}{3}{0}{pass}}
    & \textbf{Post-closure step}: e.g., effectiveness of a corrective action,
      on a dotted continuation after the terminal.
    & \HPBmac{hpbpoststep} \\
    \midrule
    \HPBglyph[lane=0.75cm]{\hpbmainline*{1.3}{4.6}
      \node[hpb/planned] at (4,0) {};
      \hpbgate[0.8]{4}{0}{-1}{}
      \hpbterminal{target}{iss}{1.5}{-1}
      \hpbtarget{1.5}{-1}{4}}
    & \textbf{Gate}: every dashed continuation entering the band is a
      targeted branch; the baseline cannot be published while one of them
      is unverified (\cref{rule:gate}).
    & \HPBmac{hpbgate} \\
    \HPBglyph[col=0.6cm,lane=0.7cm]{\hpbnode{imp}{0}{0}\hpbnode{iss}{4}{0}
      \hpbdepends{0,0}{4,0}{}\hpbnode{iss}{0}{-0.9}\hpbnode{iss}{4}{-0.9}
      \hpbpropagates{0,-0.9}{4,-0.9}{}}
    & \textbf{Dependency} (dotted, open head) and \textbf{propagation}
      (dashed, filled head) between branches or baselines.
    & \HPBmac{hpbdepends}, \HPBmac{hpbpropagates} \\
    \bottomrule
  \end{tabularx}
\end{table}

\subsection*{Reading rules}

\begin{enumerate}[label=\textbf{Q\arabic*},leftmargin=*,itemsep=2pt,topsep=3pt]
  \item \textbf{Find the product.} The thick horizontal line is the product
        line; its tags are its versions, oldest at the left.
  \item \textbf{Read the sides.} Above the line: how the product gets
        better (improvements nearest, laboratory work further out). Below
        the line: what is wrong with it and how it is kept correct (issues
        nearest, tasks further out).
  \item \textbf{Read the NOW marker.} Left of NOW is recorded history and
        does not change; right of NOW is the plan and is expected to change.
  \item \textbf{What does a version contain?} Follow every path that
        rejoins its node: these are the branches it closed, $C(v)$
        (\cref{prop:changelog}).
  \item \textbf{Where was a defect observed?} Follow an $\mathsf{X}$
        branch back to the baseline it forks from: that baseline is the
        version in which the issue was observed, $E(v)$
        (\cref{prop:escaped}).
  \item \textbf{What is in progress?} Count the heads on the NOW marker:
        hollow heads are open work; heads with a dashed continuation are
        committed to the planned baseline at its end.
  \item \textbf{Can the next version ship?} Look at its gate band: every
        dashed continuation that enters it belongs to a targeted branch,
        and any of them without a passed verification step blocks the
        release (\cref{rule:gate}).
  \item \textbf{What was abandoned?} Gray paths ending in a cross were
        closed without a version; their labels or comments give the
        rationale.
  \item \textbf{Where did a reference come from?} A diamond with
        \emph{new ref.}\ starts a new product reference; follow the
        vertical connector down to the parent line. Ticks mark life-cycle
        milestones, and a bar labeled \textsf{EoS} ends a retired line.
        Each line or stream is read relative to its own main line.
  \item \textbf{Do not measure lengths.} The time axis is ordinal: one
        column per event. Durations are computed, not read off the drawing
        (\cref{rem:ordinal}).
\end{enumerate}

% keep the whole card together, before the next appendix
\clearpage

% ---------------------------------------------------------------------
%  appendix-c-markup.tex -- HPB Markup v1.0 language reference
%  Kept in line with docs/language.md in the repository (sections 3, 5 and 10).
% ---------------------------------------------------------------------
\section{HPB Markup Language Reference}
\label{app:markup}

This appendix is the reference for HPB Markup v1.0, the text language of
\cref{sec:tooling}. The toolkit documentation at \hpbrepo{} carries the
same reference with more examples. \Cref{tab:markup} lists every
statement, \cref{tab:markup-diag} the main diagnostics, and the grammar
closes the appendix.

\paragraph{Lexical conventions.}
One statement per line. Keywords are case-insensitive. Identifiers match
\texttt{[A-Za-z0-9.\_-]+}, so versions such as \texttt{v1.0.1} and
references such as \texttt{Rev-B} or \texttt{X-118} need no quoting. Titles
and rationales are strings in double quotes; \texttt{\textbackslash"} and
\texttt{\textbackslash\textbackslash} escape a quote and a backslash.
\texttt{\#} and \texttt{\%\%} start a comment that runs to the end of the
line. Blank lines and indentation are ignored. Dates are written
\texttt{@YYYY-MM-DD} or \texttt{@YYYY-MM}.

\paragraph{Time.}
Statement order is time order: every statement that creates an event takes
the next column. The merge terminals of the branches closed by one
\texttt{release \dots{} merges} statement share one column just before the
baseline (they lie on different lanes; see layout rule~\ref{lay:ordinal}).
When \texttt{now} is omitted, NOW is drawn after the last historical event
only if the diagram has something in the future (planned baselines, open
branches or merges awaiting release); a pure history has no NOW marker.
Dates only produce calendar ticks; a date earlier than a previous one
produces a warning.

\paragraph{The gate rule in the Markup.}
The Markup checks a structural version of \cref{rule:gate}: no targeted
branch may still be open when its baseline is released. Verification
evidence is not modeled. A \texttt{release \dots{} merges} statement is the
publication of \cref{def:gate}, with the listed branches as its targeted
set. A \texttt{force} release records a deliberate violation: the model is
then not well formed under \cref{rule:gate}, and the diagram shows it with
the warning \texttt{W-R4-FORCED} and a ``!'' badge.

\providecommand{\HPBkw}[1]{\texttt{#1}}
\providecommand{\HPBph}[1]{\ensuremath{\langle}\textit{#1}\ensuremath{\rangle}}

\begin{table}[p]
  \centering
  \footnotesize
  \caption{HPB Markup v1.0 statements. Square brackets mark optional parts;
    the optional clauses of a statement may be written in any order, each
    at most once. \HPBph{vid} is a baseline (version) identifier,
    \HPBph{bid} a branch identifier, \HPBph{id} a line identifier.}
  \label{tab:markup}
  \renewcommand{\arraystretch}{1.12}
  \begin{tabularx}{\linewidth}{@{}>{\raggedright\arraybackslash}p{5.6cm}
      >{\raggedright\arraybackslash}X@{}}
    \toprule
    \textbf{Statement} & \textbf{Meaning} \\
    \midrule
    \multicolumn{2}{@{}l}{\textit{Document and lines}} \\
    \HPBkw{hpb} [\HPBph{string}]
      & Optional header with the diagram title. \\
    \HPBkw{config} \HPBph{key} [\HPBkw{=}] \HPBph{value}
      & \HPBkw{time} = \HPBkw{ordinal}\,|\,\HPBkw{linear};
        \HPBkw{labels} = \HPBkw{full}\,|\,\HPBkw{ref}\,|\,\HPBkw{none};
        \HPBkw{theme} = \HPBkw{default}\,|\,\HPBkw{mono}\,|\,\HPBkw{dark};
        \HPBkw{showClosed} = \HPBkw{true}\,|\,\HPBkw{false}. \\
    \HPBkw{line} \HPBph{id} [\HPBph{string}] [\HPBkw{from} \HPBph{vid}]
      & Declares a product line and makes it current; the first line is the
        default. With \HPBkw{from}, a child line forked from the released
        \emph{hardware} baseline \HPBph{vid} (line fork, \cref{rule:lineage}). \\
    \HPBkw{use} \HPBph{id}
      & Makes another declared line current. \\
    \HPBkw{eos} [\HPBph{id}] [\HPBph{date}]
      & End-of-support bar on a line. \\
    \midrule
    \multicolumn{2}{@{}l}{\textit{Baselines}} \\
    \HPBkw{release} \HPBph{vid} [\HPBph{kind}] [\HPBph{date}]
      [\HPBkw{merges} \HPBph{bid}, \dots] [\HPBkw{force}] [\HPBkw{on} \HPBph{id}]
      & Released baseline; publishes it and closes the listed branches
        into it. \HPBph{kind} is \HPBkw{firmware} (default),
        \HPBkw{software}, \HPBkw{hardware} (diamond) or \HPBkw{other}.
        \HPBkw{force} records a release with open targeted branches as a
        violation (warning badge). \\
    \HPBkw{plan} \HPBph{vid} [\HPBph{kind}] [\HPBph{date}]
      [\HPBkw{targets} \HPBph{bid}, \dots] [\HPBkw{on} \HPBph{id}]
      & Planned baseline, right of NOW; the listed branches are targeted at
        it. \\
    \HPBkw{collapse} \HPBph{vid} \HPBph{n}
      & Count badge: $n$ merged branches hidden at this baseline. \\
    \midrule
    \multicolumn{2}{@{}l}{\textit{Branches}} \\
    \HPBph{type} \HPBph{bid} [\HPBph{string}] [\HPBkw{from} \HPBph{vid}]
      [\HPBkw{on} \HPBph{id}] [\HPBkw{targets}\,|\,\HPBkw{target} \HPBph{vid}]
      [\HPBph{date}]
      & Opens a branch. \HPBph{type} is \HPBkw{improvement}
        (\HPBkw{imp}, \HPBkw{feature}), \HPBkw{issue} (\HPBkw{bug},
        \HPBkw{defect}, \HPBkw{x}), \HPBkw{task} (\HPBkw{t}) or \HPBkw{lab}
        (\HPBkw{l}, \HPBkw{evolution}). Without \HPBkw{from}, it forks from
        the latest released baseline of its line, the default of
        \cref{def:branch}; an issue observed in an older version needs
        \HPBkw{from} (\cref{rule:origin}). \\
    \HPBkw{step} \HPBph{bid} \HPBkw{done}\,|\,\HPBkw{todo}\,|\,%
      \HPBkw{pass}\,|\,\HPBkw{fail} [\HPBph{string}] [\HPBph{date}]
      & Intermediate step on a branch (e.g., a laboratory test). \\
    \HPBkw{merge} \HPBph{bid} \HPBkw{into} \HPBph{vid} [\HPBph{date}]
      & \HPBph{vid} must still be planned. Marks the branch as resolved,
        with its terminal at this column (the verification time of
        \cref{app:formal-ordinal}); the branch is merged when \HPBph{vid}
        is released. A baseline already released cannot receive later
        merges (\texttt{E-MERGE-PAST}). \\
    \HPBkw{discard} \HPBph{bid} [\HPBph{string}] [\HPBph{date}]
      & Discard terminal; the rationale is printed in italics. \\
    \HPBkw{target} \HPBph{bid} \HPBph{vid}
      & (Re)targets an open branch at a planned baseline; a former target
        is drawn faded. \\
    \midrule
    \multicolumn{2}{@{}l}{\textit{Markers}} \\
    \HPBkw{now}
      & The NOW column, at most once. Only planning statements may follow:
        \HPBkw{plan}, \HPBkw{target}, \HPBkw{note}, \HPBkw{collapse},
        \HPBkw{config} and \HPBkw{use}. \\
    \HPBkw{note} \HPBph{vid}\,|\,\HPBph{bid} \HPBph{string}
      & Callout attached to a node. \\
    \bottomrule
  \end{tabularx}
\end{table}

\paragraph{Diagnostics.}
Every diagnostic carries a severity, a stable code, a line, a column and a
message (\cref{tab:markup-diag}). The core checks rules R1--R4 directly and
reports them with the rule as code; R5 and R6 are protected by the
grammar and by the line-fork check. Errors stop the diagram from being
drawn as valid; by default the renderer shows them in an error box instead
of throwing. The full list of codes is in the toolkit documentation.

\begin{table}[!htbp]
  \centering
  \small
  \caption{Main diagnostics of HPB Markup v1.0, their codes and the rule
    each one protects.}
  \label{tab:markup-diag}
  \renewcommand{\arraystretch}{1.1}
  \begin{tabularx}{\linewidth}{@{}>{\raggedright\arraybackslash}X
      >{\raggedright\arraybackslash}p{3.3cm}
      >{\centering\arraybackslash}p{1.1cm}
      >{\centering\arraybackslash}p{1.0cm}@{}}
    \toprule
    \textbf{Situation} & \textbf{Code} & \textbf{Level} & \textbf{Rule} \\
    \midrule
    Branch with no released baseline to fork from & \texttt{R1} & error & R1 \\
    Branch closed twice, or changed after its terminal & \texttt{R2} & error & R2 \\
    Merge into a baseline that does not succeed the fork, including the fork itself & \texttt{R3} & error & R3 \\
    Release with open targeted branches (no \texttt{force}) & \texttt{R4} & error & R4 \\
    Release with \texttt{force} & \texttt{W-R4-FORCED} & warning & R4 \\
    \texttt{line \dots{} from} a baseline that is unknown, not released or not \texttt{hardware} & \texttt{E-LINE-FORK} & error & R5 \\
    Unknown statement, branch type or step state & \texttt{E-SYNTAX} & error & R6 \\
    Merge into a baseline released earlier & \texttt{E-MERGE-PAST} & error & -- \\
    Undeclared line, branch or baseline & \texttt{E-UNKNOWN-*} (\texttt{LINE}, \texttt{BRANCH}, \texttt{BASELINE}) & error & -- \\
    Historical event after \texttt{now} & \texttt{E-AFTER-NOW} & error & -- \\
    Date earlier than a previous date & \texttt{W-DATE-ORDER} & warning & -- \\
    \bottomrule
  \end{tabularx}
\end{table}

\paragraph{Grammar.}
The grammar below uses an EBNF-like notation: \texttt{?} optional,
\texttt{|} alternative, \texttt{*} repetition, literals in double quotes,
and \texttt{\{ \dots{} \}} any of the listed clauses, in any order, each at
most once.

\begin{lstlisting}[language=hpbplain,basicstyle=\ttfamily\scriptsize]
document   := header? statement*
header     := "hpb" string?
statement  := config | line | use | release | plan | branch | step
            | merge | discard | target | now | eos | collapse | note

config     := "config" key "="? value
line       := "line" id string? ("from" vid)?
use        := "use" id
release    := "release" vid { kind | date | "merges" idlist | "force" | "on" id }
plan       := "plan" vid { kind | date | "targets" idlist | "on" id }
kind       := "firmware" | "software" | "hardware" | "other"
branch     := btype bid { string | "from" vid | "on" id
                        | ("targets" | "target") vid | date }
btype      := "improvement" | "imp" | "feature" | "issue" | "bug" | "defect"
            | "x" | "task" | "t" | "lab" | "l" | "evolution"
step       := "step" bid ("done" | "todo" | "pass" | "fail") { string | date }
merge      := "merge" bid "into" vid date?
discard    := "discard" bid { string | date }
target     := "target" bid vid
now        := "now"
eos        := "eos" { id | date }
collapse   := "collapse" vid integer
note       := "note" (vid | bid) string

date       := "@" YYYY "-" MM ("-" DD)?
idlist     := id ("," id)*
id, vid, bid := [A-Za-z0-9._-]+
string     := '"' ( '\' any | not '"' )* '"'
comment    := ("#" | "%%") to end of line
\end{lstlisting}

\paragraph{Semantics of \texttt{release \dots{} merges}.}
The statement \lstinline[style=hpb]|release v1.1 merges I-21, X-102|
publishes the baseline \texttt{v1.1} on the current line, in the sense of
\cref{def:gate}, with
\texttt{I-21} and \texttt{X-102} as its targeted set: both are merged into
it at that moment, and their terminals share the column just before it.
To record that a branch was resolved before the release, write
\texttt{merge I-21 into v1.1} while \texttt{v1.1} is still planned; the
branch is then drawn as resolved and awaiting release, and it counts as
merged when \texttt{v1.1} is released.

\paragraph{A second example: hardware revision and line fork.}
The text below describes the hardware situation of
\cref{fig:hardware-fork}, with its own illustrative items: a laboratory
evolution that fails and then passes its EMC test, a hardware baseline
\texttt{RevB} that creates a new reference, the child line that starts from
it, and the old line that keeps receiving firmware fixes until its end of
support.

\begin{lstlisting}[style=hpb,basicstyle=\ttfamily\footnotesize]
hpb "DMS-4000: hardware revision"
line RA "DMS-4000 Rev A"
release v1.2 @2026-01
lab L-5 "Sensor board v2"
step L-5 fail "EMC test"
improvement I-30 "Low-power mode"
step L-5 pass "EMC retest"
release v1.3
release RevB hardware merges L-5, I-30
line RB "DMS-4000 Rev B" from RevB
release v2.0 on RB
use RA
release v1.4
issue X-118 "Relay chatter"
release v1.4.1 merges X-118
eos RA @2026-12
\end{lstlisting}

% flush this appendix's floats before the next appendix
\clearpage

% ---------------------------------------------------------------------
%  appendix-b-package.tex -- using the hpb.sty package (API v1.0)
% ---------------------------------------------------------------------
\section{Using the \texttt{hpb.sty} Package}
\label{app:package}

The notation of \cref{sec:notation} is distributed as a single \LaTeX{}
package, \texttt{hpb.sty} (version~1.1, MIT license). Every HPB diagram in
this paper is drawn with it. The point is standardization: glyph sizes,
stroke widths, colors, fonts and the placement of tags, labels and terminals
are fixed by the package, so diagrams drawn by different people look the
same and follow the layout rules of \cref{sec:notation}. Use it for papers,
reports and specifications; for web pages and documentation that change
often, HPB Markup (\cref{sec:tooling,app:markup}) draws the same notation
from text. This appendix summarizes the package interface.

\paragraph{Installation and requirements.}
Copy \texttt{hpb.sty} next to the main \texttt{.tex} file (on Overleaf,
upload it to the project root) and load it with \texttt{\textbackslash
usepackage\{hpb\}}. The package requires only TikZ, \texttt{xcolor} and
\texttt{xparse}, which are part of every current \TeX{} distribution; it
loads the TikZ libraries \texttt{positioning}, \texttt{calc},
\texttt{shapes.geometric}, \texttt{arrows.meta},
\texttt{decorations.pathmorphing}, \texttt{backgrounds} and \texttt{fit}.
It uses no engine-specific features, needs no shell escape and no external
files, and is intended for pdf\LaTeX{} as used by Overleaf and arXiv.
Compile once with the engine of your submission before relying on it.

\paragraph{Canvas and coordinates.}
Every diagram is an \texttt{hpbgraph} environment, a TikZ picture in which
the $x$ unit is one ordinal column and the $y$ unit is one lane. All
macros therefore take plain numbers: columns (integers, or decimals for
fine adjustment) and lanes, with $0$ for the main line, positive lanes
above it and negative lanes below it (\cref{sec:notation}). A child line
created by a line fork is simply another $y$ value (e.g., $-3.5$). The
environment accepts three keys: \texttt{col} (column width, default
1.1\,cm), \texttt{lane} (lane height, default 1.0\,cm; tags are anchored at
a fixed distance above their node, and 1.0\,cm is the smallest value that
keeps a tag clear of a branch on lane $+1$) and \texttt{legend} (width of
the legend and origin boxes, default 2.3\,cm). This paper sets one canvas
for all diagrams with \texttt{\textbackslash hpbset\{col=0.75cm, lane=1cm,
legend=1.7cm\}}, because one document uses one canvas. Any other key is passed to the
underlying \texttt{tikzpicture}; \texttt{\textbackslash hpbset} sets
defaults for a whole document. Only spacing scales with these keys; glyph
sizes are fixed in points by the notation. Plain TikZ commands can be mixed
freely with the package macros in the same coordinates. Internally the
package draws on six layers (background for gate and bundle bands, guides,
branch paths and connectors, main lines and streams, nodes, and a front
layer for tags and labels), so that the order of commands in the source
does not affect the result.

\begin{figure}[tbp]
  \centering
  \begin{minipage}[c]{0.41\linewidth}
\begin{lstlisting}[style=hpblatex,basicstyle=\ttfamily\scriptsize,
  moretexcs={hpbgate,hpbtarget}]
\begin{hpbgraph}[col=0.7cm, legend=1.5cm]
  % lanes and main line
  \hpblane{1}{imp}{Improv. (I)}{0}{8.6}
  \hpblane{-1}{iss}{Issue (X)}{0}{8.6}
  \hpbmainline{0}{6}
  \hpbmainline*{6}{8.6}
  \hpborigin{0}{Product A}
  % baselines
  \hpbrelease{1}{v2.0}
  \hpbrelease{5}{v2.0.1}
  \hpbplanned{8}{v2.1}{Jun}
  % branches
  \hpbbranch[label={I-40 Export}]
    {imp}{1}{1}{2}{6}{target}{8}
  \hpbbranch[label={X-77 Login}]
    {iss}{1}{-1}{3}{4}{merge}{5}
  % markers
  \hpbnow{6}{1}{-1}
  \hpbgate[1.5]{8}{1}{-1}{gate v2.1}
\end{hpbgraph}
\end{lstlisting}
  \end{minipage}\hfill
  \begin{minipage}[c]{0.57\linewidth}
    \centering
    % natural size (about 8.3cm wide): no scaling, glyphs at normative size
    \begin{hpbgraph}[col=0.7cm, legend=1.5cm]
      % lanes and main line
      \hpblane{1}{imp}{Improv. (I)}{0}{8.6}
      \hpblane{-1}{iss}{Issue (X)}{0}{8.6}
      \hpbmainline{0}{6}
      \hpbmainline*{6}{8.6}
      \hpborigin{0}{Product A}
      % baselines
      \hpbrelease{1}{v2.0}
      \hpbrelease{5}{v2.0.1}
      \hpbplanned{8}{v2.1}{Jun}
      % branches
      \hpbbranch[label={I-40 Export}]
        {imp}{1}{1}{2}{6}{target}{8}
      \hpbbranch[label={X-77 Login}]
        {iss}{1}{-1}{3}{4}{merge}{5}
      % markers
      \hpbnow{6}{1}{-1}
      \hpbgate[1.5]{8}{1}{-1}{gate v2.1}
    \end{hpbgraph}
  \end{minipage}
  \caption{A complete \texttt{hpb.sty} example (illustrative data): the
    source (left) and its rendering (right). Columns are ordinal
    (one event per column) and lanes follow \cref{sec:notation}. The issue X-77,
    observed in \texttt{v2.0}, is merged into the patch baseline
    \texttt{v2.0.1}; the improvement I-40 is open at NOW and targeted at the
    planned baseline \texttt{v2.1}, whose gate it keeps closed until it is
    verified, re-targeted or discarded.}
  \label{fig:package-example}
\end{figure}

\paragraph{Interface.}
\Cref{tab:api} lists the public macros and \cref{tab:api-keys} the options
of \texttt{\textbackslash hpbbranch}, the central macro: one call draws a
complete branch, i.e., the fork from its baseline, the start node with the
type letter, the run along the lane, the terminal and, for merged and
targeted branches, the connection to the closing or target baseline. Types
are written \texttt{imp}, \texttt{iss}, \texttt{task}, \texttt{lab} (or
\texttt{I}, \texttt{X}, \texttt{T}, \texttt{L}); an unknown type, terminal
or step state stops compilation with an explicit error message rather than
producing a silently wrong diagram. The color names (\texttt{hpbInk},
\texttt{hpbMain}, \texttt{hpbImp}, \texttt{hpbIss}, \texttt{hpbTask},
\texttt{hpbLab}, \texttt{hpbHw}, \texttt{hpbDisc}, \texttt{hpbNow},
\texttt{hpbGuide}, \texttt{hpbLate}) are part of the interface, so that
non-HPB figures, such as charts and state diagrams, can use the same palette.

\providecommand{\HPBmac}[1]{\texttt{\textbackslash #1}}
\providecommand{\HPBarg}[1]{\ensuremath{\langle}\textit{#1}\ensuremath{\rangle}}

\begin{table}[tbp]
  \centering
  \small
  \caption{Options of \texttt{\textbackslash hpbbranch} (all optional).}
  \label{tab:api-keys}
  \renewcommand{\arraystretch}{1.0}
  \begin{tabularx}{\linewidth}{@{}>{\raggedright\arraybackslash}p{3.0cm}
      >{\raggedright\arraybackslash}p{2.0cm}
      >{\raggedright\arraybackslash}X@{}}
    \toprule
    Key & Default & Effect \\
    \midrule
    \texttt{line=}\HPBarg{y} & \texttt{0}
      & $y$ of the line the branch forks from (child line or stream). \\
    \texttt{mergeline=}\HPBarg{y} & = \texttt{line}
      & $y$ of the line the branch merges into. \\
    \texttt{forkoff=}\HPBarg{cols} & \texttt{0.5}
      & Offset of the vertical drop from the fork baseline (\cref{sec:notation}). \\
    \texttt{mergeoff=}\HPBarg{cols} & \texttt{0.5}
      & Offset of the vertical rise before the merge baseline. \\
    \texttt{label=}\HPBarg{text} & none
      & Branch label (reference and title) at the start node. \\
    \texttt{labelpos=}\HPBarg{pos} & \texttt{above}
      & Label position, as for \HPBmac{hpblabel}. \\
    \texttt{letter=}\HPBarg{text} & type letter
      & Overrides the letter of the start node. \\
    \texttt{nostart} & --
      & Omits the start node. \\
    \texttt{faded} & --
      & Draws at 35\,\% opacity (former state, e.g., an old target). \\
    \bottomrule
  \end{tabularx}
\end{table}

\begin{table}[p]
  \centering
  \footnotesize
  \caption{Public macros of \texttt{hpb.sty} v1.1. Columns and lanes are
    plain numbers in diagram units; optional arguments in square brackets.
    \HPBarg{type}: \texttt{imp}, \texttt{iss}, \texttt{task}, \texttt{lab}
    (legend boxes also accept \texttt{main}, \texttt{hw}, \texttt{disc}).}
  \label{tab:api}
  \renewcommand{\arraystretch}{1.0}
  \begin{tabularx}{\linewidth}{@{}>{\raggedright\arraybackslash}p{3.0cm}
      >{\raggedright\arraybackslash}p{5.6cm}
      >{\raggedright\arraybackslash}X@{}}
    \toprule
    Macro & Arguments & Draws \\
    \midrule
    \multicolumn{3}{@{}l}{\textit{Canvas}} \\
    \texttt{hpbgraph} (env.) & \texttt{[col=, lane=, legend=, }\HPBarg{tikz}\texttt{]}
      & The canvas: $x$ = column, $y$ = lane. \\
    \HPBmac{hpbset} & \texttt{\{}\HPBarg{keys}\texttt{\}}
      & Document-wide defaults for the canvas keys. \\
    \HPBmac{hpblane} & \HPBarg{lane}\HPBarg{type}\HPBarg{label}\HPBarg{from}\HPBarg{to}
      & Dotted lane guide and legend box at its left end. \\
    \midrule
    \multicolumn{3}{@{}l}{\textit{Lines and baselines}} \\
    \HPBmac{hpbmainline}[\texttt{*}] & \HPBarg{from}\HPBarg{to}[\HPBarg{y}]
      & Main line, 3\,pt; starred: dashed (plan). \\
    \HPBmac{hpborigin} & \HPBarg{col}\HPBarg{label}[\HPBarg{y}]
      & Creation dot (not a baseline) and product box. \\
    \HPBmac{hpbrelease} & [\HPBarg{caption}]\HPBarg{col}\HPBarg{name}[\HPBarg{y}]
      & Released baseline with tag and arrow; optional caption (e.g., applies-to set). \\
    \HPBmac{hpbplanned}[\texttt{*}] & [\HPBarg{caption}]\HPBarg{col}\HPBarg{name}\HPBarg{date}[\HPBarg{y}]
      & Planned baseline, dashed tag, target date; starred: overdue (``!'' badge, red). \\
    \HPBmac{hpbhardware} & [\HPBarg{caption}]\HPBarg{col}\HPBarg{name}[\HPBarg{y}]
      & Hardware diamond, indigo tag, caption (default \emph{new ref.}). \\
    \HPBmac{hpbprocess} & [\HPBarg{caption}]\HPBarg{col}\HPBarg{name}[\HPBarg{y}]
      & Process baseline (square). \\
    \HPBmac{hpbstream}[\texttt{*}] & [\HPBarg{offset}]\HPBarg{fork col}\HPBarg{y}\HPBarg{to col}[\HPBarg{from y}]
      & Release stream leaving a baseline; starred: planned. \\
    \HPBmac{hpbmilestone} & \HPBarg{col}\HPBarg{label}[\HPBarg{y}]
      & Life-cycle milestone tick. \\
    \HPBmac{hpbseparator} & \HPBarg{y}\HPBarg{from}\HPBarg{to}
      & Thin rule between lane bands. \\
    \HPBmac{hpbcollapsed} & \HPBarg{col}\HPBarg{count}[\HPBarg{y}]
      & Count badge of collapsed history. \\
    \HPBmac{hpbeos} & \HPBarg{col}[\HPBarg{y}]
      & End-of-support bar labeled \textsf{EoS}. \\
    \HPBmac{hpblinefork} & \HPBarg{col}\HPBarg{from y}\HPBarg{to y}
      & Lineage connector to a child line. \\
    \midrule
    \multicolumn{3}{@{}l}{\textit{Branches and events}} \\
    \HPBmac{hpbbranch} & [\HPBarg{keys}]\HPBarg{type}\HPBarg{fork col}\HPBarg{lane}%
      \HPBarg{start col}\HPBarg{end col}\HPBarg{end}\HPBarg{merge col}
      & Complete branch; \HPBarg{end} $\in$ \{\texttt{merge}, \texttt{target},
        \texttt{open}, \texttt{discard}, \texttt{none}\}; \HPBarg{merge col}
        is empty for \texttt{open}, \texttt{discard}, \texttt{none}. \\
    \HPBmac{hpbtarget} & [\HPBarg{keys}]\HPBarg{from col}\HPBarg{lane}\HPBarg{merge col}
      & Additional dashed target continuation (re-targeting). \\
    \HPBmac{hpbstart} & \HPBarg{type}\HPBarg{col}\HPBarg{lane}\HPBarg{letter}
      & Start node with letter (stand-alone). \\
    \HPBmac{hpbnode} & \HPBarg{type}\HPBarg{col}\HPBarg{lane}
      & Plain event node. \\
    \HPBmac{hpbterminal} & \HPBarg{kind}\HPBarg{type}\HPBarg{col}\HPBarg{lane}
      & Terminal glyph (stand-alone). \\
    \HPBmac{hpbstep} & \HPBarg{type}\HPBarg{col}\HPBarg{lane}\HPBarg{state}
      & Step: \texttt{done}, \texttt{todo}, \texttt{pass}, \texttt{fail}. \\
    \HPBmac{hpbpoststep} & \HPBarg{type}\HPBarg{from col}\HPBarg{col}\HPBarg{lane}\HPBarg{state}
      & Step after closure, on a dotted continuation. \\
    \HPBmac{hpbdepends} & [\HPBarg{tikz}]\HPBarg{x,y}\HPBarg{x,y}\HPBarg{label}
      & Dependency connector (dotted, open head). \\
    \HPBmac{hpbpropagates} & [\HPBarg{tikz}]\HPBarg{x,y}\HPBarg{x,y}\HPBarg{label}
      & Propagation connector (dashed, filled head). \\
    \HPBmac{hpblabel} & [\HPBarg{pos}]\HPBarg{col}\HPBarg{lane}\HPBarg{text}
      & Label with white halo (default position \texttt{above}). \\
    \midrule
    \multicolumn{3}{@{}l}{\textit{Markers and annotations}} \\
    \HPBmac{hpbnow} & \HPBarg{col}\HPBarg{lane top}\HPBarg{lane bottom}
      & NOW marker across the given lanes. \\
    \HPBmac{hpbgate} & [\HPBarg{right}]\HPBarg{col}\HPBarg{lane top}\HPBarg{lane bottom}\HPBarg{label}
      & Gate band around a planned baseline. \\
    \HPBmac{hpbbundle} & \HPBarg{x,y}\HPBarg{x,y}\HPBarg{label}
      & Launch-bundle frame across lines. \\
    \HPBmac{hpbfit} & \HPBarg{picture}
      & Scales a picture down to the line width only if wider. \\
    \HPBmac{hpbcallout} & [\HPBarg{tikz}]\HPBarg{x,y}\HPBarg{text}\HPBarg{x,y}
      & Annotation box with leader line. \\
    \HPBmac{hpbaxis} & \HPBarg{from}\HPBarg{to}\HPBarg{y}
      & Ordinal time axis. \\
    \HPBmac{hpbtick} & \HPBarg{col}\HPBarg{text}\HPBarg{y}
      & Calendar tick under the axis. \\
    \bottomrule
  \end{tabularx}

\end{table}

\paragraph{Conventions that the package does not enforce.}
Two layout rules depend on the data and are left to the author (or to a
layout algorithm such as \cref{alg:layout}; the HPB Markup renderer
computes columns and lanes for you). First, each event must be given its own column;
the geometry
of a branch requires
$\textit{fork col} + \textit{forkoff} < \textit{start col} \le
\textit{end col} \le \textit{merge col} - 1$.
Second, when several branches leave or enter the same baseline on the same
side, the branch on the farther lane must be given the smaller
\texttt{forkoff} or \texttt{mergeoff}, so that the corners nest
(\cref{sec:notation}). For a figure in a one-column article, the natural width of
the picture should not much exceed the text width, so that, after
\texttt{\textbackslash hpbfit\{\dots\}} (which scales down only when
needed), the smallest text (6.5\,pt) stays above 6\,pt; a
practical estimate of the natural width is
$\textit{legend} + 0.35\,\mathrm{cm} + (\textit{last col} + 0.8) \cdot
\textit{col}$.

\paragraph{Example.}
\Cref{fig:package-example} shows a complete diagram and its source: a
product line with a released baseline \texttt{v2.0}, a defect fixed by the
patch baseline \texttt{v2.0.1}, an improvement targeted at the planned
baseline \texttt{v2.1}, the NOW marker and the gate of \texttt{v2.1}. The
diagram needs twelve macro calls, none of which contains a coordinate in
centimeters.



\end{document}
